\documentclass[11pt,a4paper]{article}
\usepackage[utf8]{inputenc}
\usepackage[vietnamese,english]{babel}
\usepackage{amsmath,amssymb}
\usepackage{geometry}
\geometry{a4paper, margin=2cm}
\usepackage{enumitem}
\usepackage{xcolor}
\usepackage{hyperref}
\usepackage{listings}

\title{\textbf{PFP191 - Programming Fundamentals with Python}\\\Large Complete Question Bank}
\author{Google Antigravity Image-to-Text Pipeline}
\date{\today}

\begin{document}
\maketitle
\tableofcontents
\newpage

\begin{minipage}{\textwidth}
\subsection*{Question 1 \quad \footnotesize{[0001_q001_Question_1.png]}}
\noindent \textbf{Which of the following is a comment in Python?}

\begin{enumerate}[label=\Alph*.]
    \item // This is a test
    \item /* This is a test */
    \item \textbf{# This is a test} \hfill \textcolor{blue}{\textbf{[Correct]}}
    \item $ This is a test
\end{enumerate}

\noindent \textbf{Correct Answer:} \textcolor{blue}{\textbf{C}} \\[4pt]
\noindent \textbf{Explanation:}\\
Trong Python, ký tự thăng `#` (hash / pound) được dùng để bắt đầu một dòng chú thích (comment). Trình thông dịch Python sẽ bỏ qua toàn bộ phần văn bản nằm sau ký tự `#` cho đến hết dòng đó.\newline
- Cú pháp `//` là chú thích dòng đơn trong C/C++/Java/JavaScript (trong Python `//` là toán tử chia lấy phần nguyên).\newline
- Cú pháp `/* ... */` là chú thích nhiều dòng trong C/C++/Java.\newline
- Ký tự `$` không được dùng để tạo chú thích trong Python.\newline
Do đó, `# This is a test` là cú pháp chú thích chính xác.
\end{minipage}
\vspace{0.8cm}
\hrule
\vspace{0.5cm}

\begin{minipage}{\textwidth}
\subsection*{Question 2 \quad \footnotesize{[0002_q002_Question_2.png]}}
\noindent \textbf{What is the result of this statement?\newline

\newline
\newline

\newline
```python\newline

\newline
print(ord('foo'))\newline

\newline
```}

\begin{enumerate}[label=\Alph*.]
    \item 102
    \item \textbf{It raises an exception} \hfill \textcolor{blue}{\textbf{[Correct]}}
    \item 102 111 111
    \item 324
\end{enumerate}

\noindent \textbf{Correct Answer:} \textcolor{blue}{\textbf{B}} \\[4pt]
\noindent \textbf{Explanation:}\\
Hàm tích hợp sẵn `ord(c)` trong Python nhận đầu vào là một chuỗi ký tự có độ dài đúng bằng 1 (`string of length 1`) biểu diễn một ký tự Unicode, và trả về mã số nguyên Unicode tương ứng của ký tự đó.\newline
Trong đoạn mã trên, `'foo'` là một chuỗi có độ dài bằng 3 ký tự. Khi gọi `ord('foo')`, Python sẽ ném ra ngoại lệ `TypeError: ord() expected a character, but string of length 3 found`.\newline
Do đó, câu lệnh này gây ra lỗi (It raises an exception).
\end{minipage}
\vspace{0.8cm}
\hrule
\vspace{0.5cm}

\begin{minipage}{\textwidth}
\subsection*{Question 3 \quad \footnotesize{[0003_q003_Question_3.png]}}
\noindent \textbf{Once again consider this dictionary:\newline

\newline
\newline

\newline
```python\newline

\newline
d = 'foo': 100, 'bar': 200, 'baz': 300\newline

\newline
```\newline

\newline
\newline

\newline
What is the result of this statement:\newline

\newline
\newline

\newline
```python\newline

\newline
d['bar':'baz']\newline

\newline
```}

\begin{enumerate}[label=\Alph*.]
    \item \textbf{It raises an exception} \hfill \textcolor{blue}{\textbf{[Correct]}}
    \item 200 300
    \item [200, 300]
    \item (200, 300)
\end{enumerate}

\noindent \textbf{Correct Answer:} \textcolor{blue}{\textbf{A}} \\[4pt]
\noindent \textbf{Explanation:}\\
Trong Python, kiểu dữ liệu từ điển (`dict`) là một cấu trúc dữ liệu ánh xạ (mapping), truy xuất giá trị bằng các khóa duy nhất (keys) theo cú pháp `d[key]`. Dictionary không phải là kiểu tuần tự có thứ tự đánh chỉ mục (sequence) nên không hỗ trợ cú pháp cắt lát (slicing) `[start:stop]`.\newline
Khi thực hiện `d['bar':'baz']`, Python cố gắng truyền một đối tượng `slice('bar', 'baz', None)` làm khóa cho từ điển. Vì kiểu `slice` là unhashable (không thể băm), Python sẽ ném ra lỗi `TypeError: unhashable type: 'slice'`.\newline
Do đó, câu lệnh này gây ra lỗi ngoại lệ (It raises an exception).
\end{minipage}
\vspace{0.8cm}
\hrule
\vspace{0.5cm}

\begin{minipage}{\textwidth}
\subsection*{Question 4 \quad \footnotesize{[0004_q004_Question_4.png]}}
\noindent \textbf{What will the following Python code print out?\newline

\newline
\newline

\newline
```python\newline

\newline
friends = [ 'Joseph', 'Glenn', 'Sally' ]\newline

\newline
friends.sort()\newline

\newline
print friends[1]\newline

\newline
```}

\begin{enumerate}[label=\Alph*.]
    \item Glenn
    \item \textbf{Joseph} \hfill \textcolor{blue}{\textbf{[Correct]}}
    \item Sally
    \item Friends
\end{enumerate}

\noindent \textbf{Correct Answer:} \textcolor{blue}{\textbf{B}} \\[4pt]
\noindent \textbf{Explanation:}\\
Các bước thực thi đoạn mã:\newline
1. Khởi tạo danh sách chuỗi: `friends = ['Joseph', 'Glenn', 'Sally']`.\newline
2. Phương thức `friends.sort()` sắp xếp danh sách tại chỗ (in-place) theo thứ tự bảng chữ cái (thứ tự từ điển):\newline
   - 'Glenn' (bắt đầu bằng 'G') -> chỉ số 0\newline
   - 'Joseph' (bắt đầu bằng 'J') -> chỉ số 1\newline
   - 'Sally' (bắt đầu bằng 'S') -> chỉ số 2\newline
   Danh sách sau khi sắp xếp: `['Glenn', 'Joseph', 'Sally']`.\newline
3. Lệnh `friends[1]` truy xuất phần tử tại chỉ số 1, chính là `'Joseph'`.\newline
Vì vậy, kết quả in ra là `Joseph`.
\end{minipage}
\vspace{0.8cm}
\hrule
\vspace{0.5cm}

\begin{minipage}{\textwidth}
\subsection*{Question 5 \quad \footnotesize{[0005_q005_Question_5.png]}}
\noindent \textbf{How do you insert something on a new line in a file?}

\begin{enumerate}[label=\Alph*.]
    \item write.newline('x')
    \item \textbf{n} \hfill \textcolor{blue}{\textbf{[Correct]}}
    \item You cannot do this
    \item type the content on the line below
\end{enumerate}

\noindent \textbf{Correct Answer:} \textcolor{blue}{\textbf{B}} \\[4pt]
\noindent \textbf{Explanation:}\\
Trong Python và xử lý văn bản nói chung, ký tự thoát `\n` (newline character) đại diện cho việc kết thúc dòng hiện tại và bắt đầu một dòng mới.\newline
Khi ghi dữ liệu vào tệp tin bằng phương thức `file.write()` trong Python, trình thông dịch không tự động thêm ký tự xuống dòng. Do đó, lập trình viên cần chủ động chèn ký tự `\n` vào chuỗi để dữ liệu tiếp theo được ghi trên một dòng mới.
\end{minipage}
\vspace{0.8cm}
\hrule
\vspace{0.5cm}

\begin{minipage}{\textwidth}
\subsection*{Question 6 \quad \footnotesize{[0006_q006_Question_6.png]}}
\noindent \textbf{What is the incorrect ways to create a function in Python? (Choose all the apply.)}

\begin{enumerate}[label=\Alph*.]
    \item \textbf{create myFunction():} \hfill \textcolor{blue}{\textbf{[Correct]}}
    \item def myFunction():
    \item \textbf{func myfunction():} \hfill \textcolor{blue}{\textbf{[Correct]}}
    \item \textbf{def function myfunction():} \hfill \textcolor{blue}{\textbf{[Correct]}}
\end{enumerate}

\noindent \textbf{Correct Answer:} \textcolor{blue}{\textbf{A, C, D}} \\[4pt]
\noindent \textbf{Explanation:}\\
Trong Python, cú pháp chuẩn xác và duy nhất để định nghĩa một hàm là sử dụng từ khóa `def`, theo sau là tên hàm, danh sách tham số trong ngoặc đơn và dấu hai chấm `:`, ví dụ: `def myFunction():`.\newline
- Đáp án B (`def myFunction():`) là cú pháp hoàn toàn ĐÚNG.\newline
- Các đáp án A (`create`), C (`func`), và D (`def function`) đều sử dụng từ khóa sai không tồn tại trong ngữ pháp Python để định nghĩa hàm và sẽ gây lỗi `SyntaxError`.\newline
Vì câu hỏi yêu cầu chọn tất cả các cách KHÔNG ĐÚNG (incorrect ways), nên ta chọn A, C và D.
\end{minipage}
\vspace{0.8cm}
\hrule
\vspace{0.5cm}

\begin{minipage}{\textwidth}
\subsection*{Question 7 \quad \footnotesize{[0007_q007_Question_7.png]}}
\noindent \textbf{Which collection does not allow duplicate members?}

\begin{enumerate}[label=\Alph*.]
    \item \textbf{SET} \hfill \textcolor{blue}{\textbf{[Correct]}}
    \item DICTIONARY
    \item TUPLE
    \item LIST
\end{enumerate}

\noindent \textbf{Correct Answer:} \textcolor{blue}{\textbf{A}} \\[4pt]
\noindent \textbf{Explanation:}\\
Trong Python, `Set` (tập hợp) là tập hợp các phần tử không có thứ tự, không có chỉ số và các phần tử là duy nhất, tức không cho phép chứa các phần tử trùng lặp (does not allow duplicate members). Nếu thêm một phần tử đã tồn tại vào set, phần tử đó sẽ tự động bị bỏ qua.\newline
- `List` và `Tuple` đều cho phép chứa các phần tử trùng lặp.\newline
- `Dictionary` chứa các cặp khóa - giá trị (`key: value`), trong đó khóa là duy nhất nhưng giá trị có thể trùng nhau; tuy nhiên khi nói về tập hợp các phần tử (collection members), `Set` là kiểu dữ liệu chuẩn xác định không cho phép trùng lặp.
\end{minipage}
\vspace{0.8cm}
\hrule
\vspace{0.5cm}

\begin{minipage}{\textwidth}
\subsection*{Question 8 \quad \footnotesize{[0008_q008_Question_8.png]}}
\noindent \textbf{What is the name of this symbol (>>>)?}

\begin{enumerate}[label=\Alph*.]
    \item \textbf{Chevrons} \hfill \textcolor{blue}{\textbf{[Correct]}}
    \item Hash tags
    \item Colon
    \item Semi Colon
\end{enumerate}

\noindent \textbf{Correct Answer:} \textcolor{blue}{\textbf{A}} \\[4pt]
\noindent \textbf{Explanation:}\\
Ký hiệu `>>>` là dấu nhắc lệnh chính (primary prompt) của trình thông dịch tương tác Python (Python interactive shell / REPL). Trong tài liệu và thuật ngữ lập trình Python, ký hiệu ba dấu lớn hơn liên tiếp này thường được gọi là "chevrons" (hoặc triple chevrons).\newline
- B: Hash tag là ký tự `#`.\newline
- C: Colon là dấu hai chấm `:`.\newline
- D: Semi Colon là dấu chấm phẩy `;`.
\end{minipage}
\vspace{0.8cm}
\hrule
\vspace{0.5cm}

\begin{minipage}{\textwidth}
\subsection*{Question 9 \quad \footnotesize{[0009_q009_Question_9.png]}}
\noindent \textbf{What is the value of the following expression? 45 % 10}

\begin{enumerate}[label=\Alph*.]
    \item 4
    \item 2
    \item \textbf{5} \hfill \textcolor{blue}{\textbf{[Correct]}}
    \item 0.42
\end{enumerate}

\noindent \textbf{Correct Answer:} \textcolor{blue}{\textbf{C}} \\[4pt]
\noindent \textbf{Explanation:}\\
Trong Python, toán tử `%` là toán tử chia lấy phần dư (modulo operator).\newline
Biểu thức `45 % 10` tính toán số dư khi chia 45 cho 10:\newline
Ta có: 45 = 10 * 4 + 5 (thương là 4 và số dư là 5).\newline
Do đó, giá trị của biểu thức là 5.
\end{minipage}
\vspace{0.8cm}
\hrule
\vspace{0.5cm}

\begin{minipage}{\textwidth}
\subsection*{Question 10 \quad \footnotesize{[0010_q010_Question_10.png]}}
\noindent \textbf{What is the purpose of the following Python code?\newline

\newline
\newline

\newline
```python\newline

\newline
fhand = open('mbox.txt')\newline

\newline
x = 0\newline

\newline
for line in fhand:\newline

\newline
    x = x + 1\newline

\newline
print x\newline

\newline
```}

\begin{enumerate}[label=\Alph*.]
    \item Remove the leading and trailing spaces from each line in mbox.txt
    \item \textbf{Count the lines in the file 'mbox.txt'} \hfill \textcolor{blue}{\textbf{[Correct]}}
    \item Convert the lines in mbox.txt to upper case
    \item Reverse the order of the lines in mbox.txt
\end{enumerate}

\noindent \textbf{Correct Answer:} \textcolor{blue}{\textbf{B}} \\[4pt]
\noindent \textbf{Explanation:}\\
Phân tích từng câu lệnh của đoạn mã:\newline
1. `fhand = open('mbox.txt')`: Mở tệp văn bản 'mbox.txt' và trả về một đối tượng tệp (file handle) có thể lặp được theo từng dòng.\newline
2. `x = 0`: Khởi tạo biến đếm `x` ban đầu bằng 0.\newline
3. `for line in fhand:`: Vòng lặp duyệt qua từng dòng trong tệp tin.\newline
4. `x = x + 1`: Với mỗi dòng được đọc, biến đếm `x` được tăng thêm 1 đơn vị.\newline
5. `print x`: In tổng số dòng đã đếm được sau khi duyệt hết toàn bộ tệp.\newline
Vì vậy, mục đích của đoạn mã là đếm số dòng có trong tệp 'mbox.txt' (Count the lines in the file 'mbox.txt').
\end{minipage}
\vspace{0.8cm}
\hrule
\vspace{0.5cm}

\begin{minipage}{\textwidth}
\subsection*{Question 11 \quad \footnotesize{[0011_q011_Question_11.png]}}
\noindent \textbf{What is the output of the following code?\newline

\newline
\newline

\newline
```python\newline

\newline
class Pokemon():\newline

\newline
\newline

\newline
    def __init__(self, name, type):\newline

\newline
        self.name = name\newline

\newline
        self.type = type\newline

\newline
\newline

\newline
    def stringPokemon(self):\newline

\newline
        print(f"Pokemon name is self.name and type is self.type")\newline

\newline
\newline

\newline
class GrassType(Pokemon):\newline

\newline
\newline

\newline
    # overrides the stringPokemon() function on 'Pokemon' class\newline

\newline
    def stringPokemon(self):\newline

\newline
        print(f"Grass type pokemon name is self.name")\newline

\newline
\newline

\newline
poke1 = GrassType('Bulbasaur', 'Grass')\newline

\newline
poke1.stringPokemon\newline

\newline
poke1.stringPokemon()\newline

\newline
poke2 = Pokemon('Charizard', 'Fire')\newline

\newline
poke2.stringPokemon\newline

\newline
poke2.stringPokemon()\newline

\newline
```}

\begin{enumerate}[label=\Alph*.]
    \item \textbf{Grass type pokemon name is Bulbasaur Pokemon name is Charizard and type is Fire} \hfill \textcolor{blue}{\textbf{[Correct]}}
    \item Pokemon name is Bulbasaur and type is Grass Pokemon name is Charizard and type is Fire
    \item Grass type pokemon name is Bulbasaur Grass type pokemon name is Charizard
    \item Error because the extending class has a stringPokemon() function which already exists.
\end{enumerate}

\noindent \textbf{Correct Answer:} \textcolor{blue}{\textbf{A}} \\[4pt]
\noindent \textbf{Explanation:}\\
Chi tiết từng bước thực thi:\newline
1. Lớp `GrassType` kế thừa từ lớp `Pokemon` và ghi đè (override) phương thức `stringPokemon()`.\newline
2. Tạo đối tượng `poke1 = GrassType('Bulbasaur', 'Grass')`.\newline
   - Dòng `poke1.stringPokemon`: chỉ tham chiếu đến phương thức mà không có dấu gọi hàm `()`, nên không thực thi và không in ra gì.\newline
   - Dòng `poke1.stringPokemon()`: thực thi phương thức đã ghi đè của lớp `GrassType`, in ra: `Grass type pokemon name is Bulbasaur`.\newline
3. Tạo đối tượng `poke2 = Pokemon('Charizard', 'Fire')`.\newline
   - Dòng `poke2.stringPokemon`: chỉ tham chiếu không gọi hàm.\newline
   - Dòng `poke2.stringPokemon()`: gọi phương thức của lớp cha `Pokemon`, in ra: `Pokemon name is Charizard and type is Fire`.\newline
Kết hợp lại, đầu ra là: `Grass type pokemon name is Bulbasaur Pokemon name is Charizard and type is Fire`.
\end{minipage}
\vspace{0.8cm}
\hrule
\vspace{0.5cm}

\begin{minipage}{\textwidth}
\subsection*{Question 12 \quad \footnotesize{[0012_q012_Question_12.png]}}
\noindent \textbf{Assume x and y are assigned as follows:\newline

\newline
\newline

\newline
```python\newline

\newline
x = 5\newline

\newline
y = -5\newline

\newline
```\newline

\newline
\newline

\newline
What is the effect of this statement:\newline

\newline
\newline

\newline
```python\newline

\newline
x, y = (y, x)[::-1]\newline

\newline
```}

\begin{enumerate}[label=\Alph*.]
    \item \textbf{The values of x and y are unchanged} \hfill \textcolor{blue}{\textbf{[Correct]}}
    \item Both x and y are -5
    \item Both x and y are 5
    \item The values of x and y are swapped
\end{enumerate}

\noindent \textbf{Correct Answer:} \textcolor{blue}{\textbf{A}} \\[4pt]
\noindent \textbf{Explanation:}\\
Phân tích biểu thức gán `x, y = (y, x)[::-1]`:\newline
1. Ban đầu: `x = 5`, `y = -5`.\newline
2. Vế phải tạo ra tuple `(y, x)` tương ứng là `(-5, 5)`.\newline
3. Phép cắt lát bước nhảy âm `[::-1]` đảo ngược tuple `(-5, 5)` trở thành `(5, -5)`.\newline
4. Thực hiện phép gán giải nén (unpacking): `x, y = (5, -5)`:\newline
   - `x` nhận giá trị 5 (không đổi so với ban đầu).\newline
   - `y` nhận giá trị -5 (không đổi so với ban đầu).\newline
Do đó, giá trị của x và y hoàn toàn không bị thay đổi (The values of x and y are unchanged).
\end{minipage}
\vspace{0.8cm}
\hrule
\vspace{0.5cm}

\begin{minipage}{\textwidth}
\subsection*{Question 13 \quad \footnotesize{[0013_q013_Question_13.png]}}
\noindent \textbf{What's the output of the following code snippet?\newline

\newline
\newline

\newline
```python\newline

\newline
>>> class Dog:\newline

\newline
...     def walk(self):\newline

\newline
...         return "*walking*"\newline

\newline
...\newline

\newline
...     def speak(self):\newline

\newline
...         return "Woof!"\newline

\newline
...\newline

\newline
>>> class JackRussellTerrier(Dog):\newline

\newline
...     def speak(self):\newline

\newline
...         return "Arff!"\newline

\newline
...\newline

\newline
>>> bobo = JackRussellTerrier()\newline

\newline
>>> bobo.walk()\newline

\newline
```}

\begin{enumerate}[label=\Alph*.]
    \item Arff!
    \item \textbf{*walking*} \hfill \textcolor{blue}{\textbf{[Correct]}}
    \item AttributeError: 'JackRussellTerrier' object has no attribute 'walk'
    \item Woof!
\end{enumerate}

\noindent \textbf{Correct Answer:} \textcolor{blue}{\textbf{B}} \\[4pt]
\noindent \textbf{Explanation:}\\
Phân tích tính kế thừa:\newline
1. Lớp `JackRussellTerrier` kế thừa từ lớp cha `Dog` bằng cú pháp `class JackRussellTerrier(Dog)`.\newline
2. Trong lớp con, chỉ có phương thức `speak(self)` được ghi đè (overridden) để trả về `"Arff!"`.\newline
3. Phương thức `walk(self)` không bị ghi đè, do đó thể hiện `bobo` của lớp con thừa hưởng hoàn toàn phương thức `walk()` từ lớp cha `Dog`.\newline
4. Khi gọi `bobo.walk()` trong phiên tương tác REPL (`>>>`), giá trị trả về của hàm là `*walking*` sẽ được hiển thị ra màn hình.
\end{minipage}
\vspace{0.8cm}
\hrule
\vspace{0.5cm}

\begin{minipage}{\textwidth}
\subsection*{Question 14 \quad \footnotesize{[0014_q014_Question_14.png]}}
\noindent \textbf{Which of the following code uses the inheritance feature of Python?}

\begin{enumerate}[label=\Alph*.]
    \item Class Foo: Pass
    \item Class Foo(object): passclass Hoo(object): pass
    \item \textbf{Class Foo: passclass Hoo(Foo): pass} \hfill \textcolor{blue}{\textbf{[Correct]}}
    \item None of the above code.
\end{enumerate}

\noindent \textbf{Correct Answer:} \textcolor{blue}{\textbf{C}} \\[4pt]
\noindent \textbf{Explanation:}\\
Trong Python, tính kế thừa (inheritance) cho phép một lớp con tái sử dụng và mở rộng các thuộc tính, phương thức của lớp cha. Cú pháp kế thừa giữa các lớp là đặt tên lớp cha trong dấu ngoặc đơn sau tên lớp con: `class SubClass(SuperClass):`.\newline
- Trong đáp án C: `Class Foo: passclass Hoo(Foo): pass`, lớp `Hoo` kế thừa từ lớp cha do người dùng định nghĩa là `Foo` thông qua cú pháp `Hoo(Foo)`.\newline
- Đáp án A chỉ định nghĩa một lớp đơn lẻ không kế thừa.\newline
- Đáp án B kế thừa từ `object` (lớp cơ sở mặc định ngầm định của Python), không thể hiện mối quan hệ kế thừa giữa các lớp tự định nghĩa với nhau.\newline
Do đó, đáp án C là đáp án chính xác.
\end{minipage}
\vspace{0.8cm}
\hrule
\vspace{0.5cm}

\begin{minipage}{\textwidth}
\subsection*{Question 15 \quad \footnotesize{[0015_q015_Question_15.png]}}
\noindent \textbf{What will be the value of x when the following statement is executed:\newline

\newline
\newline

\newline
```python\newline

\newline
x = int(98.6)\newline

\newline
```}

\begin{enumerate}[label=\Alph*.]
    \item 100
    \item 6
    \item \textbf{98} \hfill \textcolor{blue}{\textbf{[Correct]}}
    \item 99
\end{enumerate}

\noindent \textbf{Correct Answer:} \textcolor{blue}{\textbf{C}} \\[4pt]
\noindent \textbf{Explanation:}\\
Trong Python, hàm chuyển đổi kiểu dữ liệu `int()` khi nhận một số thực (float) sẽ thực hiện cắt cụt (truncation toward zero), tức là loại bỏ hoàn toàn phần thập phân phía sau dấu chấm mà không làm tròn số (rounding).\newline
Do đó, `int(98.6)` sẽ lược bỏ phần `.6` và trả về số nguyên `98`.\newline
Biến `x` nhận giá trị là `98`.
\end{minipage}
\vspace{0.8cm}
\hrule
\vspace{0.5cm}

\begin{minipage}{\textwidth}
\subsection*{Question 16 \quad \footnotesize{[0016_q016_Question_16.png]}}
\noindent \textbf{What will the following program print out:\newline

\newline
\newline

\newline
```python\newline

\newline
>>> x = 15\newline

\newline
>>> x = x + 10\newline

\newline
>>> print(x)\newline

\newline
```}

\begin{enumerate}[label=\Alph*.]
    \item \textbf{25} \hfill \textcolor{blue}{\textbf{[Correct]}}
    \item 5
    \item 15
    \item "print x"
\end{enumerate}

\noindent \textbf{Correct Answer:} \textcolor{blue}{\textbf{A}} \\[4pt]
\noindent \textbf{Explanation:}\\
Các bước thực hiện lệnh:\newline
1. `x = 15`: Khởi tạo biến `x` và gán giá trị 15.\newline
2. `x = x + 10`: Tính toán vế phải `15 + 10 = 25`, sau đó gán kết quả 25 trở lại cho biến `x`.\newline
3. `print(x)`: In giá trị hiện thời của `x` ra màn hình, kết quả nhận được là `25`.
\end{minipage}
\vspace{0.8cm}
\hrule
\vspace{0.5cm}

\begin{minipage}{\textwidth}
\subsection*{Question 17 \quad \footnotesize{[0017_q017_Question_17.png]}}
\noindent \textbf{What does the following Python sequence print out?\newline

\newline
\newline

\newline
```python\newline

\newline
x = 'From: Using the : character'\newline

\newline
y = re.findall('^F.+:', x)\newline

\newline
print(y)\newline

\newline
```}

\begin{enumerate}[label=\Alph*.]
    \item ^F.+:
    \item \textbf{['From: Using the :']} \hfill \textcolor{blue}{\textbf{[Correct]}}
    \item ['From:']
    \item From:
\end{enumerate}

\noindent \textbf{Correct Answer:} \textcolor{blue}{\textbf{B}} \\[4pt]
\noindent \textbf{Explanation:}\\
Phân tích biểu thức chính quy (Regex) `^F.+:`:\newline
- `^F`: Bắt đầu chuỗi bằng chữ cái `'F'`.\newline
- `.+`: Khớp với một hoặc nhiều ký tự bất kỳ (ngoại trừ ký tự xuống dòng).\newline
- `:`: Kết thúc bằng ký tự hai chấm `:`.\newline
Trong Python Regex, toán tử lặp `+` mặc định là tham lam (greedy), nghĩa là nó sẽ khớp chuỗi dài nhất có thể cho tới dấu hai chấm cuối cùng.\newline
Trong chuỗi `'From: Using the : character'`, có hai dấu hai chấm: một sau `From` và một sau `the `. Do tính tham lam của `.+`, chuỗi khớp được sẽ kéo dài đến dấu hai chấm thứ hai, tạo thành chuỗi `'From: Using the :'`. Hàm `re.findall()` trả về một danh sách chứa các chuỗi con khớp được: `['From: Using the :']`.
\end{minipage}
\vspace{0.8cm}
\hrule
\vspace{0.5cm}

\begin{minipage}{\textwidth}
\subsection*{Question 18 \quad \footnotesize{[0018_q018_Question_18.png]}}
\noindent \textbf{Given a code,\newline

\newline
\newline

\newline
```python\newline

\newline
class People():\newline

\newline
    def __init__(self, name):\newline

\newline
        self.name = name\newline

\newline
\newline

\newline
    def namePrint(self):\newline

\newline
        print(self.name)\newline

\newline
\newline

\newline
person1 = People("Sally")\newline

\newline
person2 = People("Louise")\newline

\newline
person1.namePrint()\newline

\newline
```\newline

\newline
\newline

\newline
What is correct answer?}

\begin{enumerate}[label=\Alph*.]
    \item person1 and person2 are two different instances of the People class.
    \item \textbf{The __init__ method is used to set initial values for attributes.} \hfill \textcolor{blue}{\textbf{[Correct]}}
    \item 'self' is not needed in def namePrint(self):
    \item person2 has a different value for 'name' than person1.
\end{enumerate}

\noindent \textbf{Correct Answer:} \textcolor{blue}{\textbf{B}} \\[4pt]
\noindent \textbf{Explanation:}\\
Trong lập trình hướng đối tượng (OOP) của Python, phương thức `__init__` là phương thức khởi tạo (constructor / initializer) được tự động gọi mỗi khi một đối tượng mới được tạo ra từ một lớp. Vai trò chính của `__init__` là gán và thiết lập các giá trị ban đầu cho các thuộc tính của đối tượng (trong đoạn mã trên là `self.name = name`).\newline
Đây là câu hỏi lý thuyết cốt lõi về vai trò của `__init__` trong chuẩn ngân hàng đề thi FPTU, do đó đáp án chuẩn xác nhất là B: "The __init__ method is used to set initial values for attributes."\newline
(Phương án C sai vì `self` bắt buộc phải có ở tham số đầu tiên của phương thức thực thể; các phương án A và D dù là các nhận xét đúng về mặt sự kiện nhưng không phải là định nghĩa trọng tâm kiến thức được kiểm tra).
\end{minipage}
\vspace{0.8cm}
\hrule
\vspace{0.5cm}

\begin{minipage}{\textwidth}
\subsection*{Question 19 \quad \footnotesize{[0019_q019_Question_19.png]}}
\noindent \textbf{For the following code, which of the following statements is true?\newline

\newline
\newline

\newline
```python\newline

\newline
def printHello():\newline

\newline
    print("Hello")\newline

\newline
\newline

\newline
a = printHello()\newline

\newline
```}

\begin{enumerate}[label=\Alph*.]
    \item \textbf{PrintHello() is a function and a is a variable. None of them are objects.} \hfill \textcolor{blue}{\textbf{[Correct]}}
    \item Both printHello() and a refer to the same object.
    \item PrintHello() and a refer to different objects.
    \item Syntax error! You cannot assign function to a variable in Python.
\end{enumerate}

\noindent \textbf{Correct Answer:} \textcolor{blue}{\textbf{A}} \\[4pt]
\noindent \textbf{Explanation:}\\
Phân tích mã nguồn:\newline
1. `def printHello():` định nghĩa một hàm mang tên `printHello`.\newline
2. `a = printHello()` thực hiện lời gọi hàm `printHello()`. Hàm in ra chữ `Hello` và trả về giá trị mặc định là `None`. Giá trị này được gán cho biến `a`.\newline
- Theo ngân hàng câu hỏi chuẩn FPTU: Đáp án A được chọn nhằm phân biệt khái niệm cơ bản giữa hàm (`function`) và biến lưu trữ (`variable`) so với các đối tượng thực thể của lớp trong OOP.\newline
- Ghi chú chuyên sâu: Trong ngôn ngữ Python chuẩn, mọi thực thể đều là first-class object (hàm là đối tượng kiểu function, `None` là đối tượng kiểu NoneType). Đáp án B sai vì `printHello` (đối tượng hàm) và `a` (giá trị `None`) là hai thực thể hoàn toàn khác nhau. Đáp án D sai vì Python hoàn toàn cho phép gán kết quả trả về của hàm cho một biến.
\end{minipage}
\vspace{0.8cm}
\hrule
\vspace{0.5cm}

\begin{minipage}{\textwidth}
\subsection*{Question 20 \quad \footnotesize{[0020_q020_Question_20.png]}}
\noindent \textbf{Which of the following elements of a mathematical expression in Python is evaluated first?}

\begin{enumerate}[label=\Alph*.]
    \item Addition +
    \item Multiplication *
    \item Subtraction -
    \item \textbf{Parentheses ( )} \hfill \textcolor{blue}{\textbf{[Correct]}}
\end{enumerate}

\noindent \textbf{Correct Answer:} \textcolor{blue}{\textbf{D}} \\[4pt]
\noindent \textbf{Explanation:}\\
Quy tắc thứ tự ưu tiên của các toán tử trong toán học cũng như trong Python tuân theo quy tắc PEMDAS:\newline
1. **P**arentheses `( )` (Dấu ngoặc đơn): Có độ ưu tiên cao nhất, các phép toán nằm trong dấu ngoặc luôn được tính toán đầu tiên.\newline
2. **E**xponentiation `**` (Lũy thừa).\newline
3. **M**ultiplication `*`, **D**ivision `/`, `//`, `%` (Phép nhân và các phép chia).\newline
4. **A**ddition `+`, **S**ubtraction `-` (Phép cộng và trừ).\newline
Do đó, trong các thành phần trên, biểu thức trong dấu ngoặc đơn `Parentheses ( )` luôn được đánh giá đầu tiên.
\end{minipage}
\vspace{0.8cm}
\hrule
\vspace{0.5cm}

\begin{minipage}{\textwidth}
\subsection*{Question 21 \quad \footnotesize{[0021_q021_Question_21.png]}}
\noindent \textbf{Which of the following variables is the "most mnemonic"?}

\begin{enumerate}[label=\Alph*.]
    \item var_173
    \item x1q3z9ocd
    \item xyz
    \item \textbf{days} \hfill \textcolor{blue}{\textbf{[Correct]}}
\end{enumerate}

\noindent \textbf{Correct Answer:} \textcolor{blue}{\textbf{D}} \\[4pt]
\noindent \textbf{Explanation:}\\
Trong lập trình Python (đặc biệt theo giáo trình Python for Everybody của Dr. Charles Severance):\newline
- **Mnemonic variable name** (tên biến có tính gợi nhớ) là tên biến được đặt sao cho bản thân từ ngữ mang ý nghĩa tự nhiên, gợi nhắc trực tiếp cho con người (lập trình viên) biết dữ liệu hoặc mục đích của biến đó đang đại diện cho điều gì.\newline
- Tên biến `days` thể hiện rõ ràng nó lưu trữ thông tin về "số ngày" hoặc danh sách các ngày, rất dễ nhớ và dễ đọc hiểu mã nguồn.\newline
- Ngược lại:\newline
  + `x1q3z9ocd` là chuỗi ký tự ngẫu nhiên, vô nghĩa và cực kỳ khó nhớ.\newline
  + `xyz` là các ký tự chung chung, không truyền tải bất kỳ ngữ cảnh nghiệp vụ nào.\newline
  + `var_173` là tên chung chung (variable 173), không mang tính gợi nhớ về dữ liệu thực tế được lưu trữ.
\end{minipage}
\vspace{0.8cm}
\hrule
\vspace{0.5cm}

\begin{minipage}{\textwidth}
\subsection*{Question 22 \quad \footnotesize{[0022_q022_Question_22.png]}}
\noindent \textbf{What is the difference between the "+" and "*" character in regular expressions?}

\begin{enumerate}[label=\Alph*.]
    \item The "+" matches upper case characters and the "*" matches lowercase characters
    \item The "+" matches the beginning of a line and the "*" matches the end of a line
    \item The "+" matches the actual plus character and the "*" matches any character
    \item \textbf{The "+" matches at least one character and the "*" matches zero or more characters} \hfill \textcolor{blue}{\textbf{[Correct]}}
\end{enumerate}

\noindent \textbf{Correct Answer:} \textcolor{blue}{\textbf{D}} \\[4pt]
\noindent \textbf{Explanation:}\\
Trong biểu thức chính quy (Regular Expressions / module `re` trong Python):\newline
- Ký tự lượng từ `+` (plus quantifier) khớp với **ít nhất một lần** (1 hoặc nhiều lần - 1 or more occurrences / at least one) ký tự hoặc nhóm đứng ngay trước nó.\newline
- Ký tự lượng từ `*` (asterisk/star quantifier) khớp với **không hoặc nhiều lần** (0 or more occurrences) ký tự hoặc nhóm đứng ngay trước nó.\newline
- Phân tích các lựa chọn sai:\newline
  + A sai: Regular expressions dùng các lớp ký tự như `[A-Z]` để khớp chữ hoa và `[a-z]` cho chữ thường.\newline
  + B sai: Ký tự khớp đầu dòng là `^` (caret) và cuối dòng là `$` (dollar).\newline
  + C sai: Để khớp ký tự dấu cộng thực tế, ta phải dùng escape ký tự `\+`.
\end{minipage}
\vspace{0.8cm}
\hrule
\vspace{0.5cm}

\begin{minipage}{\textwidth}
\subsection*{Question 23 \quad \footnotesize{[0023_q023_Question_23.png]}}
\noindent \textbf{What is a correct syntax to return the third character in a string?}

\begin{enumerate}[label=\Alph*.]
    \item \textbf{x = "Hello"[2]} \hfill \textcolor{blue}{\textbf{[Correct]}}
    \item x = sub("Hello",0,2)
    \item x = "Hello".sub(0,2)
    \item x = sub "Hello"
\end{enumerate}

\noindent \textbf{Correct Answer:} \textcolor{blue}{\textbf{A}} \\[4pt]
\noindent \textbf{Explanation:}\\
Trong Python, chuỗi ký tự (`str`) là một chuỗi tuần tự (sequence) hỗ trợ truy cập phần tử thông qua chỉ mục (indexing) bằng cặp dấu ngoặc vuông `[]`:\newline
- Hệ thống chỉ số trong Python bắt đầu từ `0` (zero-based indexing):\newline
  + Ký tự thứ nhất (1st) là index 0: `"Hello"[0]` -> `'H'`\newline
  + Ký tự thứ hai (2nd) là index 1: `"Hello"[1]` -> `'e'`\newline
  + Ký tự thứ ba (3rd) là index 2: `"Hello"[2]` -> `'l'`\newline
- Do đó, cú pháp chính xác để lấy ký tự thứ 3 là `x = "Hello"[2]`.\newline
- Các lựa chọn khác sai vì Python không có hàm hoặc phương thức tích hợp sẵn nào tên là `sub()` trên đối tượng chuỗi (nếu gọi `"Hello".sub(...)` sẽ gặp `AttributeError`).
\end{minipage}
\vspace{0.8cm}
\hrule
\vspace{0.5cm}

\begin{minipage}{\textwidth}
\subsection*{Question 24 \quad \footnotesize{[0024_q024_Question_24.png]}}
\noindent \textbf{Which of the following will result in an error?}

\begin{enumerate}[label=\Alph*.]
    \item print(str1[2])
    \item \textbf{str1[1]="x"} \hfill \textcolor{blue}{\textbf{[Correct]}}
    \item print(str1[0:9])
    \item Both (b) and (c)
\end{enumerate}

\noindent \textbf{Correct Answer:} \textcolor{blue}{\textbf{B}} \\[4pt]
\noindent \textbf{Explanation:}\\
Trong Python, chuỗi (`str`) là kiểu dữ liệu **bất biến** (immutable):\newline
- Một khi chuỗi đã được khởi tạo, ta không thể thay đổi giá trị của bất kỳ ký tự nào bên trong chuỗi thông qua chỉ số (item assignment).\newline
- Câu lệnh `str1[1] = "x"` sẽ ngay lập tức gây ra ngoại lệ: `TypeError: 'str' object does not support item assignment`.\newline
- Phân tích các đáp án khác:\newline
  + `print(str1[2])` là cú pháp lấy ký tự tại index 2, hoàn toàn hợp lệ nếu độ dài chuỗi > 2.\newline
  + `print(str1[0:9])` là phép cắt lát (slicing). Trong Python, slicing ngoài phạm vi độ dài (out-of-bounds) không bao giờ ném lỗi `IndexError`, mà chỉ cắt tới cuối chuỗi hiện có.
\end{minipage}
\vspace{0.8cm}
\hrule
\vspace{0.5cm}

\begin{minipage}{\textwidth}
\subsection*{Question 25 \quad \footnotesize{[0025_q025_Question_25.png]}}
\noindent \textbf{How do you start writing a for loop in Python?}

\begin{enumerate}[label=\Alph*.]
    \item for x = 1 to y:
    \item for x > y:
    \item for each x in y:
    \item \textbf{for x in y:} \hfill \textcolor{blue}{\textbf{[Correct]}}
\end{enumerate}

\noindent \textbf{Correct Answer:} \textcolor{blue}{\textbf{D}} \\[4pt]
\noindent \textbf{Explanation:}\\
Cú pháp chuẩn để bắt đầu một vòng lặp `for` trong Python là:\newline
```python\newline
for <biến_lặp> in <iterable>:\newline
```\newline
- Lệnh này duyệt qua từng phần tử trong cấu trúc dữ liệu tuần hoàn `y` (như list, tuple, string, range, v.v.) và gán lần lượt từng giá trị cho biến `x`.\newline
- Phân tích các lựa chọn sai:\newline
  + A (`for x = 1 to y:`): Là cú pháp của các ngôn ngữ như BASIC, Pascal; không hợp lệ trong Python.\newline
  + B (`for x > y:`): Cú pháp sai, vòng lặp điều kiện sử dụng từ khóa `while x > y:`.\newline
  + C (`for each x in y:`): Là từ khóa của các ngôn ngữ như Visual Basic hay PHP; Python không có từ khóa `each`.
\end{minipage}
\vspace{0.8cm}
\hrule
\vspace{0.5cm}

\begin{minipage}{\textwidth}
\subsection*{Question 26 \quad \footnotesize{[0026_q026_Question_26.png]}}
\noindent \textbf{What does the following Python code print out?\newline

\newline
\newline

\newline
```python\newline

\newline
a = [1, 2, 3, 8]\newline

\newline
b = [4, 5, 6]\newline

\newline
c = a + b\newline

\newline
print(c)\newline

\newline
```}

\begin{enumerate}[label=\Alph*.]
    \item 28
    \item \textbf{[1, 2, 3, 8, 4, 5, 6]} \hfill \textcolor{blue}{\textbf{[Correct]}}
    \item It raises an exception
    \item 7
\end{enumerate}

\noindent \textbf{Correct Answer:} \textcolor{blue}{\textbf{B}} \\[4pt]
\noindent \textbf{Explanation:}\\
Từng bước thực thi của đoạn mã:\newline
1. Khởi tạo danh sách `a = [1, 2, 3, 8]`.\newline
2. Khởi tạo danh sách `b = [4, 5, 6]`.\newline
3. Toán tử `+` khi thực hiện giữa hai danh sách trong Python đại diện cho phép nối danh sách (list concatenation), tạo ra một danh sách mới chứa tất cả các phần tử của danh sách thứ nhất theo sau là các phần tử của danh sách thứ hai.\newline
   -> `c = [1, 2, 3, 8] + [4, 5, 6] = [1, 2, 3, 8, 4, 5, 6]`.\newline
4. Lệnh `print(c)` in kết quả ra màn hình: `[1, 2, 3, 8, 4, 5, 6]`.\newline
- Lưu ý: Toán tử `+` không cộng từng phần tử toán học, cũng không tính tổng số phần tử nên A (28) và D (7) là sai.
\end{minipage}
\vspace{0.8cm}
\hrule
\vspace{0.5cm}

\begin{minipage}{\textwidth}
\subsection*{Question 27 \quad \footnotesize{[0027_q027_Question_27.png]}}
\noindent \textbf{You look at the following text:\newline

\newline
\newline

\newline
```python\newline

\newline
if x == 6 :\newline

\newline
    print('Is 6')\newline

\newline
    print('Is Still 6')\newline

\newline
    print('Third 6')\newline

\newline
```\newline

\newline
\newline

\newline
It looks perfect but Python is giving you an 'Indentation Error' on the second print statement. What is the most likely reason?}

\begin{enumerate}[label=\Alph*.]
    \item In order to make humans feel inadequate, Python randomly emits 'Indentation Errors' on perfectly good code - after about an hour the error will just go away without any changes to your program
    \item Python has reached its limit on the largest Python program that can be run
    \item Python thinks 'Still' is a mis-spelled word in the string
    \item \textbf{You have mixed tabs and spaces in the file} \hfill \textcolor{blue}{\textbf{[Correct]}}
\end{enumerate}

\noindent \textbf{Correct Answer:} \textcolor{blue}{\textbf{D}} \\[4pt]
\noindent \textbf{Explanation:}\\
Giải thích chi tiết về lỗi thụt đầu dòng trong Python:\newline
- Python sử dụng mức độ thụt đầu dòng (indentation) để xác định các khối lệnh (code blocks).\newline
- Khi nhìn trong một số trình soạn thảo văn bản, 1 ký tự Tab (`\t`) có thể hiển thị độ rộng tương đương 4 dấu cách (`    `), khiến mắt người nhìn thấy các dòng hoàn toàn thẳng hàng.\newline
- Tuy nhiên, trình thông dịch Python coi Tab và Space là hai ký tự hoàn toàn khác nhau. Việc trộn lẫn cả phím Tab lẫn phím Space để thụt đầu dòng trong cùng một khối mã là nguyên nhân hàng đầu gây ra lỗi `IndentationError` (hoặc cụ thể hơn là `TabError: inconsistent use of tabs and spaces in indentation`).\newline
- Khuyến nghị theo chuẩn PEP 8 là luôn dùng 4 dấu cách (spaces) cho mỗi cấp độ thụt dòng và không dùng tab.
\end{minipage}
\vspace{0.8cm}
\hrule
\vspace{0.5cm}

\begin{minipage}{\textwidth}
\subsection*{Question 28 \quad \footnotesize{[0028_q028_Question_28.png]}}
\noindent \textbf{Which of the parts of a computer actually executes the program instructions?}

\begin{enumerate}[label=\Alph*.]
    \item RAM
    \item \textbf{CPU} \hfill \textcolor{blue}{\textbf{[Correct]}}
    \item Input/Output Devices
    \item Secondary Memory
\end{enumerate}

\noindent \textbf{Correct Answer:} \textcolor{blue}{\textbf{B}} \\[4pt]
\noindent \textbf{Explanation:}\\
Trong kiến trúc máy tính Von Neumann căn bản:\newline
- **CPU** (Central Processing Unit - Bộ xử lý trung tâm) là thành phần phần cứng cốt lõi thực hiện chu trình Fetch-Decode-Execute (Lấy lệnh, Giải mã và Thực thi) các chỉ thị của chương trình.\newline
- Phân tích các thành phần khác:\newline
  + **RAM** (Random Access Memory - Main Memory): Lưu trữ tạm thời các chương trình và dữ liệu mà CPU cần xử lý ngay, nhưng RAM không tự thực thi lệnh.\newline
  + **Secondary Memory** (Bộ nhớ thứ cấp như ổ cứng HDD, SSD, thẻ nhớ): Lưu trữ thông tin lâu dài (không mất đi khi tắt máy).\newline
  + **Input/Output Devices** (Thiết bị vào/ra như bàn phím, màn hình, chuột): Giúp máy tính giao tiếp với người dùng và môi trường bên ngoài.
\end{minipage}
\vspace{0.8cm}
\hrule
\vspace{0.5cm}

\begin{minipage}{\textwidth}
\subsection*{Question 29 \quad \footnotesize{[0029_q029_Question_29.png]}}
\noindent \textbf{Which of the following slicing operations will produce the list [12, 3]?\newline

\newline
\newline

\newline
```python\newline

\newline
t = [9, 41, 12, 3, 74, 15]\newline

\newline
```}

\begin{enumerate}[label=\Alph*.]
    \item t[1:3]
    \item t[2:2]
    \item t[:]
    \item \textbf{t[2:4]} \hfill \textcolor{blue}{\textbf{[Correct]}}
\end{enumerate}

\noindent \textbf{Correct Answer:} \textcolor{blue}{\textbf{D}} \\[4pt]
\noindent \textbf{Explanation:}\\
Phép cắt lát danh sách (list slicing) trong Python có cú pháp `t[start:stop]`, lấy các phần tử từ chỉ mục `start` đến `stop - 1`:\newline
- Xác định chỉ số của danh sách `t`:\newline
  + `t[0]` = 9\newline
  + `t[1]` = 41\newline
  + `t[2]` = 12\newline
  + `t[3]` = 3\newline
  + `t[4]` = 74\newline
  + `t[5]` = 15\newline
- Mục tiêu là tạo ra danh sách con `[12, 3]`, gồm hai phần tử ở index 2 và index 3.\newline
- Do đó, ta cần `start = 2` và `stop = 4` (`t[2:4]`), lấy các phần tử tại vị trí 2 và 3 -> `[12, 3]`.\newline
- Kiểm tra các phương án khác:\newline
  + `t[1:3]` lấy index 1 và 2 -> `[41, 12]`\newline
  + `t[2:2]` lát cắt có start == stop trả về danh sách rỗng `[]`\newline
  + `t[:]` sao chép nông toàn bộ danh sách `[9, 41, 12, 3, 74, 15]`.
\end{minipage}
\vspace{0.8cm}
\hrule
\vspace{0.5cm}

\begin{minipage}{\textwidth}
\subsection*{Question 30 \quad \footnotesize{[0030_q030_Question_30.png]}}
\noindent \textbf{Python scripts (files) have names that end with:}

\begin{enumerate}[label=\Alph*.]
    \item .exe
    \item \textbf{.py} \hfill \textcolor{blue}{\textbf{[Correct]}}
    \item .png
    \item .doc
\end{enumerate}

\noindent \textbf{Correct Answer:} \textcolor{blue}{\textbf{B}} \\[4pt]
\noindent \textbf{Explanation:}\\
Theo quy ước chuẩn của ngôn ngữ Python:\newline
- Các tệp kịch bản mã nguồn Python luôn có đuôi mở rộng là `.py` (ví dụ: `script.py`, `app.py`).\newline
- Phân tích các phần mở rộng khác:\newline
  + `.exe`: Tệp thực thi nhị phân độc lập trên hệ điều hành Microsoft Windows.\newline
  + `.png`: Tệp hình ảnh đồ họa định dạng Portable Network Graphics.\newline
  + `.doc`: Định dạng tệp tài liệu văn bản của ứng dụng Microsoft Word.
\end{minipage}
\vspace{0.8cm}
\hrule
\vspace{0.5cm}

\begin{minipage}{\textwidth}
\subsection*{Question 31 \quad \footnotesize{[0031_q031_Question_31.png]}}
\noindent \textbf{Which of the following are valid ways to specify the string literal foo'bar in Python:}

\begin{enumerate}[label=\Alph*.]
    \item \textbf{"foo'bar"} \hfill \textcolor{blue}{\textbf{[Correct]}}
    \item 'foo"bar'
    \item """"foo'bar""""
    \item 'foo'bar'
\end{enumerate}

\noindent \textbf{Correct Answer:} \textcolor{blue}{\textbf{A}} \\[4pt]
\noindent \textbf{Explanation:}\\
Trong Python, chuỗi ký tự có thể được khai báo bằng cặp dấu ngoặc đơn `'...'` hoặc cặp dấu ngoặc kép `"..."`:\newline
- Để chuỗi biểu diễn chứa ký tự nháy đơn `'` mà không cần sử dụng ký tự thoát (escape backslash `\'`), ta có thể bọc toàn bộ chuỗi trong cặp dấu nháy kép: `"foo'bar"`. Giá trị chuỗi thu được chính xác là `foo'bar`.\newline
- Phân tích các lựa chọn còn lại:\newline
  + B: `'foo"bar'` tạo ra chuỗi `foo"bar` (chứa dấu nháy kép chứ không phải nháy đơn).\newline
  + C: `""""foo'bar""""` có 4 dấu ngoặc kép liên tiếp ở hai đầu, sai cú pháp chuỗi nhiều dòng (triple-quotes chỉ gồm đúng 3 dấu `"""`).\newline
  + D: `'foo'bar'` bị lỗi `SyntaxError: invalid syntax` vì dấu nháy đơn thứ hai làm kết thúc chuỗi sớm, đoạn `bar'` phía sau không hợp lệ.
\end{minipage}
\vspace{0.8cm}
\hrule
\vspace{0.5cm}

\begin{minipage}{\textwidth}
\subsection*{Question 32 \quad \footnotesize{[0032_q032_Question_32.png]}}
\noindent \textbf{Which of the following statement about Python is not correct?}

\begin{enumerate}[label=\Alph*.]
    \item The python ecosystem is very rich and provides easy to use tools for data science
    \item \textbf{Due to its proprietary nature, database access from python is not available for many database} \hfill \textcolor{blue}{\textbf{[Correct]}}
    \item There are libraries and APIs available to access many of the popular databases from python
    \item Python is a popular scripting language for connecting and accessing databases
\end{enumerate}

\noindent \textbf{Correct Answer:} \textcolor{blue}{\textbf{B}} \\[4pt]
\noindent \textbf{Explanation:}\\
Đánh giá tính chính xác của các phát biểu về Python:\newline
- Đáp án B là phát biểu **sai (not correct)** vì:\newline
  + Python là phần mềm nguồn mở (open-source) dưới giấy phép Python Software Foundation License, hoàn toàn **không mang tính độc quyền** ("proprietary").\newline
  + Python có khả năng kết nối cơ sở dữ liệu cực kỳ mạnh mẽ và toàn diện (chuẩn PEP 249 Python DB-API, các driver cho SQLite, PostgreSQL, MySQL, Oracle, MongoDB, v.v.), chứ không hề bị hạn chế kết nối cơ sở dữ liệu.\newline
- Các khẳng định A, C, D đều hoàn toàn chính xác về thế mạnh xử lý dữ liệu và hệ sinh thái thư viện kết nối cơ sở dữ liệu của Python.
\end{minipage}
\vspace{0.8cm}
\hrule
\vspace{0.5cm}

\begin{minipage}{\textwidth}
\subsection*{Question 33 \quad \footnotesize{[0033_q033_Question_33.png]}}
\noindent \textbf{How do you start writing an if statement in Python?}

\begin{enumerate}[label=\Alph*.]
    \item \textbf{if x > y:} \hfill \textcolor{blue}{\textbf{[Correct]}}
    \item if (x > y)
    \item if x > y then:
    \item x > y?:
\end{enumerate}

\noindent \textbf{Correct Answer:} \textcolor{blue}{\textbf{A}} \\[4pt]
\noindent \textbf{Explanation:}\\
Cú pháp câu lệnh điều kiện `if` trong Python:\newline
- Cấu trúc chuẩn: `if <biểu_thức_điều_kiện>:`\newline
- Cuối dòng khai báo mệnh đề `if` bắt buộc phải có dấu hai chấm `:` để mở đầu cho khối mã lệnh được thụt lề thực thi phía sau.\newline
- Phân tích các phương án sai:\newline
  + B (`if (x > y)`): Thiếu dấu `:` ở cuối dòng -> gây ra lỗi `SyntaxError: expected ':'`.\newline
  + C (`if x > y then:`): Python không sử dụng từ khóa `then` (như Pascal, Ruby, hay Basic).\newline
  + D (`x > y?:`): Cú pháp không hợp lệ trong Python (toán tử ba ngôi trong Python được viết dưới dạng `a if condition else b`).
\end{minipage}
\vspace{0.8cm}
\hrule
\vspace{0.5cm}

\begin{minipage}{\textwidth}
\subsection*{Question 34 \quad \footnotesize{[0034_q034_Question_34.png]}}
\noindent \textbf{```python\newline

\newline
class Dog:\newline

\newline
    def __init__(self, name, age):\newline

\newline
        self.name = name\newline

\newline
        self.age = age\newline

\newline
\newline

\newline
class JackRussellTerrier(Dog):\newline

\newline
    pass\newline

\newline
\newline

\newline
class Dachshund(Dog):\newline

\newline
    pass\newline

\newline
\newline

\newline
class Bulldog(Dog):\newline

\newline
    pass\newline

\newline
\newline

\newline
miles = JackRussellTerrier("Miles", 4)\newline

\newline
buddy = Dachshund("Buddy", 9)\newline

\newline
jack = Bulldog("Jack", 3)\newline

\newline
jim = Bulldog("Jim", 5)\newline

\newline
```\newline

\newline
\newline

\newline
Given the above code snippet, which of the following REPL outputs are correct? (Select all that apply.)}

\begin{enumerate}[label=\Alph*.]
    \item \textbf{isinstance(jack, Dog)} \hfill \textcolor{blue}{\textbf{[Correct]}}
    \item \textbf{isinstance(miles, Dog)} \hfill \textcolor{blue}{\textbf{[Correct]}}
    \item isinstance(jack, Dachshund)
    \item isinstance(miles, Bulldog)
\end{enumerate}

\noindent \textbf{Correct Answer:} \textcolor{blue}{\textbf{A, B}} \\[4pt]
\noindent \textbf{Explanation:}\\
Hàm built-in `isinstance(object, classinfo)` trong Python kiểm tra xem đối tượng `object` có phải là thể hiện của lớp `classinfo` hoặc bất kỳ lớp cha nào mà lớp đó kế thừa hay không:\newline
- Cấu trúc phân cấp kế thừa trong đoạn mã:\newline
  + `JackRussellTerrier` là lớp con kế thừa từ `Dog`.\newline
  + `Dachshund` là lớp con kế thừa từ `Dog`.\newline
  + `Bulldog` là lớp con kế thừa từ `Dog`.\newline
- Phân tích từng biểu thức:\newline
  + **A. `isinstance(jack, Dog)`**: `jack` là thể hiện của `Bulldog`, mà `Bulldog` kế thừa từ `Dog`. Do đó biểu thức trả về `True` -> **ĐÚNG**.\newline
  + **B. `isinstance(miles, Dog)`**: `miles` là thể hiện của `JackRussellTerrier`, mà `JackRussellTerrier` kế thừa từ `Dog`. Do đó biểu thức trả về `True` -> **ĐÚNG**.\newline
  + **C. `isinstance(jack, Dachshund)`**: `jack` là đối tượng `Bulldog`, hoàn toàn không liên quan đến nhánh kế thừa `Dachshund` -> trả về `False`.\newline
  + **D. `isinstance(miles, Bulldog)`**: `miles` là đối tượng `JackRussellTerrier`, không liên quan đến nhánh `Bulldog` -> trả về `False`.\newline
- Do đó, hai biểu thức A và B là các biểu thức cho kết quả đúng (`True`).
\end{minipage}
\vspace{0.8cm}
\hrule
\vspace{0.5cm}

\begin{minipage}{\textwidth}
\subsection*{Question 35 \quad \footnotesize{[0035_q035_Question_35.png]}}
\noindent \textbf{Which of the following functions do you use to write data in the binary format?}

\begin{enumerate}[label=\Alph*.]
    \item write
    \item output
    \item \textbf{dump} \hfill \textcolor{blue}{\textbf{[Correct]}}
    \item send
\end{enumerate}

\noindent \textbf{Correct Answer:} \textcolor{blue}{\textbf{C}} \\[4pt]
\noindent \textbf{Explanation:}\\
Trong lập trình Python, để ghi tuần tự hóa dữ liệu đối tượng phức tạp thành định dạng nhị phân (binary serialization / pickling):\newline
- Module `pickle` cung cấp hàm `pickle.dump(obj, file)` dùng để tuần tự hóa một đối tượng và ghi trực tiếp vào một tệp tin mở ở chế độ nhị phân (`'wb'`).\newline
- Để đọc dữ liệu nhị phân này ngược lại thành đối tượng Python, ta dùng hàm `pickle.load(file)`.\newline
- Phân tích các lựa chọn khác:\newline
  + `write`: Là phương thức của đối tượng tệp (`file.write()`), nhận vào chuỗi hoặc bytes đơn thuần chứ không phải là hàm chuyên dụng để serialize dữ liệu đối tượng sang binary format.\newline
  + `output`, `send`: Không phải là hàm chuẩn trong thư viện chuẩn Python dùng cho việc ghi dữ liệu nhị phân đối tượng.
\end{minipage}
\vspace{0.8cm}
\hrule
\vspace{0.5cm}

\begin{minipage}{\textwidth}
\subsection*{Question 36 \quad \footnotesize{[0036_q036_Question_36.png]}}
\noindent \textbf{In Python, a class is .... for a concrete object.}

\begin{enumerate}[label=\Alph*.]
    \item nuisance
    \item instance
    \item \textbf{blueprint} \hfill \textcolor{blue}{\textbf{[Correct]}}
    \item distraction
\end{enumerate}

\noindent \textbf{Correct Answer:} \textcolor{blue}{\textbf{C}} \\[4pt]
\noindent \textbf{Explanation:}\\
Trong lập trình hướng đối tượng (OOP) và ngôn ngữ Python:\newline
- Một lớp (`class`) được định nghĩa là một "bản thiết kế" (blueprint) hoặc khuôn mẫu mô tả các thuộc tính (dữ liệu) và hành vi (phương thức) mà các đối tượng thuộc lớp đó sẽ sở hữu.\newline
- Một đối tượng cụ thể (`concrete object`) được tạo ra từ bản thiết kế này được gọi là một thực thể (`instance`).\newline
- Do đó, câu hoàn chỉnh mang đúng ngữ nghĩa chuẩn là: *"In Python, a class is blueprint for a concrete object."*\newline
- Các từ `nuisance` (sự phiền toái) và `distraction` (sự xao nhãng) hoàn toàn vô nghĩa trong ngữ cảnh kỹ thuật.
\end{minipage}
\vspace{0.8cm}
\hrule
\vspace{0.5cm}

\begin{minipage}{\textwidth}
\subsection*{Question 37 \quad \footnotesize{[0037_q037_Question_37.png]}}
\noindent \textbf{Which method can be used to return a string in upper case letters?}

\begin{enumerate}[label=\Alph*.]
    \item \textbf{upper()} \hfill \textcolor{blue}{\textbf{[Correct]}}
    \item uppercase()
    \item toUpperCase()
    \item upperCase()
\end{enumerate}

\noindent \textbf{Correct Answer:} \textcolor{blue}{\textbf{A}} \\[4pt]
\noindent \textbf{Explanation:}\\
Trong Python, kiểu dữ liệu chuỗi (`str`) cung cấp phương thức tích hợp sẵn `.upper()`:\newline
- Cú pháp: `string.upper()` trả về một chuỗi mới với toàn bộ các chữ cái được chuyển đổi thành chữ in hoa.\newline
- Ví dụ: `'python'.upper()` -> `'PYTHON'`.\newline
- Phân tích các đáp án sai:\newline
  + `toUpperCase()`: Là phương thức của ngôn ngữ Java và JavaScript, không tồn tại trong Python (nếu gọi sẽ sinh ra ngoại lệ `AttributeError: 'str' object has no attribute 'toUpperCase'`).\newline
  + `uppercase()`, `upperCase()`: Không phải là các phương thức của kiểu chuỗi trong Python.
\end{minipage}
\vspace{0.8cm}
\hrule
\vspace{0.5cm}

\begin{minipage}{\textwidth}
\subsection*{Question 38 \quad \footnotesize{[0038_q038_Question_38.png]}}
\noindent \textbf{If you a class is derived from two different classes, it's called _______}

\begin{enumerate}[label=\Alph*.]
    \item Multilevel Inheritance
    \item \textbf{Multiple Inheritance} \hfill \textcolor{blue}{\textbf{[Correct]}}
    \item Hierarchical Inheritance
    \item Python Inheritance
\end{enumerate}

\noindent \textbf{Correct Answer:} \textcolor{blue}{\textbf{B}} \\[4pt]
\noindent \textbf{Explanation:}\\
Trong lập trình hướng đối tượng (OOP) và Python:\newline
- **Multiple Inheritance** (Đa kế thừa) là thuật ngữ chỉ việc một lớp con được kế thừa trực tiếp từ hai hoặc nhiều lớp cha khác nhau (ví dụ: `class C(A, B):`). Python là một trong những ngôn ngữ hỗ trợ đầy đủ đa kế thừa thông qua cơ chế MRO (Method Resolution Order).\newline
- Phân biệt với các loại kế thừa khác:\newline
  + **Multilevel Inheritance** (Kế thừa nhiều cấp): Một lớp kế thừa từ lớp con của một lớp khác (ví dụ: C kế thừa B, B kế thừa A).\newline
  + **Hierarchical Inheritance** (Kế thừa thứ bậc): Nhiều lớp con cùng kế thừa từ một lớp cha duy nhất (ví dụ: B và C đều kế thừa từ A).
\end{minipage}
\vspace{0.8cm}
\hrule
\vspace{0.5cm}

\begin{minipage}{\textwidth}
\subsection*{Question 39 \quad \footnotesize{[0039_q039_Question_39.png]}}
\noindent \textbf{You have the following dictionary definition:\newline

\newline
\newline

\newline
```python\newline

\newline
d = 'foo': 100, 'bar': 200, 'baz': 300\newline

\newline
```\newline

\newline
\newline

\newline
What method call will delete the entry whose value is 200?}

\begin{enumerate}[label=\Alph*.]
    \item \textbf{d.pop("bar")} \hfill \textcolor{blue}{\textbf{[Correct]}}
    \item d.remove("bar")
    \item d.remove("200")
    \item d.pop("200")
\end{enumerate}

\noindent \textbf{Correct Answer:} \textcolor{blue}{\textbf{A}} \\[4pt]
\noindent \textbf{Explanation:}\\
Trong từ điển Python (`dict`), các phần tử được lưu trữ dưới dạng các cặp `key: value`:\newline
- Trong dictionary `d = {'foo': 100, 'bar': 200, 'baz': 300}`, phần tử có giá trị `200` tương ứng với khóa (key) là `'bar'`.\newline
- Phương thức chuẩn của `dict` để xóa một phần tử theo khóa (key) và trả về giá trị của phần tử đó là `d.pop(key)`. Do đó câu lệnh chính xác là `d.pop("bar")`.\newline
- Phân tích các phương án sai:\newline
  + `d.remove(...)`: Đối tượng dictionary trong Python không hề có phương thức `remove()` (phương thức này thuộc về kiểu `list` hoặc `set`). Gọi `d.remove()` sẽ bị lỗi `AttributeError: 'dict' object has no attribute 'remove'`.\newline
  + `d.pop("200")`: Phương thức `pop()` nhận tham số là khóa (key) chứ không phải giá trị. Trong `d` không tồn tại key nào là chuỗi `"200"`, nên lệnh này sẽ ném ngoại lệ `KeyError: '200'`.
\end{minipage}
\vspace{0.8cm}
\hrule
\vspace{0.5cm}

\begin{minipage}{\textwidth}
\subsection*{Question 40 \quad \footnotesize{[0040_q040_Question_40.png]}}
\noindent \textbf{What does the following Python code print out?\newline

\newline
\newline

\newline
```python\newline

\newline
a = [1, 2, 3]\newline

\newline
b = [4, 5, 6]\newline

\newline
c = a + b\newline

\newline
print len(c)\newline

\newline
```}

\begin{enumerate}[label=\Alph*.]
    \item 21
    \item [4, 5, 6]
    \item 15
    \item \textbf{6} \hfill \textcolor{blue}{\textbf{[Correct]}}
\end{enumerate}

\noindent \textbf{Correct Answer:} \textcolor{blue}{\textbf{D}} \\[4pt]
\noindent \textbf{Explanation:}\\
Từng bước thực thi của đoạn mã:\newline
1. Khởi tạo danh sách `a = [1, 2, 3]` gồm 3 phần tử.\newline
2. Khởi tạo danh sách `b = [4, 5, 6]` gồm 3 phần tử.\newline
3. Toán tử `+` thực hiện phép nối hai danh sách: `c = a + b = [1, 2, 3, 4, 5, 6]`.\newline
4. Hàm `len(c)` đếm tổng số lượng phần tử hiện có trong danh sách `c`. Danh sách `c` có 6 phần tử, vì vậy `len(c)` trả về số nguyên `6`.\newline
5. Lệnh in đưa giá trị `6` ra màn hình.\newline
*(Ghi chú: Cú pháp `print len(c)` là cú pháp câu lệnh print trong Python 2; trong Python 3 được viết dưới dạng hàm `print(len(c))`, kết quả hiển thị vẫn là `6`)*.
\end{minipage}
\vspace{0.8cm}
\hrule
\vspace{0.5cm}

\begin{minipage}{\textwidth}
\subsection*{Question 41 \quad \footnotesize{[0041_q041_Question_41.png]}}
\noindent \textbf{What Python function would you use if you wanted to prompt the user for a file name to open?}

\begin{enumerate}[label=\Alph*.]
    \item file_input()
    \item read()
    \item \textbf{input()} \hfill \textcolor{blue}{\textbf{[Correct]}}
    \item cin
\end{enumerate}

\noindent \textbf{Correct Answer:} \textcolor{blue}{\textbf{C}} \\[4pt]
\noindent \textbf{Explanation:}\\
Trong Python, hàm tích hợp sẵn (built-in function) `input()` được sử dụng để hiển thị một lời nhắc (prompt) ra màn hình và chờ người dùng nhập vào một dòng văn bản từ bàn phím dưới dạng chuỗi (string).\newline
- Cú pháp thông dụng: `filename = input('Enter file name: ')`\newline
- Sau khi nhận được tên tệp từ người dùng, lập trình viên có thể mở tệp đó bằng hàm `open(filename)`.\newline
- Phân tích các đáp án khác:\newline
  + A. `file_input()`: Không phải là hàm chuẩn trong Python để nhận nhập liệu từ người dùng (trong thư viện chuẩn có module `fileinput` để duyệt qua các dòng từ nhiều file, nhưng không phải hàm prompt).\newline
  + B. `read()`: Là phương thức của đối tượng file (file object) dùng để đọc dữ liệu từ tệp tin đã mở, không dùng để nhận nhập liệu từ bàn phím.\newline
  + D. `cin`: Là đối tượng nhập dữ liệu luồng chuẩn trong ngôn ngữ C++ (`std::cin`), không tồn tại trong Python.
\end{minipage}
\vspace{0.8cm}
\hrule
\vspace{0.5cm}

\begin{minipage}{\textwidth}
\subsection*{Question 42 \quad \footnotesize{[0042_q042_Question_42.png]}}
\noindent \textbf{Which of these collections defines a SET?}

\begin{enumerate}[label=\Alph*.]
    \item "name": "apple", "color": "green"
    \item ["apple", "banana", "cherry"]
    \item \textbf{"apple", "banana", "cherry"} \hfill \textcolor{blue}{\textbf{[Correct]}}
    \item ("apple", "banana", "cherry")
\end{enumerate}

\noindent \textbf{Correct Answer:} \textcolor{blue}{\textbf{C}} \\[4pt]
\noindent \textbf{Explanation:}\\
Trong Python, tập hợp (`set`) là một kiểu dữ liệu chứa các phần tử không trùng lặp và không có thứ tự, được định nghĩa bằng cặp dấu ngoặc nhọn `{}` chứa các phần tử phân cách bởi dấu phẩy: `{"apple", "banana", "cherry"}`.\newline
- Phân tích chi tiết các phương án:\newline
  + A. `{"name": "apple", "color": "green"}`: Đây là kiểu từ điển (`dict` - Dictionary) vì chứa các cặp khóa - giá trị (`key: value`).\newline
  + B. `["apple", "banana", "cherry"]`: Đây là kiểu danh sách (`list` - List) được định nghĩa bằng cặp ngoặc vuông `[]`.\newline
  + C. `{"apple", "banana", "cherry"}`: Đây chính xác là kiểu tập hợp (`set`).\newline
  + D. `("apple", "banana", "cherry")`: Đây là kiểu bộ giá trị (`tuple` - Tuple) được định nghĩa bằng cặp ngoặc tròn `()`.\newline
- Lưu ý: Để tạo một tập hợp rỗng, ta phải dùng `set()`, vì `{}` mặc định sẽ tạo ra một dictionary rỗng.
\end{minipage}
\vspace{0.8cm}
\hrule
\vspace{0.5cm}

\begin{minipage}{\textwidth}
\subsection*{Question 43 \quad \footnotesize{[0043_q043_Question_43.png]}}
\noindent \textbf{Find the output of the given Python program\newline

\newline
\newline

\newline
```python\newline

\newline
a = 25\newline

\newline
if a < 15:\newline

\newline
    print("Hi")\newline

\newline
if a <= 30:\newline

\newline
    print("Hello")\newline

\newline
else:\newline

\newline
    print("Know Program")\newline

\newline
```}

\begin{enumerate}[label=\Alph*.]
    \item Hi
    \item \textbf{Hello} \hfill \textcolor{blue}{\textbf{[Correct]}}
    \item Hi
    \item Hi
\end{enumerate}

\noindent \textbf{Correct Answer:} \textcolor{blue}{\textbf{B}} \\[4pt]
\noindent \textbf{Explanation:}\\
Từng bước thực thi chương trình:\newline
1. Gán giá trị biến `a = 25`.\newline
2. Cấu trúc rẽ nhánh thứ nhất:\newline
   `if a < 15:`\newline
   Điều kiện `25 < 15` trả về `False`, do đó lệnh `print("Hi")` không được thực hiện.\newline
3. Cấu trúc rẽ nhánh thứ hai (đây là một cấu trúc `if...else` độc lập với khối `if` phía trước):\newline
   `if a <= 30:`\newline
   Điều kiện `25 <= 30` trả về `True`, do đó khối lệnh bên trong `if` được thực thi: `print("Hello")` in ra dòng chữ `Hello`.\newline
   Nhánh `else:` tương ứng bị bỏ qua và không in `Know Program`.\newline
Kết luận: Kết quả hiển thị trên màn hình là `Hello`. Do đó đáp án B là chính xác.
\end{minipage}
\vspace{0.8cm}
\hrule
\vspace{0.5cm}

\begin{minipage}{\textwidth}
\subsection*{Question 44 \quad \footnotesize{[0044_q044_Question_44.png]}}
\noindent \textbf{Which of the following elements of a mathematical expression in Python is evaluated last?}

\begin{enumerate}[label=\Alph*.]
    \item \textbf{Addition} \hfill \textcolor{blue}{\textbf{[Correct]}}
    \item Multiplication
    \item Parentheses
    \item Division
\end{enumerate}

\noindent \textbf{Correct Answer:} \textcolor{blue}{\textbf{A}} \\[4pt]
\noindent \textbf{Explanation:}\\
Theo quy tắc thứ tự ưu tiên của các toán tử trong biểu thức toán học của Python (Quy tắc PEMDAS / BODMAS):\newline
1. **Parentheses** `()` (Dấu ngoặc đơn): Có độ ưu tiên cao nhất, các biểu thức trong ngoặc luôn được tính toán đầu tiên.\newline
2. **Exponentiation** `**` (Toán tử lũy thừa).\newline
3. **Multiplication và Division** (`*`, `/`, `//`, `%`): Nhân và chia có cùng cấp độ ưu tiên, được đánh giá từ trái sang phải.\newline
4. **Addition và Subtraction** (`+`, `-`): Cộng và trừ có độ ưu tiên thấp nhất trong các phép toán số học cơ bản.\newline
\newline
Trong 4 đáp án đưa ra (Addition, Multiplication, Parentheses, Division), phép cộng (**Addition**) có mức độ ưu tiên thấp nhất, do đó nó được đánh giá sau cùng (evaluated last). Đáp án đúng là A.
\end{minipage}
\vspace{0.8cm}
\hrule
\vspace{0.5cm}

\begin{minipage}{\textwidth}
\subsection*{Question 45 \quad \footnotesize{[0045_q045_Question_45.png]}}
\noindent \textbf{Given the file dog_breeds.txt, which of the following is the correct way to open the file for reading as a text file? Select all that apply.}

\begin{enumerate}[label=\Alph*.]
    \item open('dog_breeds.txt', 'wb')
    \item open('dog_breeds.txt', 'rb')
    \item \textbf{open('dog_breeds.txt', 'r')} \hfill \textcolor{blue}{\textbf{[Correct]}}
    \item open('dog_breeds.txt', 'w')
\end{enumerate}

\noindent \textbf{Correct Answer:} \textcolor{blue}{\textbf{C}} \\[4pt]
\noindent \textbf{Explanation:}\\
Hàm `open(file, mode)` trong Python dùng để mở một tệp. Tham số `mode` chỉ định chế độ mở file:\newline
- Chế độ `'r'` (Read text): Mở tệp để đọc dưới dạng văn bản (text mode). Đây cũng là chế độ mặc định của hàm `open()`. Vì vậy `open('dog_breeds.txt', 'r')` là cú pháp chính xác để mở file đọc dưới dạng text.\newline
- Phân tích các chế độ khác:\newline
  + A. `'wb'`: Ghi ở chế độ nhị phân (Write Binary), sẽ xóa sạch nội dung cũ của file hoặc tạo file mới nếu chưa có.\newline
  + B. `'rb'`: Đọc ở chế độ nhị phân (Read Binary), dữ liệu đọc được sẽ là dạng `bytes` thay vì văn bản (`str`).\newline
  + D. `'w'`: Mở file để ghi ở dạng văn bản (Write text), sẽ ghi đè nội dung file.\newline
Do đó, chỉ có đáp án C là mở file để đọc dưới dạng tệp văn bản.
\end{minipage}
\vspace{0.8cm}
\hrule
\vspace{0.5cm}

\begin{minipage}{\textwidth}
\subsection*{Question 46 \quad \footnotesize{[0046_q046_Question_46.png]}}
\noindent \textbf{A bytes object is immutable. If you want a binary sequence that can be modified, you should use a bytearray.\newline

\newline
\newline

\newline
Is it possible to create a bytearray object directly from a bytes object, as shown below:\newline

\newline
\newline

\newline
```python\newline

\newline
array_of_bytes = bytearray(b'15x 80a#')\newline

\newline
```}

\begin{enumerate}[label=\Alph*.]
    \item \textbf{Yes} \hfill \textcolor{blue}{\textbf{[Correct]}}
    \item No
\end{enumerate}

\noindent \textbf{Correct Answer:} \textcolor{blue}{\textbf{A}} \\[4pt]
\noindent \textbf{Explanation:}\\
Trong Python:\newline
- Đối tượng `bytes` là chuỗi byte bất biến (immutable sequence of bytes).\newline
- Đối tượng `bytearray` là chuỗi byte có thể thay đổi được (mutable sequence of bytes).\newline
- Hàm khởi tạo (constructor) `bytearray()` cho phép truyền vào nhiều kiểu dữ liệu nguồn khác nhau, bao gồm: một đối tượng `bytes`, một chuỗi kèm định dạng mã hóa (encoding), một danh sách các số nguyên từ 0 đến 255, hoặc kích thước buffer (int).\newline
- Vì vậy, việc truyền trực tiếp một bytes literal như `b'15\x80a#'` vào `bytearray(b'15\x80a#')` là hoàn toàn hợp lệ trong Python:\newline
  `array_of_bytes = bytearray(b'15\x80a#')`\newline
  sẽ tạo ra một `bytearray` có các phần tử tương ứng và cho phép lập trình viên chỉnh sửa nội dung các byte đó.\newline
Đáp án đúng là **Yes** (A).
\end{minipage}
\vspace{0.8cm}
\hrule
\vspace{0.5cm}

\begin{minipage}{\textwidth}
\subsection*{Question 47 \quad \footnotesize{[0047_q047_Question_47.png]}}
\noindent \textbf{How would you use string slicing [:] to print out 'uct' from the following string?\newline

\newline
\newline

\newline
```python\newline

\newline
x = 'From marquard@uct.ac.za'\newline

\newline
```}

\begin{enumerate}[label=\Alph*.]
    \item print(x[14/17])
    \item \textbf{print(x[14:17])} \hfill \textcolor{blue}{\textbf{[Correct]}}
    \item print(x[14:3])
    \item print(x[14+3])
\end{enumerate}

\noindent \textbf{Correct Answer:} \textcolor{blue}{\textbf{B}} \\[4pt]
\noindent \textbf{Explanation:}\\
Cú pháp cắt chuỗi (string slicing) trong Python là `string[start:end]`, trong đó:\newline
- `start`: Chỉ số bắt đầu (được bao gồm trong kết quả).\newline
- `end`: Chỉ số kết thúc (không được bao gồm, dừng trước vị trí này).\newline
\newline
Xác định vị trí các ký tự trong chuỗi `x = 'From marquard@uct.ac.za'`:\newline
- Chỉ số 0 đến 3: `'From'`\newline
- Chỉ số 4: `' '` (dấu cách)\newline
- Chỉ số 5 đến 12: `'marquard'`\newline
- Chỉ số 13: `'@'`\newline
- Chỉ số 14: `'u'`\newline
- Chỉ số 15: `'c'`\newline
- Chỉ số 16: `'t'`\newline
- Chỉ số 17: `'.'`\newline
\newline
Để trích xuất chuỗi con `'uct'` gồm 3 ký tự tại vị trí index 14, 15, 16:\newline
- Vị trí bắt đầu: `start = 14`\newline
- Vị trí kết thúc: `end = 17` (vì `end` không lấy, sẽ lấy từ 14 đến 16)\newline
Cú pháp chính xác là `x[14:17]`. Do đó lệnh `print(x[14:17])` in ra `'uct'`.\newline
Đáp án đúng là B.
\end{minipage}
\vspace{0.8cm}
\hrule
\vspace{0.5cm}

\begin{minipage}{\textwidth}
\subsection*{Question 48 \quad \footnotesize{[0048_q048_Question_48.png]}}
\noindent \textbf{```python\newline

\newline
class Sales:\newline

\newline
    def __init__(self, id):\newline

\newline
        self.id = id\newline

\newline
        id = 100\newline

\newline
\newline

\newline
val = Sales(123)\newline

\newline
print (val.id)\newline

\newline
```}

\begin{enumerate}[label=\Alph*.]
    \item SyntaxError, this program will not run
    \item 100
    \item \textbf{123} \hfill \textcolor{blue}{\textbf{[Correct]}}
    \item None of the above
\end{enumerate}

\noindent \textbf{Correct Answer:} \textcolor{blue}{\textbf{C}} \\[4pt]
\noindent \textbf{Explanation:}\\
Phân tích từng dòng thực thi mã nguồn:\newline
1. Lớp `Sales` có phương thức khởi tạo `__init__(self, id)` nhận tham số `self` (đối tượng hiện tại) và `id` (tham số truyền vào).\newline
2. Dòng `self.id = id`: Gán giá trị của tham số `id` cho thuộc tính của đối tượng `self.id`.\newline
3. Dòng `id = 100`: Chỉ gán lại giá trị cho biến cục bộ (local variable) `id` trong phạm vi của phương thức `__init__`. Biến này sẽ bị hủy khi phương thức kết thúc, và phép gán này hoàn toàn không làm thay đổi thuộc tính `self.id` của đối tượng.\newline
4. Lệnh `val = Sales(123)` tạo một thực thể mới của lớp `Sales` với đối số `123`. Trong `__init__`, `self.id` được gán bằng `123`.\newline
5. Khi gọi `print(val.id)`, chương trình truy cập vào thuộc tính của thực thể `val`, giá trị này vẫn là `123`.\newline
Do đó kết quả in ra màn hình là `123`. Đáp án C là chính xác.
\end{minipage}
\vspace{0.8cm}
\hrule
\vspace{0.5cm}

\begin{minipage}{\textwidth}
\subsection*{Question 49 \quad \footnotesize{[0049_q049_Question_49.png]}}
\noindent \textbf{The _________ keyword defines a template indicating the data that will be in an object of the class and the functions that can be called on an object of the class.}

\begin{enumerate}[label=\Alph*.]
    \item Class
    \item object
    \item \textbf{class} \hfill \textcolor{blue}{\textbf{[Correct]}}
    \item instance
\end{enumerate}

\noindent \textbf{Correct Answer:} \textcolor{blue}{\textbf{C}} \\[4pt]
\noindent \textbf{Explanation:}\\
Trong lập trình hướng đối tượng (OOP) với Python:\newline
- Một **lớp** (class) đóng vai trò là bản thiết kế mẫu (template / blueprint) định nghĩa các thuộc tính dữ liệu (data attributes) và các hàm / phương thức (methods) mà các đối tượng (objects) thuộc lớp đó sở hữu.\newline
- Từ khóa trong Python dùng để định nghĩa một lớp là `class` (viết thường toàn bộ).\newline
- Ví dụ:\newline
  ```python\newline
  class Dog:\newline
      def __init__(self, name):\newline
          self.name = name\newline
  ```\newline
- Lưu ý: Python phân biệt chữ hoa chữ thường (case-sensitive), do đó `Class` (viết hoa chữ C ở đáp án A) là sai cú pháp. Các từ `object` và `instance` là thuật ngữ chỉ đối tượng hoặc thực thể, không phải từ khóa để định nghĩa template lớp.\newline
Đáp án đúng là C (`class`).
\end{minipage}
\vspace{0.8cm}
\hrule
\vspace{0.5cm}

\begin{minipage}{\textwidth}
\subsection*{Question 50 \quad \footnotesize{[0050_q050_Question_50.png]}}
\noindent \textbf{Which of the following commands can be used to read the next line in a file using the file object ?}

\begin{enumerate}[label=\Alph*.]
    \item tmpfile.read(n)
    \item tmpfile.read()
    \item \textbf{tmpfile.readline()} \hfill \textcolor{blue}{\textbf{[Correct]}}
    \item tmpfile.readlines()
\end{enumerate}

\noindent \textbf{Correct Answer:} \textcolor{blue}{\textbf{C}} \\[4pt]
\noindent \textbf{Explanation:}\\
Trong Python, khi làm việc với đối tượng tệp (file object):\newline
- Phương thức `readline()` đọc duy nhất **một dòng tiếp theo** từ vị trí con trỏ hiện tại trong tệp cho đến khi gặp ký tự xuống dòng `\n` (hoặc hết tệp), và trả về một chuỗi (`str`).\newline
- Phân tích các phương thức khác của đối tượng file:\newline
  + A. `read(n)`: Đọc tối đa `n` ký tự (hoặc bytes) từ file, không giới hạn theo từng dòng.\newline
  + B. `read()`: Đọc toàn bộ nội dung còn lại của file thành một chuỗi duy nhất.\newline
  + C. `readline()`: Đọc chính xác một dòng tiếp theo của file.\newline
  + D. `readlines()`: Đọc toàn bộ các dòng còn lại trong file và trả về một danh sách các chuỗi (list of strings).\newline
Do đó, lệnh dùng để đọc dòng tiếp theo là `tmpfile.readline()`. Đáp án đúng là C.
\end{minipage}
\vspace{0.8cm}
\hrule
\vspace{0.5cm}

\begin{minipage}{\textwidth}
\subsection*{Question 51 \quad \footnotesize{[0051_q051_Question_51.png]}}
\noindent \textbf{What is "self" typically used for in a Python method within a class?}

\begin{enumerate}[label=\Alph*.]
    \item To terminate a loop
    \item To set the residual value in an expression where the method is used
    \item The number of parameters to the method
    \item \textbf{To refer to the instance in which the method is being called} \hfill \textcolor{blue}{\textbf{[Correct]}}
\end{enumerate}

\noindent \textbf{Correct Answer:} \textcolor{blue}{\textbf{D}} \\[4pt]
\noindent \textbf{Explanation:}\\
Trong Python:\newline
- `self` là tham số đầu tiên theo quy ước trong các phương thức của một lớp (instance methods).\newline
- `self` đại diện cho chính đối tượng / thực thể cụ thể (instance) mà phương thức đó đang được gọi.\newline
- Nhờ có `self`, phương thức có thể truy cập, đọc và ghi các thuộc tính dữ liệu (`self.attribute`) cũng như gọi các phương thức khác (`self.other_method()`) của chính thực thể đó.\newline
- Khi gọi phương thức qua đối tượng (ví dụ: `my_obj.method()`), Python tự động truyền đối tượng `my_obj` vào vị trí tham số `self`.\newline
- Phân tích các phương án khác:\newline
  + A sai: Lệnh kết thúc vòng lặp là `break`.\newline
  + B, C sai: Không liên quan đến ý nghĩa và chức năng của `self`.\newline
Do đó, `self` dùng để tham chiếu đến thực thể mà phương thức đang được gọi trên đó (To refer to the instance in which the method is being called). Đáp án đúng là D.
\end{minipage}
\vspace{0.8cm}
\hrule
\vspace{0.5cm}

\begin{minipage}{\textwidth}
\subsection*{Question 52 \quad \footnotesize{[0052_q052_Question_52.png]}}
\noindent \textbf{What would the following Python code print out?\newline

\newline
\newline

\newline
```python\newline

\newline
stuff = dict()\newline

\newline
print(stuff.get('candy',-1))\newline

\newline
```}

\begin{enumerate}[label=\Alph*.]
    \item \textbf{-1} \hfill \textcolor{blue}{\textbf{[Correct]}}
    \item The program would fail with a traceback
    \item 0
    \item 'candy'
\end{enumerate}

\noindent \textbf{Correct Answer:} \textcolor{blue}{\textbf{A}} \\[4pt]
\noindent \textbf{Explanation:}\\
Chi tiết từng bước thực thi:\newline
1. `stuff = dict()`: Tạo một từ điển rỗng (empty dictionary) gán cho biến `stuff`.\newline
2. Phương thức `dict.get(key, default)`:\newline
   - Dùng để truy xuất giá trị tương ứng với khóa `key` trong từ điển mà không làm phát sinh lỗi `KeyError` nếu khóa không tồn tại.\newline
   - Nếu `key` có trong từ điển: trả về giá trị của `key`.\newline
   - Nếu `key` không có trong từ điển: trả về giá trị mặc định được chỉ định ở tham số thứ hai `default`.\newline
3. Trong biểu thức `stuff.get('candy', -1)`:\newline
   - Khóa `'candy'` không tồn tại trong từ điển `stuff`.\newline
   - Giá trị mặc định được chỉ định là `-1`.\newline
   - Do đó, phương thức trả về `-1`.\newline
4. Hàm `print()` in kết quả `-1` ra màn hình.\newline
Đáp án đúng là A.
\end{minipage}
\vspace{0.8cm}
\hrule
\vspace{0.5cm}

\begin{minipage}{\textwidth}
\subsection*{Question 53 \quad \footnotesize{[0053_q053_Question_53.png]}}
\noindent \textbf{Which of the following reserved keyword is used to end a function, return generator?}

\begin{enumerate}[label=\Alph*.]
    \item break
    \item switch
    \item Return
    \item \textbf{yield} \hfill \textcolor{blue}{\textbf{[Correct]}}
\end{enumerate}

\noindent \textbf{Correct Answer:} \textcolor{blue}{\textbf{D}} \\[4pt]
\noindent \textbf{Explanation:}\\
Trong Python:\newline
- Từ khóa `yield` được sử dụng bên trong một hàm để biến hàm đó thành một **hàm tạo (generator function)**.\newline
- Khi hàm generator được gọi, nó không thực thi toàn bộ phần thân hàm ngay lập tức mà trả về một đối tượng **generator**.\newline
- Mỗi khi gặp câu lệnh `yield`, hàm sẽ tạm dừng thực thi tại vị trí đó, lưu lại trạng thái cục bộ của nó và trả về giá trị biểu thức đi kèm `yield` cho bên ngoài (thông qua hàm `next()`). Khi được gọi tiếp tục, hàm sẽ chạy tiếp từ dòng lệnh ngay sau `yield`.\newline
- Phân tích các đáp án khác:\newline
  + A. `break`: Dùng để thoát sớm khỏi vòng lặp `for` hoặc `while`.\newline
  + B. `switch`: Không phải từ khóa trong Python (Python dùng `match...case` từ phiên bản 3.10).\newline
  + C. `Return`: Viết hoa chữ cái đầu là sai cú pháp trong Python (`return` viết thường), đồng thời `return` thông thường trả về một giá trị duy nhất và kết thúc hàm hoàn toàn, không tạo generator.\newline
Đáp án đúng là D (`yield`).
\end{minipage}
\vspace{0.8cm}
\hrule
\vspace{0.5cm}

\begin{minipage}{\textwidth}
\subsection*{Question 54 \quad \footnotesize{[0054_q054_Question_54.png]}}
\noindent \textbf{What will be the value of x when the following statement is executed:\newline

\newline
\newline

\newline
```python\newline

\newline
x = round (98.6)\newline

\newline
```}

\begin{enumerate}[label=\Alph*.]
    \item 100
    \item 6
    \item 98
    \item \textbf{99} \hfill \textcolor{blue}{\textbf{[Correct]}}
\end{enumerate}

\noindent \textbf{Correct Answer:} \textcolor{blue}{\textbf{D}} \\[4pt]
\noindent \textbf{Explanation:}\\
Hàm tích hợp sẵn `round(number[, ndigits])` trong Python:\newline
- Dùng để làm tròn một số thực đến một độ chính xác nhất định.\newline
- Khi không truyền đối số thứ hai (`ndigits`), hàm `round()` sẽ làm tròn số thực về số nguyên gần nhất (integer).\newline
- Đối với số `98.6`, phần thập phân là `0.6` (lớn hơn 0.5), do đó số nguyên gần nhất là `99`.\newline
- Vì vậy, `round(98.6)` trả về giá trị `99`, và biến `x` nhận giá trị `99`.\newline
Đáp án đúng là D.
\end{minipage}
\vspace{0.8cm}
\hrule
\vspace{0.5cm}

\begin{minipage}{\textwidth}
\subsection*{Question 55 \quad \footnotesize{[0055_q055_Question_55.png]}}
\noindent \textbf{What does the Python dir() function show when we pass an object into it as a parameter?}

\begin{enumerate}[label=\Alph*.]
    \item It shows the parent class
    \item It shows the number of parameters to the constructor
    \item \textbf{It shows the methods and attributes of the object} \hfill \textcolor{blue}{\textbf{[Correct]}}
    \item It shows the type of the object
\end{enumerate}

\noindent \textbf{Correct Answer:} \textcolor{blue}{\textbf{C}} \\[4pt]
\noindent \textbf{Explanation:}\\
Theo tài liệu chuẩn của Python (Python Documentation):\newline
- Hàm `dir([object])` khi được truyền vào một đối tượng sẽ cố gắng trả về một danh sách (`list`) chứa tên các thuộc tính và phương thức hợp lệ (valid attributes and methods) của đối tượng đó.\newline
- Danh sách này bao gồm cả các thuộc tính đặc biệt (magic/dunder methods như `__init__`, `__str__`, v.v.) và các thuộc tính, phương thức thông thường do người dùng tự định nghĩa hoặc kế thừa từ lớp cha.\newline
- Phân tích các đáp án:\newline
  + A: Muốn xem lớp cha thường dùng `cls.__bases__`.\newline
  + B: `dir()` không hiển thị số lượng tham số của hàm khởi tạo (constructor).\newline
  + C: Hiển thị các phương thức và thuộc tính của đối tượng (It shows the methods and attributes of the object) - Đây là câu trả lời chính xác.\newline
  + D: Muốn xem kiểu của đối tượng ta dùng hàm `type(object)`.\newline
Đáp án đúng là C.
\end{minipage}
\vspace{0.8cm}
\hrule
\vspace{0.5cm}

\begin{minipage}{\textwidth}
\subsection*{Question 56 \quad \footnotesize{[0056_q056_Question_56.png]}}
\noindent \textbf{When Python is running in the interactive mode and displaying the chevron prompt (>>>) - what question is Python asking you?}

\begin{enumerate}[label=\Alph*.]
    \item What Python script would you like me to run?
    \item What is the next machine language instruction to run?
    \item What is your favorite color?
    \item \textbf{What Python statement would you like me to run?} \hfill \textcolor{blue}{\textbf{[Correct]}}
\end{enumerate}

\noindent \textbf{Correct Answer:} \textcolor{blue}{\textbf{D}} \\[4pt]
\noindent \textbf{Explanation:}\\
Theo giáo trình "Python for Everybody" của GS. Charles Severance (giáo trình chuẩn cho học phần PFP191):\newline
- Dấu nhắc lệnh `>>>` (được gọi là chevron prompt) xuất hiện khi bạn chạy trình thông dịch Python ở chế độ tương tác (interactive mode / Python REPL).\newline
- Dấu nhắc này ngụ ý câu hỏi mà Python đang gửi tới bạn: **"What Python statement would you like me to run?"** (Bạn muốn tôi chạy câu lệnh Python nào tiếp theo?).\newline
- Bạn nhập vào một câu lệnh Python hợp lệ (ví dụ: `print('Hello')`), nhấn Enter, Python sẽ thực thi ngay câu lệnh đó và hiển thị kết quả trực tiếp, sau đó lại đưa ra dấu nhắc `>>>` để chờ câu lệnh kế tiếp.\newline
Đáp án đúng là D.
\end{minipage}
\vspace{0.8cm}
\hrule
\vspace{0.5cm}

\begin{minipage}{\textwidth}
\subsection*{Question 57 \quad \footnotesize{[0057_q057_Question_57.png]}}
\noindent \textbf{When you have multiple lines in an if block, how do you indicate the end of the if block?}

\begin{enumerate}[label=\Alph*.]
    \item You omit the semicolon ; on the last line of the if block
    \item \textbf{You de-indent the next line past the if block to the same level of indent as the original if statement} \hfill \textcolor{blue}{\textbf{[Correct]}}
    \item You put the colon : character on a line by itself to indicate we are done with the if block
    \item You use a curly brace  after the last line of the if block
\end{enumerate}

\noindent \textbf{Correct Answer:} \textcolor{blue}{\textbf{B}} \\[4pt]
\noindent \textbf{Explanation:}\\
Trong Python, cấu trúc phân cấp và các khối lệnh (blocks of code) được định nghĩa hoàn toàn dựa trên quy tắc **thụt lề (indentation)** thay vì sử dụng các cặp ngoặc nhọn `{}` như C/C++/Java hay các từ khóa như `endif` trong Pascal/PHP.\newline
- Tất cả các câu lệnh nằm trong cùng một khối `if` phải có cùng độ thụt lề (thông thường là 4 dấu cách).\newline
- Để kết thúc khối lệnh `if`, bạn chỉ cần **giảm mức thụt lề (de-indent / unindent)** ở dòng lệnh kế tiếp về bằng với mức thụt lề của câu lệnh `if` ban đầu. Dòng lệnh có mức thụt lề bằng với `if` sẽ nằm ngoài khối `if` và được thực thi sau khi cấu trúc `if` kết thúc.\newline
- Phân tích các đáp án sai:\newline
  + A sai: Python không dùng dấu chấm phẩy `;` để báo hiệu kết thúc khối lệnh.\newline
  + C sai: Dấu hai chấm `:` chỉ dùng ở cuối câu lệnh rẽ nhánh `if <condition>:`, không bao giờ đứng một mình trên một dòng để kết thúc block.\newline
  + D sai: Python không dùng dấu ngoặc nhọn `{}` để đóng khối lệnh.\newline
Đáp án đúng là B.
\end{minipage}
\vspace{0.8cm}
\hrule
\vspace{0.5cm}

\begin{minipage}{\textwidth}
\subsection*{Question 58 \quad \footnotesize{[0058_q058_Question_58.png]}}
\noindent \textbf{Identify the data structure we use in the code given below?\newline

\newline
# Code starts\newline

\newline
\newline

\newline
asthetic1 = 'quizOrbit'\newline

\newline
asthetic2 = 'QuizOrbit'\newline

\newline
\newline

\newline
print(asthetic1!=asthetic2)\newline

\newline
\newline

\newline
# Code Ends}

\begin{enumerate}[label=\Alph*.]
    \item HashMap
    \item Tree
    \item Dictionary
    \item \textbf{True} \hfill \textcolor{blue}{\textbf{[Correct]}}
\end{enumerate}

\noindent \textbf{Correct Answer:} \textcolor{blue}{\textbf{D}} \\[4pt]
\noindent \textbf{Explanation:}\\
Phân tích đoạn mã và ngữ cảnh câu hỏi thi PFP191:\newline
1. Đoạn mã khởi tạo hai chuỗi ký tự:\newline
   - `asthetic1 = 'quizOrbit'`\newline
   - `asthetic2 = 'QuizOrbit'`\newline
2. Lệnh `print(asthetic1!=asthetic2)`:\newline
   - Thực hiện phép toán so sánh khác (`!=`) giữa hai chuỗi ký tự.\newline
   - Vì Python phân biệt chữ hoa và chữ thường (case-sensitive), chuỗi `'quizOrbit'` có ký tự đầu là chữ thường `'q'` (mã ASCII là 113) khác với ký tự `'Q'` hoa (mã ASCII là 81) của chuỗi `'QuizOrbit'`.\newline
   - Do hai chuỗi này khác nhau, phép so sánh `asthetic1 != asthetic2` trả về giá trị kiểu boolean là `True`.\newline
   - Lệnh `print` xuất ra màn hình giá trị `True`.\newline
3. Ngữ cảnh đề thi ngân hàng (FUOverflow / EOS): Tiêu đề câu hỏi ghi "Identify the data structure..." là lỗi sao chép đề từ hệ thống đề thi, tuy nhiên trong các phương án chọn lựa:\newline
   - A. HashMap\newline
   - B. Tree\newline
   - C. Dictionary\newline
   - D. True\newline
   Phương án D chính là kết quả đầu ra (output) của biểu thức `print(asthetic1!=asthetic2)`.\newline
Do đó, đáp án đúng theo hệ thống chấm thi là D (`True`).
\end{minipage}
\vspace{0.8cm}
\hrule
\vspace{0.5cm}

\begin{minipage}{\textwidth}
\subsection*{Question 59 \quad \footnotesize{[0059_q059_Question_59.png]}}
\noindent \textbf{What keyword is used to indicate the start of a method in a Python class?}

\begin{enumerate}[label=\Alph*.]
    \item \textbf{def} \hfill \textcolor{blue}{\textbf{[Correct]}}
    \item break
    \item function
    \item continue
\end{enumerate}

\noindent \textbf{Correct Answer:} \textcolor{blue}{\textbf{A}} \\[4pt]
\noindent \textbf{Explanation:}\\
Trong Python:\newline
- Một phương thức (method) trong lớp thực chất là một hàm (function) được định nghĩa bên trong khối lệnh của lớp đó.\newline
- Cả hàm thông thường và phương thức của lớp đều được bắt đầu khai báo bằng từ khóa `def` (viết tắt của từ "define").\newline
- Ví dụ về khai báo phương thức trong lớp:\newline
  ```python\newline
  class PartyAnimal:\newline
      def party(self):\newline
          print('Having a party!')\newline
  ```\newline
- Phân tích các phương án khác:\newline
  + B. `break`: Dùng để kết thúc sớm một vòng lặp (`for`/`while`).\newline
  + C. `function`: Không phải từ khóa trong Python (đây là từ khóa khai báo hàm trong JavaScript).\newline
  + D. `continue`: Dùng để bỏ qua lần lặp hiện tại và chuyển sang lần lặp kế tiếp của vòng lặp.\newline
Đáp án đúng là A (`def`).
\end{minipage}
\vspace{0.8cm}
\hrule
\vspace{0.5cm}

\begin{minipage}{\textwidth}
\subsection*{Question 60 \quad \footnotesize{[0060_q060_Question_60.png]}}
\noindent \textbf{What is the value of the following expression? 42 % 10}

\begin{enumerate}[label=\Alph*.]
    \item \textbf{2} \hfill \textcolor{blue}{\textbf{[Correct]}}
    \item 420
    \item 4210
    \item 42
\end{enumerate}

\noindent \textbf{Correct Answer:} \textcolor{blue}{\textbf{A}} \\[4pt]
\noindent \textbf{Explanation:}\\
Trong Python, toán tử `%` là toán tử chia lấy phần dư (modulus operator).\newline
- Biểu thức `42 % 10` thực hiện phép chia số nguyên 42 cho 10:\newline
  + Phần nguyên (thương): `42 // 10 = 4`\newline
  + Phép nhân ngược lại: `4 * 10 = 40`\newline
  + Phần dư còn lại: `42 - 40 = 2`\newline
- Do đó, giá trị của biểu thức `42 % 10` là `2`.\newline
Đáp án đúng là A.
\end{minipage}
\vspace{0.8cm}
\hrule
\vspace{0.5cm}

\begin{minipage}{\textwidth}
\subsection*{Question 61 \quad \footnotesize{[0061_q061_Question_61.png]}}
\noindent \textbf{What is the iteration variable in the following Python code?\newline

\newline
\newline

\newline
```python\newline

\newline
for letter in 'banana' :\newline

\newline
    print(letter)\newline

\newline
```}

\begin{enumerate}[label=\Alph*.]
    \item 'banana'
    \item for
    \item \textbf{letter} \hfill \textcolor{blue}{\textbf{[Correct]}}
    \item print
\end{enumerate}

\noindent \textbf{Correct Answer:} \textcolor{blue}{\textbf{C}} \\[4pt]
\noindent \textbf{Explanation:}\\
Trong cấu trúc vòng lặp `for <variable> in <iterable>:`, biến đứng ngay sau từ khóa `for` được gọi là biến lặp (iteration variable). Ở mỗi lần lặp, biến này sẽ nhận một giá trị (từng ký tự) kế tiếp từ đối tượng có thể lặp (iterable, ở đây là chuỗi `'banana'`). Do đó, `letter` chính là biến lặp.\newline
\newline
- `'banana'` là chuỗi ký tự được duyệt qua (iterable sequence).\newline
- `for` là từ khóa cú pháp bắt đầu vòng lặp.\newline
- `print` là hàm tích hợp sẵn (built-in function) dùng để xuất dữ liệu.
\end{minipage}
\vspace{0.8cm}
\hrule
\vspace{0.5cm}

\begin{minipage}{\textwidth}
\subsection*{Question 62 \quad \footnotesize{[0062_q062_Question_62.png]}}
\noindent \textbf{What do we use the second parameter of the open() call to indicate}

\begin{enumerate}[label=\Alph*.]
    \item How large we expect the file to be
    \item The list of folders to be searched to find the file we want to open
    \item \textbf{Whether we want to read data from the file or write data to the file} \hfill \textcolor{blue}{\textbf{[Correct]}}
    \item What disk drive the file is stored on
\end{enumerate}

\noindent \textbf{Correct Answer:} \textcolor{blue}{\textbf{C}} \\[4pt]
\noindent \textbf{Explanation:}\\
Hàm `open(file, mode='r')` trong Python nhận tham số thứ nhất là tên/đường dẫn tệp tin và tham số thứ hai là `mode` (chế độ mở file). Chế độ này xác định hành động thao tác với tệp: đọc (`'r'`), ghi đè (`'w'`), ghi nối tiếp (`'a'`), hoặc nhị phân (`'b'`), v.v. Do đó, tham số thứ hai dùng để chỉ định xem chúng ta muốn đọc dữ liệu từ tệp tin hay ghi dữ liệu vào tệp tin ("Whether we want to read data from the file or write data to the file").
\end{minipage}
\vspace{0.8cm}
\hrule
\vspace{0.5cm}

\begin{minipage}{\textwidth}
\subsection*{Question 63 \quad \footnotesize{[0063_q063_Question_63.png]}}
\noindent \textbf{A list of lines is returned if using the method of .....}

\begin{enumerate}[label=\Alph*.]
    \item <readable()>
    \item <readnines()>
    \item \textbf{<readlines()>} \hfill \textcolor{blue}{\textbf{[Correct]}}
    \item <readline()>
\end{enumerate}

\noindent \textbf{Correct Answer:} \textcolor{blue}{\textbf{C}} \\[4pt]
\noindent \textbf{Explanation:}\\
Phương thức `readlines()` của đối tượng tệp tin trong Python đọc toàn bộ nội dung của tệp và trả về một danh sách (`list`) các chuỗi ký tự, trong đó mỗi phần tử là một dòng văn bản trong tệp tin (bao gồm cả ký tự xuống dòng `\n`).\newline
\newline
- `readline()` chỉ đọc duy nhất một dòng kế tiếp dưới dạng chuỗi đơn (`str`).\newline
- `readable()` trả về giá trị kiểu boolean (`True`/`False`) kiểm tra xem file có thể đọc được hay không.\newline
- `<readnines()>` không tồn tại trong Python.
\end{minipage}
\vspace{0.8cm}
\hrule
\vspace{0.5cm}

\begin{minipage}{\textwidth}
\subsection*{Question 64 \quad \footnotesize{[0064_q064_Question_64.png]}}
\noindent \textbf{Which of the following is not a Python reserved word?}

\begin{enumerate}[label=\Alph*.]
    \item \textbf{iterate} \hfill \textcolor{blue}{\textbf{[Correct]}}
    \item continue
    \item else
    \item break
\end{enumerate}

\noindent \textbf{Correct Answer:} \textcolor{blue}{\textbf{A}} \\[4pt]
\noindent \textbf{Explanation:}\\
Từ khóa dành riêng (reserved words / keywords) trong Python là các từ có ý nghĩa ngữ pháp cố định được định nghĩa trước bởi ngôn ngữ, không được sử dụng để đặt tên biến, hàm hay định danh. Có thể kiểm tra bằng danh sách `keyword.kwlist`.\newline
\newline
- `continue`, `else`, `break` đều là các từ khóa dành riêng trong Python.\newline
- `iterate` không phải là từ khóa của Python (lập trình viên hoàn toàn có thể tự do đặt tên biến hoặc hàm là `iterate`).
\end{minipage}
\vspace{0.8cm}
\hrule
\vspace{0.5cm}

\begin{minipage}{\textwidth}
\subsection*{Question 65 \quad \footnotesize{[0065_q065_Question_65.png]}}
\noindent \textbf{In which of the following does the CricketFan class correctly inherit from the PartyAnimal class?}

\begin{enumerate}[label=\Alph*.]
    \item from party import PartyAnimal
    \item \textbf{class CricketFan(PartyAnimal)} \hfill \textcolor{blue}{\textbf{[Correct]}}
    \item an = PartyAnimal()
    \item CricketFan = PartyAnimal()
\end{enumerate}

\noindent \textbf{Correct Answer:} \textcolor{blue}{\textbf{B}} \\[4pt]
\noindent \textbf{Explanation:}\\
Trong lập trình hướng đối tượng (OOP) với Python, cú pháp để định nghĩa một lớp con kế thừa từ một lớp cha là đặt tên lớp cha trong cặp dấu ngoặc đơn ngay sau tên lớp con:\newline
`class SubClassName(BaseClassName):`\newline
Do đó, cú pháp để lớp `CricketFan` kế thừa từ lớp `PartyAnimal` là `class CricketFan(PartyAnimal)`.\newline
\newline
- Phương án A là câu lệnh nhập một lớp từ module.\newline
- Phương án C và D là các câu lệnh khởi tạo đối tượng hoặc gán biến.
\end{minipage}
\vspace{0.8cm}
\hrule
\vspace{0.5cm}

\begin{minipage}{\textwidth}
\subsection*{Question 66 \quad \footnotesize{[0066_q066_Question_66.png]}}
\noindent \textbf{In the following code (numbers added) - which will be the last line to execute successfully?\newline

\newline
\newline

\newline
```python\newline

\newline
(1)  astr = 'Hello Bob'\newline

\newline
(2)  istr = int(astr)\newline

\newline
(3)  print('First', istr)\newline

\newline
(4)  astr = '123'\newline

\newline
(5)  istr = int(astr)\newline

\newline
(6)  print('Second', istr)\newline

\newline
```}

\begin{enumerate}[label=\Alph*.]
    \item 2
    \item \textbf{1} \hfill \textcolor{blue}{\textbf{[Correct]}}
    \item 6
    \item 3
\end{enumerate}

\noindent \textbf{Correct Answer:} \textcolor{blue}{\textbf{B}} \\[4pt]
\noindent \textbf{Explanation:}\\
Phân tích quá trình thực thi từng dòng lệnh:\newline
- Dòng (1): Gán chuỗi `'Hello Bob'` cho biến `astr`. Dòng này thực thi thành công.\newline
- Dòng (2): `istr = int(astr)` cố gắng chuyển chuỗi `'Hello Bob'` thành số nguyên. Do chuỗi chứa các ký tự chữ cái không thể chuyển thành số nguyên cơ số 10, Python sẽ ném ngoại lệ `ValueError: invalid literal for int() with base 10: 'Hello Bob'`.\newline
Do không có khối `try...except` để xử lý ngoại lệ này, chương trình lập tức dừng lại (crash) tại dòng (2) mà không thực thi thành công dòng (2) hay các dòng sau đó (3, 4, 5, 6).\newline
Vì vậy, dòng lệnh cuối cùng thực thi thành công là dòng (1).
\end{minipage}
\vspace{0.8cm}
\hrule
\vspace{0.5cm}

\begin{minipage}{\textwidth}
\subsection*{Question 67 \quad \footnotesize{[0067_q067_Question_67.png]}}
\noindent \textbf{For the following code:\newline

\newline
\newline

\newline
```python\newline

\newline
astr = 'Hello Bob'\newline

\newline
istr = 0\newline

\newline
try:\newline

\newline
    istr = int(astr)\newline

\newline
except:\newline

\newline
    istr = -1\newline

\newline
```\newline

\newline
\newline

\newline
What will the value be for istr variable after this code executes?}

\begin{enumerate}[label=\Alph*.]
    \item \textbf{-1} \hfill \textcolor{blue}{\textbf{[Correct]}}
    \item It will be the 'Not a number' value (i.e. NaN)
    \item It will be a random number depending on the operating system the program runs on
    \item false
\end{enumerate}

\noindent \textbf{Correct Answer:} \textcolor{blue}{\textbf{A}} \\[4pt]
\noindent \textbf{Explanation:}\\
Từng bước thực thi của đoạn mã:\newline
1. Khởi tạo biến: `astr = 'Hello Bob'` và `istr = 0`.\newline
2. Bắt đầu khối `try:`: câu lệnh `istr = int(astr)` được thực thi. Hàm `int('Hello Bob')` gây ra lỗi `ValueError` vì chuỗi ký tự không thể ép kiểu sang số nguyên.\newline
3. Khi xảy ra ngoại lệ, phép gán cho `istr` tại dòng này bị gián đoạn và quyền điều khiển nhảy ngay sang khối `except:`.\newline
4. Trong khối `except:`, câu lệnh `istr = -1` được thực thi và gán giá trị `-1` cho `istr`.\newline
Do đó, sau khi đoạn mã hoàn tất, giá trị của `istr` là `-1`.
\end{minipage}
\vspace{0.8cm}
\hrule
\vspace{0.5cm}

\begin{minipage}{\textwidth}
\subsection*{Question 68 \quad \footnotesize{[0068_q068_Question_68.png]}}
\noindent \textbf{What is the result of the following statement:\newline

\newline
\newline

\newline
```python\newline

\newline
list((b'abcde' + 'fghi')[3:6])\newline

\newline
```}

\begin{enumerate}[label=\Alph*.]
    \item [b'd', b'e', b'f']
    \item b'def'
    \item [100, 101, 102]
    \item \textbf{It raises an exception} \hfill \textcolor{blue}{\textbf{[Correct]}}
\end{enumerate}

\noindent \textbf{Correct Answer:} \textcolor{blue}{\textbf{D}} \\[4pt]
\noindent \textbf{Explanation:}\\
Trong Python 3, kiểu dữ liệu `bytes` (`b'abcde'`) và kiểu dữ liệu chuỗi `str` (`'fghi'`) là hai kiểu dữ liệu hoàn toàn khác biệt.\newline
Toán tử cộng chuỗi (`+`) không hỗ trợ việc nối giữa một đối tượng `bytes` và một đối tượng `str`. Biểu thức `b'abcde' + 'fghi'` sẽ trực tiếp ném ra ngoại lệ `TypeError: can't concat str to bytes` (hoặc `TypeError: can't concat bytes to str`).\newline
Vì vậy, câu lệnh sẽ phát sinh một ngoại lệ ("It raises an exception").
\end{minipage}
\vspace{0.8cm}
\hrule
\vspace{0.5cm}

\begin{minipage}{\textwidth}
\subsection*{Question 69 \quad \footnotesize{[0069_q069_Question_69.png]}}
\noindent \textbf{What are the names of the two types of mode that are used in Python?}

\begin{enumerate}[label=\Alph*.]
    \item \textbf{Interactive mode & Script mode} \hfill \textcolor{blue}{\textbf{[Correct]}}
    \item Interactive mode & Scratch mode
    \item Interactive mode & Imperial mode
    \item Imperial mode & Script mode
\end{enumerate}

\noindent \textbf{Correct Answer:} \textcolor{blue}{\textbf{A}} \\[4pt]
\noindent \textbf{Explanation:}\\
Python có hai chế độ hoạt động cơ bản để thực thi mã lệnh:\newline
1. **Interactive mode (Chế độ tương tác)**: Người dùng nhập trực tiếp từng câu lệnh vào trình thông dịch Python (qua giao diện dòng lệnh với dấu nhắc `>>>`) và nhận kết quả phản hồi ngay lập tức sau mỗi dòng lệnh.\newline
2. **Script mode (Chế độ kịch bản)**: Lập trình viên viết toàn bộ chương trình và lưu thành một tệp mã nguồn (file có đuôi mở rộng `.py`), sau đó chuyển tệp này cho trình thông dịch Python thực thi từ đầu đến cuối.\newline
Vì vậy, hai chế độ là Interactive mode & Script mode.
\end{minipage}
\vspace{0.8cm}
\hrule
\vspace{0.5cm}

\begin{minipage}{\textwidth}
\subsection*{Question 70 \quad \footnotesize{[0070_q070_Question_70.png]}}
\noindent \textbf{Assume the variable x has been initialized to an integer value (e.g., x = 3). What does the following statement do?\newline

\newline
\newline

\newline
```python\newline

\newline
x = x + 2\newline

\newline
```}

\begin{enumerate}[label=\Alph*.]
    \item This would fail as it is a syntax error
    \item Produce the value "false" because "x" can never equal "x+2"
    \item Create a function called "x" and put the value 2 in the function
    \item \textbf{Retrieve the current value for x, add two to it and put the sum back into x} \hfill \textcolor{blue}{\textbf{[Correct]}}
\end{enumerate}

\noindent \textbf{Correct Answer:} \textcolor{blue}{\textbf{D}} \\[4pt]
\noindent \textbf{Explanation:}\\
Trong Python, dấu `=` là toán tử gán (assignment operator), không phải dấu bằng so sánh trong toán học (`==`).\newline
Khi câu lệnh `x = x + 2` được thực thi:\newline
1. Biểu thức ở vế phải (`x + 2`) được đánh giá trước: Python truy xuất giá trị hiện tại của biến `x`, sau đó cộng thêm 2 vào giá trị đó.\newline
2. Kết quả tổng tính được sau đó sẽ được gán ngược lại (lưu vào) biến `x` ở vế trái.\newline
Do đó, phát biểu đúng là: "Retrieve the current value for x, add two to it and put the sum back into x".
\end{minipage}
\vspace{0.8cm}
\hrule
\vspace{0.5cm}

\begin{minipage}{\textwidth}
\subsection*{Question 71 \quad \footnotesize{[0071_q071_Question_71.png]}}
\noindent \textbf{What does the following code print out?\newline

\newline
\newline

\newline
```python\newline

\newline
print("123" + "abc")\newline

\newline
```}

\begin{enumerate}[label=\Alph*.]
    \item \textbf{123abc} \hfill \textcolor{blue}{\textbf{[Correct]}}
    \item 123+abc
    \item This is a syntax error because you cannot add strings
    \item hello world
\end{enumerate}

\noindent \textbf{Correct Answer:} \textcolor{blue}{\textbf{A}} \\[4pt]
\noindent \textbf{Explanation:}\\
Trong Python, cả `"123"` và `"abc"` đều là các chuỗi ký tự (`str`).\newline
Toán tử cộng (`+`) khi áp dụng giữa hai chuỗi có chức năng là ghép chuỗi (string concatenation), tạo ra một chuỗi mới bằng cách nối nội dung của chuỗi thứ hai vào sau chuỗi thứ nhất.\newline
Kết quả của `"123" + "abc"` là `"123abc"`. Lệnh `print()` in chuỗi này ra màn hình.
\end{minipage}
\vspace{0.8cm}
\hrule
\vspace{0.5cm}

\begin{minipage}{\textwidth}
\subsection*{Question 72 \quad \footnotesize{[0072_q072_Question_72.png]}}
\noindent \textbf{Which of these collections defines a DICTIONARY?}

\begin{enumerate}[label=\Alph*.]
    \item ("apple", "banana", "cherry")
    \item \textbf{"name": "apple", "color": "green"} \hfill \textcolor{blue}{\textbf{[Correct]}}
    \item "apple", "banana", "cherry"
    \item ["apple", "banana", "cherry"]
\end{enumerate}

\noindent \textbf{Correct Answer:} \textcolor{blue}{\textbf{B}} \\[4pt]
\noindent \textbf{Explanation:}\\
Trong Python:\newline
- `dict` (từ điển / dictionary) được định nghĩa trong cặp dấu ngoặc nhọn `{}` với các cặp khóa - giá trị dạng `key: value`, phân tách bằng dấu phẩy. Ví dụ: `{"name": "apple", "color": "green"}`.\newline
- `("apple", "banana", "cherry")` sử dụng dấu ngoặc đơn `()` là kiểu `tuple`.\newline
- `{"apple", "banana", "cherry"}` chứa các phần tử đơn lẻ trong dấu `{}` là kiểu `set` (tập hợp).\newline
- `["apple", "banana", "cherry"]` sử dụng dấu ngoặc vuông `[]` là kiểu `list` (danh sách).
\end{minipage}
\vspace{0.8cm}
\hrule
\vspace{0.5cm}

\begin{minipage}{\textwidth}
\subsection*{Question 73 \quad \footnotesize{[0073_q073_Question_73.png]}}
\noindent \textbf{What is output of the sample code?\newline

\newline
\newline

\newline
```python\newline

\newline
class A:\newline

\newline
    def __init__(self):\newline

\newline
        self.calcI(30)\newline

\newline
\newline

\newline
    def calcI(self, i):\newline

\newline
        self.i = 2 * i\newline

\newline
\newline

\newline
class B(A):\newline

\newline
    def __init__(self):\newline

\newline
        super().__init__()\newline

\newline
        print("i from B is", self.i)\newline

\newline
\newline

\newline
    def calcI(self, i):\newline

\newline
        self.i = 3 * i\newline

\newline
\newline

\newline
b = B()\newline

\newline
```}

\begin{enumerate}[label=\Alph*.]
    \item The __init__ method of only class B gets invoked.
    \item The __init__ method of class A gets invoked and it displays "i from B is 0".
    \item The __init__ method of class A gets invoked and it displays "i from B is 60"
    \item \textbf{The __init__ method of class A gets invoked and it displays "i from B is 90".} \hfill \textcolor{blue}{\textbf{[Correct]}}
\end{enumerate}

\noindent \textbf{Correct Answer:} \textcolor{blue}{\textbf{D}} \\[4pt]
\noindent \textbf{Explanation:}\\
Phân tích cơ chế kế thừa và tính đa hình (polymorphism / dynamic dispatch) trong Python:\newline
1. Khi khởi tạo đối tượng `b = B()`, hàm khởi tạo `B.__init__(self)` được gọi với `self` là đối tượng thuộc lớp `B`.\newline
2. Trong `B.__init__`, lệnh `super().__init__()` kích hoạt hàm khởi tạo `A.__init__(self)` của lớp cha `A`. Do đó, hàm `__init__` của lớp A được gọi (The `__init__` method of class A gets invoked).\newline
3. Trong `A.__init__`, câu lệnh `self.calcI(30)` được gọi. Vì đối tượng `self` thực tế là một thể hiện của lớp con `B`, và lớp `B` đã ghi đè (override) phương thức `calcI`, cơ chế phân giải phương thức động của Python sẽ gọi `B.calcI(self, 30)` chứ không gọi `A.calcI`.\newline
4. Trong `B.calcI`: `self.i = 3 * 30 = 90`.\newline
5. Sau khi `super().__init__()` thực thi xong, luồng điều khiển quay lại `B.__init__` và thực thi dòng tiếp theo: `print("i from B is", self.i)` -> In ra `i from B is 90`.\newline
Vì vậy, hàm `__init__` của lớp A được kích hoạt và in ra "i from B is 90".
\end{minipage}
\vspace{0.8cm}
\hrule
\vspace{0.5cm}

\begin{minipage}{\textwidth}
\subsection*{Question 74 \quad \footnotesize{[0074_q074_Question_74.png]}}
\noindent \textbf{The following code sequence fails with a traceback when the user enters a file that does not exist. How would you avoid the traceback and make it so you could print out your own error message when a bad file name was entered?\newline

\newline
\newline

\newline
```python\newline

\newline
fname = raw_input('Enter the file name: ')\newline

\newline
fhand = open(fname)\newline

\newline
```}

\begin{enumerate}[label=\Alph*.]
    \item \textbf{try / except} \hfill \textcolor{blue}{\textbf{[Correct]}}
    \item signal handlers
    \item try / catch / finally
    \item on error resume next
\end{enumerate}

\noindent \textbf{Correct Answer:} \textcolor{blue}{\textbf{A}} \\[4pt]
\noindent \textbf{Explanation:}\\
Trong Python, cấu trúc để bắt và xử lý ngoại lệ (exceptions) nhằm ngăn chặn traceback làm dừng đột ngột chương trình là `try / except`.\newline
Bằng cách đặt lệnh mở file nguy cơ lỗi `open(fname)` vào khối `try:`, nếu tệp không tồn tại thì Python sẽ ném ngoại lệ (`FileNotFoundError` / `IOError`). Khối `except:` sẽ bắt ngoại lệ này, cho phép lập trình viên in thông báo lỗi tùy chỉnh thân thiện với người dùng.\newline
\newline
- `try / catch / finally` là cú pháp của các ngôn ngữ như Java, JavaScript, C# (Python dùng từ khóa `except` thay vì `catch`).\newline
- `on error resume next` là câu lệnh của Visual Basic.\newline
- `signal handlers` dùng để xử lý tín hiệu hệ điều hành, không phải xử lý ngoại lệ file thông thường trong Python.
\end{minipage}
\vspace{0.8cm}
\hrule
\vspace{0.5cm}

\begin{minipage}{\textwidth}
\subsection*{Question 75 \quad \footnotesize{[0075_q075_Question_75.png]}}
\noindent \textbf{What is the type of the return value of the re.findall() method?}

\begin{enumerate}[label=\Alph*.]
    \item \textbf{A list of strings} \hfill \textcolor{blue}{\textbf{[Correct]}}
    \item An integer
    \item A boolean
    \item A single character
\end{enumerate}

\noindent \textbf{Correct Answer:} \textcolor{blue}{\textbf{A}} \\[4pt]
\noindent \textbf{Explanation:}\\
Theo tài liệu chuẩn của Python cho mô-đun biểu thức chính quy `re`:\newline
Hàm `re.findall(pattern, string, flags=0)` quét toàn bộ chuỗi `string` và tìm tất cả các chuỗi con khớp không chồng lấn với biểu thức chính quy `pattern`. Kết quả trả về luôn là một danh sách các chuỗi ký tự (`list` of `str`). Nếu không có kết quả nào khớp, nó trả về một danh sách rỗng `[]`.\newline
Vì vậy, kiểu dữ liệu trả về là "A list of strings".
\end{minipage}
\vspace{0.8cm}
\hrule
\vspace{0.5cm}

\begin{minipage}{\textwidth}
\subsection*{Question 76 \quad \footnotesize{[0076_q076_Question_76.png]}}
\noindent \textbf{For the following list, how would you print out 'Sally'?\newline

\newline
\newline

\newline
```python\newline

\newline
friends = ['Joseph', 'Glenn', 'Sally']\newline

\newline
```}

\begin{enumerate}[label=\Alph*.]
    \item print(friends['Sally'])
    \item print(friends[2:1])
    \item \textbf{print(friends[2])} \hfill \textcolor{blue}{\textbf{[Correct]}}
    \item print friends[3]
\end{enumerate}

\noindent \textbf{Correct Answer:} \textcolor{blue}{\textbf{C}} \\[4pt]
\noindent \textbf{Explanation:}\\
Trong Python, danh sách (`list`) được đánh chỉ số từ 0 (zero-indexed):\newline
- `friends[0]` là phần tử thứ nhất: `'Joseph'`\newline
- `friends[1]` là phần tử thứ hai: `'Glenn'`\newline
- `friends[2]` là phần tử thứ ba: `'Sally'`\newline
Do đó, câu lệnh chính xác để in ra `'Sally'` là `print(friends[2])`.\newline
\newline
- `friends['Sally']`: Báo lỗi `TypeError` vì chỉ số của list phải là số nguyên hoặc slice.\newline
- `friends[2:1]`: Cắt lát với chỉ số kết thúc nhỏ hơn bắt đầu nên trả về danh sách rỗng `[]`.\newline
- `friends[3]`: Báo lỗi `IndexError: list index out of range` do danh sách chỉ có 3 phần tử (chỉ số tối đa là 2).
\end{minipage}
\vspace{0.8cm}
\hrule
\vspace{0.5cm}

\begin{minipage}{\textwidth}
\subsection*{Question 77 \quad \footnotesize{[0077_q077_Question_77.png]}}
\noindent \textbf{What would the following Python code print out?\newline

\newline
\newline

\newline
```python\newline

\newline
fruit = 'Banana'\newline

\newline
fruit[0] = 'b'\newline

\newline
print(fruit)\newline

\newline
```}

\begin{enumerate}[label=\Alph*.]
    \item b
    \item B
    \item banana
    \item \textbf{Nothing would print - the program fails with a traceback error} \hfill \textcolor{blue}{\textbf{[Correct]}}
\end{enumerate}

\noindent \textbf{Correct Answer:} \textcolor{blue}{\textbf{D}} \\[4pt]
\noindent \textbf{Explanation:}\\
Trong Python, chuỗi ký tự (`str`) là một kiểu dữ liệu bất biến (immutable). Sau khi một chuỗi được tạo, ta không thể thay đổi, xóa hay gán lại giá trị cho từng phần tử/ký tự thông qua chỉ mục (item assignment).\newline
Khi dòng lệnh `fruit[0] = 'b'` chạy, Python sẽ ném ngoại lệ:\newline
`TypeError: 'str' object does not support item assignment`\newline
Vì chương trình gặp lỗi tại dòng này và không có khối `try...except` xử lý, chương trình sẽ dừng thực thi ngay lập tức với một thông báo lỗi traceback trước khi câu lệnh `print(fruit)` có cơ hội chạy. Do đó, không có gì được in ra màn hình.
\end{minipage}
\vspace{0.8cm}
\hrule
\vspace{0.5cm}

\begin{minipage}{\textwidth}
\subsection*{Question 78 \quad \footnotesize{[0078_q078_Question_78.png]}}
\noindent \textbf{What is the output of the following code?\newline

\newline
\newline

\newline
```python\newline

\newline
class Point:\newline

\newline
    def __init__(self, x = 0, y = 0):\newline

\newline
        self.x = x+1\newline

\newline
        self.y = y+1\newline

\newline
\newline

\newline
p1 = Point()\newline

\newline
print(p1.x, p1.y)\newline

\newline
```}

\begin{enumerate}[label=\Alph*.]
    \item None None
    \item Initializes the class for use.
    \item 0 0
    \item \textbf{1 1} \hfill \textcolor{blue}{\textbf{[Correct]}}
    \item X y
\end{enumerate}

\noindent \textbf{Correct Answer:} \textcolor{blue}{\textbf{D}} \\[4pt]
\noindent \textbf{Explanation:}\\
Phân tích quá trình thực thi của đoạn mã:\newline
1. Lớp `Point` định nghĩa hàm khởi tạo `__init__(self, x = 0, y = 0)` có 2 tham số mặc định `x = 0` và `y = 0`.\newline
2. Khi gọi `p1 = Point()`, không có đối số nào được truyền vào nên các tham số mặc định được áp dụng: `x = 0`, `y = 0`.\newline
3. Trong hàm `__init__`:\newline
   - `self.x = x + 1` -> `0 + 1 = 1`\newline
   - `self.y = y + 1` -> `0 + 1 = 1`\newline
4. Lệnh `print(p1.x, p1.y)` in hai giá trị thuộc tính phân tách bởi dấu cách mặc định:\newline
   Kết quả hiển thị trên màn hình là: `1 1`.\newline
Đáp án đúng là D.
\end{minipage}
\vspace{0.8cm}
\hrule
\vspace{0.5cm}

\begin{minipage}{\textwidth}
\subsection*{Question 79 \quad \footnotesize{[0079_q079_Question_79.png]}}
\noindent \textbf{What is the output from this print() function call:\newline

\newline
\newline

\newline
```python\newline

\newline
print(\newline

\newline
    '100 200 300'.count(''),\newline

\newline
    '100 200 300'.count('', 5, 10),\newline

\newline
    '100 200 300'.count('', 5)\newline

\newline
)\newline

\newline
```\newline

\newline
*(Lưu ý: Trong đề thi gốc, do chuỗi chứa các ký tự `$` nên hệ thống web đã tự động render thành các thẻ KaTeX/MathML `<span class="katex">...`)*}

\begin{enumerate}[label=\Alph*.]
    \item 310
    \item 311
    \item 321
    \item \textbf{312} \hfill \textcolor{blue}{\textbf{[Correct]}}
\end{enumerate}

\noindent \textbf{Correct Answer:} \textcolor{blue}{\textbf{D}} \\[4pt]
\noindent \textbf{Explanation:}\\
Chuỗi xét là `s = '$100 $200 $300'`.\newline
Các vị trí xuất hiện của ký tự `$` trong chuỗi gồm có chỉ số 0, 5 và 10.\newline
Phương thức `str.count(sub[, start[, end]])` trong Python đếm số lần xuất hiện không chồng lặp của chuỗi con `sub` trong phạm vi từ chỉ số `start` đến `end - 1`:\newline
1. `s.count('$')`: Đếm trên toàn bộ chuỗi -> Có 3 ký tự `$` (tại các vị trí 0, 5, 10) => Kết quả là `3`.\newline
2. `s.count('$', 5, 10)`: Đếm trong phạm vi chỉ số `[5, 10)` (từ vị trí 5 đến 9) -> Chỉ có 1 ký tự `$` (tại vị trí 5, vị trí 10 nằm ngoài khoảng) => Kết quả là `1`.\newline
3. `s.count('$', 5)`: Đếm từ chỉ số 5 đến hết chuỗi -> Có 2 ký tự `$` (tại vị trí 5 và 10) => Kết quả là `2`.\newline
Ghép các giá trị kết quả lại theo các phương án trắc nghiệm của đề thi là `312`.\newline
\newline
*Ghi chú thêm về giao diện đề thi*: Do câu hỏi gốc chứa chuỗi có ký hiệu `$` nên trình render Markdown/LaTeX của hệ thống thi đã nhầm là công thức toán và sinh mã thẻ HTML `<span class="katex">...`. Bản chất câu hỏi là kiểm tra kiến thức về phương thức `str.count()`.
\end{minipage}
\vspace{0.8cm}
\hrule
\vspace{0.5cm}

\begin{minipage}{\textwidth}
\subsection*{Question 80 \quad \footnotesize{[0080_q080_Question_80.png]}}
\noindent \textbf{What does this Python statement do if PartyAnimal is a class? zap = PartyAnimal()}

\begin{enumerate}[label=\Alph*.]
    \item Clear out all the data in the PartyAnimal variable and put the old values for the data in zap
    \item \textbf{Use the PartyAnimal template to make a new object and assign it to zap} \hfill \textcolor{blue}{\textbf{[Correct]}}
    \item Subtract the value of the zap variable from the value in the PartyAnimal variable and put the difference in zap
    \item Copy the value from the PartyAnimal variable to the variable zap
\end{enumerate}

\noindent \textbf{Correct Answer:} \textcolor{blue}{\textbf{B}} \\[4pt]
\noindent \textbf{Explanation:}\\
Trong lập trình hướng đối tượng (OOP) với Python:\newline
- Lớp (class) đóng vai trò là một bản thiết kế hoặc khuôn mẫu (blueprint / template) xác định cấu trúc thuộc tính và hành vi.\newline
- Khi câu lệnh `zap = PartyAnimal()` được thực thi, Python tiến hành khởi tạo một đối tượng (instantiation): sử dụng khuôn mẫu `PartyAnimal` để tạo ra một đối tượng (instance) mới và gán tham chiếu của đối tượng đó cho biến `zap`.\newline
Do đó, phát biểu chính xác là: "Use the PartyAnimal template to make a new object and assign it to zap".
\end{minipage}
\vspace{0.8cm}
\hrule
\vspace{0.5cm}

\begin{minipage}{\textwidth}
\subsection*{Question 81 \quad \footnotesize{[0081_q081_Question_81.png]}}
\noindent \textbf{When reading a file using the file object, what method is best for reading the entire file into a single string?}

\begin{enumerate}[label=\Alph*.]
    \item \textbf{.read()} \hfill \textcolor{blue}{\textbf{[Correct]}}
    \item .readline()
    \item .readlines()
    \item .read_file_to_str()
\end{enumerate}

\noindent \textbf{Correct Answer:} \textcolor{blue}{\textbf{A}} \\[4pt]
\noindent \textbf{Explanation:}\\
Trong Python, đối tượng tệp (file object) cung cấp các phương thức đọc dữ liệu chính sau:\newline
- `read([size])`: Đọc toàn bộ nội dung của tệp (hoặc tối đa `size` ký tự/byte) và trả về dưới dạng một chuỗi (`str`) duy nhất nếu mở ở chế độ văn bản (text mode) hoặc `bytes` nếu mở ở chế độ nhị phân (binary mode). Đây là phương thức chuẩn và phù hợp nhất để đọc toàn bộ tệp vào một chuỗi đơn lẻ.\newline
- `readline([size])`: Chỉ đọc một dòng duy nhất tính từ vị trí con trỏ hiện tại cho đến ký tự xuống dòng `\n`.\newline
- `readlines([hint])`: Đọc toàn bộ các dòng còn lại trong tệp và trả về một danh sách (list) các chuỗi, trong đó mỗi chuỗi là một dòng.\newline
- `.read_file_to_str()`: Không phải là một phương thức tích hợp sẵn của file object trong Python.
\end{minipage}
\vspace{0.8cm}
\hrule
\vspace{0.5cm}

\begin{minipage}{\textwidth}
\subsection*{Question 82 \quad \footnotesize{[0082_q082_Question_82.png]}}
\noindent \textbf{How would you use the index operator [] to print out the letter q from the following string? x = 'From marquard@uct.ac.za'}

\begin{enumerate}[label=\Alph*.]
    \item print x[10]
    \item \textbf{print(x[8])} \hfill \textcolor{blue}{\textbf{[Correct]}}
    \item print(x[9])
    \item print(x[q])
\end{enumerate}

\noindent \textbf{Correct Answer:} \textcolor{blue}{\textbf{B}} \\[4pt]
\noindent \textbf{Explanation:}\\
Trong Python, chuỗi ký tự (string) được đánh chỉ mục (zero-based index) bắt đầu từ `0`:\newline
Xét chuỗi: `x = 'From marquard@uct.ac.za'`\newline
- `x[0]` = 'F'\newline
- `x[1]` = 'r'\newline
- `x[2]` = 'o'\newline
- `x[3]` = 'm'\newline
- `x[4]` = ' ' (khoảng trắng)\newline
- `x[5]` = 'm'\newline
- `x[6]` = 'a'\newline
- `x[7]` = 'r'\newline
- `x[8]` = 'q'\newline
- `x[9]` = 'u'\newline
- `x[10]` = 'a'\newline
\newline
Ký tự `'q'` nằm ở vị trí chỉ mục 8 (`index 8`). Cú pháp chuẩn trong Python 3 để in ký tự này là `print(x[8])`.\newline
- Đáp án A (`print x[10]`): chỉ mục 10 là ký tự `'a'`, đồng thời thiếu dấu ngoặc đơn của hàm `print()` trong Python 3.\newline
- Đáp án C (`print(x[9])`): chỉ mục 9 là ký tự `'u'`.\newline
- Đáp án D (`print(x[q])`): `q` không được đặt trong dấu ngoặc kép/đơn nên được hiểu là biến, biến `q` chưa định nghĩa sẽ gây lỗi `NameError`.
\end{minipage}
\vspace{0.8cm}
\hrule
\vspace{0.5cm}

\begin{minipage}{\textwidth}
\subsection*{Question 83 \quad \footnotesize{[0083_q083_Question_83.png]}}
\noindent \textbf{What is true about the following code segment:\newline

\newline
\newline

\newline
```python\newline

\newline
if x == 5:\newline

\newline
    print('Is 5')\newline

\newline
    print('Is Still 5')\newline

\newline
    print('Third 5')\newline

\newline
```}

\begin{enumerate}[label=\Alph*.]
    \item The string 'Is 5' will never print out regardless of the value for x.
    \item The string 'Is 5' will always print out regardless of the value for x.
    \item Only two of the three print statements will print out if the value of x is less than zero.
    \item \textbf{Depending on the value of x, either all three of the print statements will execute or none of the statements will execute} \hfill \textcolor{blue}{\textbf{[Correct]}}
\end{enumerate}

\noindent \textbf{Correct Answer:} \textcolor{blue}{\textbf{D}} \\[4pt]
\noindent \textbf{Explanation:}\\
Trong Python, các khối lệnh (code blocks) được phân định bằng thụt lề (indentation):\newline
- Cả ba câu lệnh `print` (`print('Is 5')`, `print('Is Still 5')`, và `print('Third 5')`) đều được thụt lề cùng cấp và nằm hoàn toàn bên trong khối điều kiện của `if x == 5:`.\newline
- Khi biểu thức điều kiện `x == 5` trả về `True` (nghĩa là giá trị của `x` bằng 5), toàn bộ khối lệnh gồm cả ba câu lệnh `print` sẽ được thực thi tuần tự.\newline
- Khi biểu thức điều kiện `x == 5` trả về `False` (giá trị của `x` khác 5), Python sẽ bỏ qua toàn bộ khối lệnh thụt lề này và không in ra bất kỳ dòng nào.\newline
- Vì vậy, tùy thuộc vào giá trị của `x`, hoặc cả ba câu lệnh print đều được thực thi, hoặc không có câu lệnh nào được thực thi ("Depending on the value of x, either all three of the print statements will execute or none of the statements will execute").
\end{minipage}
\vspace{0.8cm}
\hrule
\vspace{0.5cm}

\begin{minipage}{\textwidth}
\subsection*{Question 84 \quad \footnotesize{[0084_q084_Question_84.png]}}
\noindent \textbf{Which came first, the instance or the class?}

\begin{enumerate}[label=\Alph*.]
    \item \textbf{Class} \hfill \textcolor{blue}{\textbf{[Correct]}}
    \item Method
    \item Function
    \item Instance
\end{enumerate}

\noindent \textbf{Correct Answer:} \textcolor{blue}{\textbf{A}} \\[4pt]
\noindent \textbf{Explanation:}\\
Trong lập trình hướng đối tượng (Object-Oriented Programming - OOP):\newline
- Lớp (Class) là khuôn mẫu hoặc bản thiết kế (blueprint/template) định nghĩa các thuộc tính và hành vi (phương thức) chung của một nhóm đối tượng.\newline
- Thực thể (Instance) là một đối tượng cụ thể được sinh ra và cấp phát bộ nhớ dựa trên khuôn mẫu của lớp đó.\newline
- Để có thể tạo ra một thực thể cụ thể (instance), bản thiết kế (class) phải được định nghĩa trước trong mã nguồn. Do đó, "Class" luôn xuất hiện và được định nghĩa trước "Instance".
\end{minipage}
\vspace{0.8cm}
\hrule
\vspace{0.5cm}

\begin{minipage}{\textwidth}
\subsection*{Question 85 \quad \footnotesize{[0085_q085_Question_85.png]}}
\noindent \textbf{When you have multiple lines in an if block, how do you indicate the end of the if block?}

\begin{enumerate}[label=\Alph*.]
    \item You use a curly brace  after the last line of the if block
    \item You omit the semicolon ; on the last line of the if block
    \item You put the colon : character on a line by itself to indicate we are done with the if block
    \item \textbf{You de-indent the next line past the if block to the same level of indent as the original if statement} \hfill \textcolor{blue}{\textbf{[Correct]}}
\end{enumerate}

\noindent \textbf{Correct Answer:} \textcolor{blue}{\textbf{D}} \\[4pt]
\noindent \textbf{Explanation:}\\
Python không sử dụng các dấu bao đóng như cặp ngoặc nhọn `{}` hay các từ khóa kết thúc khối lệnh như `end` / `endif`:\newline
- Python dựa hoàn toàn vào khoảng trắng thụt lề (indentation) để xác định phạm vi và cấu trúc phân cấp của các khối mã lệnh.\newline
- Khi muốn kết thúc một khối lệnh (chẳng hạn như khối `if`), lập trình viên biểu thị điều này bằng cách lùi thụt lề (de-indent) ở dòng tiếp theo về lại cùng mức thụt lề với câu lệnh `if` ban đầu.\newline
- Các lựa chọn A, B, C đều sai vì không đúng quy chuẩn cú pháp của Python.
\end{minipage}
\vspace{0.8cm}
\hrule
\vspace{0.5cm}

\begin{minipage}{\textwidth}
\subsection*{Question 86 \quad \footnotesize{[0086_q086_Question_86.png]}}
\noindent \textbf{Which of these operators is not a comparison / logical operator?}

\begin{enumerate}[label=\Alph*.]
    \item ==
    \item >=
    \item <=
    \item \textbf{=} \hfill \textcolor{blue}{\textbf{[Correct]}}
\end{enumerate}

\noindent \textbf{Correct Answer:} \textcolor{blue}{\textbf{D}} \\[4pt]
\noindent \textbf{Explanation:}\\
Trong ngôn ngữ lập trình Python:\newline
- Dấu `=` là toán tử gán (assignment operator), được sử dụng để gán giá trị của biểu thức ở vế phải cho biến ở vế trái (ví dụ: `x = 10`), không mang ý nghĩa so sánh hay kiểm tra logic.\newline
- `==` là toán tử so sánh bằng (equality comparison operator), trả về `True` nếu hai vế bằng nhau và `False` nếu khác nhau.\newline
- `>=` là toán tử so sánh lớn hơn hoặc bằng (greater than or equal to).\newline
- `<=` là toán tử so sánh nhỏ hơn hoặc bằng (less than or equal to).\newline
Vì vậy, toán tử `=` không phải là toán tử so sánh hay logic.
\end{minipage}
\vspace{0.8cm}
\hrule
\vspace{0.5cm}

\begin{minipage}{\textwidth}
\subsection*{Question 87 \quad \footnotesize{[0087_q087_Question_87.png]}}
\noindent \textbf{Which of these words is a reserved word in Python ?}

\begin{enumerate}[label=\Alph*.]
    \item payroll
    \item pizza
    \item \textbf{while} \hfill \textcolor{blue}{\textbf{[Correct]}}
    \item names
\end{enumerate}

\noindent \textbf{Correct Answer:} \textcolor{blue}{\textbf{C}} \\[4pt]
\noindent \textbf{Explanation:}\\
Từ khóa dành riêng (reserved words hay keywords) trong Python là các từ có ý nghĩa cú pháp cố định được định nghĩa trước bởi trình thông dịch, không thể dùng làm tên định danh (tên biến, hàm, lớp...):\newline
- `while` là một từ khóa chuẩn trong Python dùng để khởi tạo vòng lặp điều kiện `while` (có thể kiểm tra qua `import keyword; keyword.kwlist`).\newline
- Các từ `payroll`, `pizza`, `names` chỉ là các danh từ thông thường trong tiếng Anh, không phải từ khóa của ngôn ngữ và hoàn toàn có thể được dùng làm tên biến hợp lệ.
\end{minipage}
\vspace{0.8cm}
\hrule
\vspace{0.5cm}

\begin{minipage}{\textwidth}
\subsection*{Question 88 \quad \footnotesize{[0088_q088_Question_88.png]}}
\noindent \textbf{Which of the following can be used to invoke the __init__ method in B from A, where A is a subclass of B?}

\begin{enumerate}[label=\Alph*.]
    \item \textbf{super().__init__()} \hfill \textcolor{blue}{\textbf{[Correct]}}
    \item B.__init__()
    \item super().__init__(self)
\end{enumerate}

\noindent \textbf{Correct Answer:} \textcolor{blue}{\textbf{A}} \\[4pt]
\noindent \textbf{Explanation:}\\
Trong lập trình hướng đối tượng với Python 3:\newline
- Khi lớp `A` là lớp con kế thừa từ lớp `B` (`class A(B):`), để gọi phương thức khởi tạo `__init__` của lớp cha `B` từ lớp con `A`, phương pháp chuẩn và khuyến nghị nhất là sử dụng hàm tích hợp sẵn `super()`: `super().__init__()`.\newline
- Đối tượng `super()` tự động liên kết với thực thể hiện tại (`self`) và truyền nó vào phương thức của lớp cha, vì vậy không cần và không được truyền lại `self` vào trong ngoặc đơn. Do đó đáp án C `super().__init__(self)` là sai và sẽ gây lỗi `TypeError`.\newline
- Đối với đáp án B `B.__init__()`: Nếu gọi trực tiếp qua tên lớp `B`, phương thức `__init__` là một hàm thông thường yêu cầu phải truyền đối số thực thể đầu tiên là `self`: `B.__init__(self)`. Nếu viết `B.__init__()` mà không truyền `self`, Python sẽ ném lỗi `TypeError: __init__() missing 1 required positional argument: 'self'`.\newline
Vì vậy, chỉ có đáp án A là chính xác.
\end{minipage}
\vspace{0.8cm}
\hrule
\vspace{0.5cm}

\begin{minipage}{\textwidth}
\subsection*{Question 89 \quad \footnotesize{[0089_q089_Question_89.png]}}
\noindent \textbf{Which statement is used to stop a loop?}

\begin{enumerate}[label=\Alph*.]
    \item exit
    \item return
    \item stop
    \item \textbf{break} \hfill \textcolor{blue}{\textbf{[Correct]}}
\end{enumerate}

\noindent \textbf{Correct Answer:} \textcolor{blue}{\textbf{D}} \\[4pt]
\noindent \textbf{Explanation:}\\
Trong Python:\newline
- Câu lệnh `break` là câu lệnh điều khiển vòng lặp được thiết kế đặc thù để ngay lập tức chấm dứt (dừng hoàn toàn) vòng lặp `for` hoặc `while` gần nhất chứa nó. Khi gặp `break`, luồng thực thi sẽ thoát khỏi vòng lặp và chuyển đến câu lệnh tiếp theo ngay sau khối lặp.\newline
- `exit`: Là hàm dùng để kết thúc toàn bộ chương trình hoặc tiến trình Python (như `sys.exit()` hoặc built-in `exit()`), không phải câu lệnh chỉ để dừng một vòng lặp.\newline
- `return`: Dùng để kết thúc thực thi của một hàm và trả về giá trị cho nơi gọi hàm (chỉ dùng hợp lệ bên trong thân hàm).\newline
- `stop`: Không phải là một từ khóa hay câu lệnh của Python.
\end{minipage}
\vspace{0.8cm}
\hrule
\vspace{0.5cm}

\begin{minipage}{\textwidth}
\subsection*{Question 90 \quad \footnotesize{[0090_q090_Question_90.png]}}
\noindent \textbf{You want to create a bytes object consisting of five null (0x00) bytes. All of the following will work except one. Select the one that doesn't work:}

\begin{enumerate}[label=\Alph*.]
    \item \textbf{bytes(0, 0, 0, 0, 0)} \hfill \textcolor{blue}{\textbf{[Correct]}}
    \item bytes(5)
    \item bytes([0] * 5)
    \item bytes('x 00x 00x 00x 00x 00', 'utf-8')
\end{enumerate}

\noindent \textbf{Correct Answer:} \textcolor{blue}{\textbf{A}} \\[4pt]
\noindent \textbf{Explanation:}\\
Hàm khởi tạo đối tượng `bytes` trong Python có cú pháp `bytes([source[, encoding[, errors]]])`, nhận tối đa 3 tham số:\newline
- Đáp án A (`bytes(0, 0, 0, 0, 0)`): Truyền 5 tham số số nguyên riêng biệt. Do hàm `bytes` không chấp nhận 5 tham số, Python sẽ lập tức báo lỗi `TypeError: bytes() takes at most 3 arguments (5 given)`. Do đó đây là cách KHÔNG hoạt động (doesn't work).\newline
- Đáp án B (`bytes(5)`): Khi truyền một số nguyên `n` vào `bytes()`, Python sẽ tạo một đối tượng `bytes` có độ dài `n` với tất cả các byte đều có giá trị bằng 0 (`b'\x00\x00\x00\x00\x00'`). Cách này hoạt động tốt.\newline
- Đáp án C (`bytes([0] * 5)`): Biểu thức `[0] * 5` tạo ra danh sách `[0, 0, 0, 0, 0]`. Khi truyền một iterable chứa các số nguyên trong khoảng 0-255 vào `bytes()`, nó sẽ tạo ra đối tượng bytes tương ứng `b'\x00\x00\x00\x00\x00'`. Cách này hoạt động tốt.\newline
- Đáp án D (`bytes('\x00\x00\x00\x00\x00', 'utf-8')`): Mã hóa chuỗi ký tự null thành bytes sử dụng bảng mã utf-8, tạo ra kết quả `b'\x00\x00\x00\x00\x00'`. Cách này hoạt động tốt.
\end{minipage}
\vspace{0.8cm}
\hrule
\vspace{0.5cm}

\begin{minipage}{\textwidth}
\subsection*{Question 91 \quad \footnotesize{[0091_q091_Question_91.png]}}
\noindent \textbf{You want to create a bytes object consisting of five null (0x00) bytes. Choose all the apply.}

\begin{enumerate}[label=\Alph*.]
    \item \textbf{bytes('x 00x 00x 00x 00x 00', 'utf-8')} \hfill \textcolor{blue}{\textbf{[Correct]}}
    \item bytes(0, 0, 0, 0, 0)
    \item bytes[0:5]
    \item \textbf{bytes(5)} \hfill \textcolor{blue}{\textbf{[Correct]}}
\end{enumerate}

\noindent \textbf{Correct Answer:} \textcolor{blue}{\textbf{A, D}} \\[4pt]
\noindent \textbf{Explanation:}\\
Để tạo một đối tượng bytes gồm 5 byte có giá trị null (`0x00`), các phương án hợp lệ gồm:\newline
- Đáp án A (`bytes('\x00\x00\x00\x00\x00', 'utf-8')`): Mã hóa chuỗi gồm 5 ký tự null `\x00` sang định dạng bytes với mã hóa utf-8. Kết quả thu được là `b'\x00\x00\x00\x00\x00'`, gồm chính xác 5 byte null.\newline
- Đáp án D (`bytes(5)`): Truyền số nguyên `5` vào constructor `bytes()`, Python sẽ khởi tạo một đối tượng `bytes` có độ dài bằng 5 với các byte được lấp đầy bằng giá trị 0 (`b'\x00\x00\x00\x00\x00'`).\newline
- Đáp án B sai vì `bytes()` chỉ nhận tối đa 3 tham số, truyền 5 tham số rời rạc sẽ gây ngoại lệ `TypeError: bytes() takes at most 3 arguments (5 given)`.\newline
- Đáp án C sai vì `bytes[0:5]` là cú pháp slicing trên lớp `bytes`, không phải cú pháp gọi hàm tạo đối tượng, không tạo ra đối tượng bytes có 5 byte null.
\end{minipage}
\vspace{0.8cm}
\hrule
\vspace{0.5cm}

\begin{minipage}{\textwidth}
\subsection*{Question 92 \quad \footnotesize{[0092_q092_Question_92.png]}}
\noindent \textbf{Which of the following statement about Python is not correct?}

\begin{enumerate}[label=\Alph*.]
    \item There are libraries and APIs available to access many of the popular databases from python
    \item \textbf{Due to its proprietary nature, database access from python is not available for many database} \hfill \textcolor{blue}{\textbf{[Correct]}}
    \item Python is a popular scripting language for connecting and accessing databases
    \item The python ecosystem is very rich and provides easy to use tools for data science
\end{enumerate}

\noindent \textbf{Correct Answer:} \textcolor{blue}{\textbf{B}} \\[4pt]
\noindent \textbf{Explanation:}\\
Phát biểu B là phát biểu KHÔNG chính xác (not correct) vì:\newline
- Python là một ngôn ngữ lập trình mã nguồn mở hoàn toàn miễn phí (open-source), được quản lý bởi Python Software Foundation, hoàn toàn KHÔNG phải là một ngôn ngữ "độc quyền" (proprietary).\newline
- Hơn nữa, Python có hệ sinh thái thư viện kết nối cơ sở dữ liệu cực kỳ phong phú và mạnh mẽ thông qua chuẩn Python Database API Specification (PEP 249), cho phép kết nối đến hầu như mọi cơ sở dữ liệu phổ biến hiện nay như SQLite, PostgreSQL, MySQL, SQL Server, Oracle, v.v.\newline
- Các khẳng định A, C, D đều mô tả rất chính xác về ngôn ngữ và hệ sinh thái Python.
\end{minipage}
\vspace{0.8cm}
\hrule
\vspace{0.5cm}

\begin{minipage}{\textwidth}
\subsection*{Question 93 \quad \footnotesize{[0093_q093_Question_93.png]}}
\noindent \textbf{What is the output of print list[1:3] if list = [ 'abcd', 786 , 2.23, 'john', 70.2 ]?}

\begin{enumerate}[label=\Alph*.]
    \item [ 'abcd', 786 , 2.23, 'john', 70.2 ]
    \item None of the above.
    \item \textbf{[786, 2.23]} \hfill \textcolor{blue}{\textbf{[Correct]}}
    \item abcd
\end{enumerate}

\noindent \textbf{Correct Answer:} \textcolor{blue}{\textbf{C}} \\[4pt]
\noindent \textbf{Explanation:}\\
Trong Python, cú pháp cắt lát danh sách (list slice) có dạng `list[start:stop]`:\newline
- Lát cắt sẽ lấy các phần tử từ chỉ mục `start` đến chỉ mục `stop - 1` (không bao gồm vị trí `stop`).\newline
- Xét danh sách: `list = ['abcd', 786, 2.23, 'john', 70.2]`:\newline
  + Vị trí chỉ mục 0: `'abcd'`\newline
  + Vị trí chỉ mục 1: `786`\newline
  + Vị trí chỉ mục 2: `2.23`\newline
  + Vị trí chỉ mục 3: `'john'`\newline
  + Vị trí chỉ mục 4: `70.2`\newline
- Biểu thức `list[1:3]` sẽ trích xuất các phần tử ở vị trí index `1` và `2`, trả về danh sách con mới: `[786, 2.23]`.
\end{minipage}
\vspace{0.8cm}
\hrule
\vspace{0.5cm}

\begin{minipage}{\textwidth}
\subsection*{Question 94 \quad \footnotesize{[0094_q094_Question_94.png]}}
\noindent \textbf{In the following code, x = 9 What is x?}

\begin{enumerate}[label=\Alph*.]
    \item A pointer
    \item \textbf{A variable} \hfill \textcolor{blue}{\textbf{[Correct]}}
    \item A parameter
    \item A constant
\end{enumerate}

\noindent \textbf{Correct Answer:} \textcolor{blue}{\textbf{B}} \\[4pt]
\noindent \textbf{Explanation:}\\
Trong câu lệnh Python `x = 9`:\newline
- `x` là tên của một biến (a variable) được sử dụng để định danh và tham chiếu đến đối tượng số nguyên trong bộ nhớ.\newline
- `=` là toán tử gán (assignment operator).\newline
- `9` là giá trị hằng số nguyên (literal constant).\newline
- Phân tích các lựa chọn khác:\newline
  + Pointer (con trỏ): Python quản lý tham chiếu ngầm định, lập trình viên không trực tiếp thao tác biến như một con trỏ rõ ràng như trong C/C++.\newline
  + Parameter (tham số): `x` không nằm trong danh sách khai báo tham số của một định nghĩa hàm `def func(...)`.\newline
  + Constant (hằng số): `x` là biến số có thể gán lại giá trị khác bất kỳ lúc nào, không phải là hằng số.
\end{minipage}
\vspace{0.8cm}
\hrule
\vspace{0.5cm}

\begin{minipage}{\textwidth}
\subsection*{Question 95 \quad \footnotesize{[0095_q095_Question_95.png]}}
\noindent \textbf{```python\newline

\newline
class Person:\newline

\newline
    def __init__(self, name, age):\newline

\newline
        self.name = name\newline

\newline
        self.age = age\newline

\newline
```\newline

\newline
\newline

\newline
The correct way to instantiate the above Person class is:}

\begin{enumerate}[label=\Alph*.]
    \item \textbf{Person ("Rufus", 3)} \hfill \textcolor{blue}{\textbf{[Correct]}}
    \item Person ()
    \item Person.create("Rufus", 3)
    \item Person.init("Rufus", 3)
\end{enumerate}

\noindent \textbf{Correct Answer:} \textcolor{blue}{\textbf{A}} \\[4pt]
\noindent \textbf{Explanation:}\\
Trong Python:\newline
- Để khởi tạo một thể hiện (instantiate an object) từ một lớp, ta sử dụng cú pháp gọi hàm với tên lớp: `ClassName(arguments...)`.\newline
- Phương thức `__init__` của lớp `Person` được định nghĩa với các tham số: `(self, name, age)`. Tham số `self` được Python tự động truyền vào đại diện cho đối tượng mới tạo, vì vậy người gọi cần truyền đủ 2 đối số tương ứng với `name` và `age`.\newline
- Cú pháp chính xác là: `Person("Rufus", 3)`.\newline
- Phân tích các đáp án khác:\newline
  + B (`Person()`): Không truyền đủ 2 tham số bắt buộc `name` và `age`, sẽ gây lỗi `TypeError: __init__() missing 2 required positional arguments: 'name' and 'age'`.\newline
  + C (`Person.create("Rufus", 3)`): Lớp `Person` không định nghĩa phương thức `create`, sẽ gây lỗi `AttributeError`.\newline
  + D (`Person.init("Rufus", 3)`): Phương thức khởi tạo của Python có tên là `__init__` (với 2 dấu gạch dưới ở mỗi bên), không phải `init`.
\end{minipage}
\vspace{0.8cm}
\hrule
\vspace{0.5cm}

\begin{minipage}{\textwidth}
\subsection*{Question 96 \quad \footnotesize{[0096_q096_Question_96.png]}}
\noindent \textbf{Which of the following data structures in python are immutable?}

\begin{enumerate}[label=\Alph*.]
    \item Linked List
    \item All of the above
    \item \textbf{Tuple} \hfill \textcolor{blue}{\textbf{[Correct]}}
    \item Tree
\end{enumerate}

\noindent \textbf{Correct Answer:} \textcolor{blue}{\textbf{C}} \\[4pt]
\noindent \textbf{Explanation:}\\
Trong Python:\newline
- Một đối tượng hoặc cấu trúc dữ liệu được gọi là bất biến (immutable) nếu nội dung, giá trị của nó không thể bị thay đổi sau khi được tạo ra trong bộ nhớ.\newline
- `Tuple` là cấu trúc dữ liệu tích hợp sẵn (built-in sequence) có tính chất bất biến (immutable). Một khi tuple đã được khởi tạo, ta không thể thêm, bớt hoặc gán lại giá trị cho các phần tử của nó (`TypeError: 'tuple' object does not support item assignment`).\newline
- Linked List (Danh sách liên kết) và Tree (Cây) là các cấu trúc dữ liệu động có thể thêm, sửa, xóa các nút (nodes) và liên kết giữa chúng, vì vậy chúng có tính chất biến đổi (mutable).\newline
- Do đó, Tuple là đáp án chính xác.
\end{minipage}
\vspace{0.8cm}
\hrule
\vspace{0.5cm}

\begin{minipage}{\textwidth}
\subsection*{Question 97 \quad \footnotesize{[0097_q097_Question_97.png]}}
\noindent \textbf{Which data structure is the output of the code given below?\newline

\newline
\newline

\newline
```python\newline

\newline
# Code starts\newline

\newline
col = (2,4,6,3,5,8,7,6,4,3,0,9)\newline

\newline
l = list(col)\newline

\newline
\newline

\newline
finalAns = list(set(col))\newline

\newline
print(type(finalAns))\newline

\newline
# Code Ends\newline

\newline
```}

\begin{enumerate}[label=\Alph*.]
    \item \textbf{List} \hfill \textcolor{blue}{\textbf{[Correct]}}
    \item Set
    \item Dictionary
    \item Tree
\end{enumerate}

\noindent \textbf{Correct Answer:} \textcolor{blue}{\textbf{A}} \\[4pt]
\noindent \textbf{Explanation:}\\
Phân tích từng bước thực thi của đoạn mã:\newline
1. `col = (2,4,6,3,5,8,7,6,4,3,0,9)`: Tạo một tuple gồm các số nguyên.\newline
2. `l = list(col)`: Chuyển đổi tuple thành một list `l`.\newline
3. `set(col)`: Chuyển đổi tuple `col` thành một tập hợp (set), thao tác này loại bỏ các phần tử trùng lặp: `{0, 2, 3, 4, 5, 6, 7, 8, 9}`.\newline
4. `finalAns = list(...)`: Hàm `list()` nhận tập hợp vừa tạo và chuyển đổi nó thành kiểu danh sách (`list`).\newline
5. `print(type(finalAns))`: Hàm `type(finalAns)` kiểm tra kiểu dữ liệu của biến `finalAns`, kết quả in ra màn hình là `<class 'list'>`.\newline
Do đó, cấu trúc dữ liệu kết quả là List.
\end{minipage}
\vspace{0.8cm}
\hrule
\vspace{0.5cm}

\begin{minipage}{\textwidth}
\subsection*{Question 98 \quad \footnotesize{[0098_q098_Question_98.png]}}
\noindent \textbf{In Python, a function within a class definition is called a:}

\begin{enumerate}[label=\Alph*.]
    \item a class function
    \item an operation
    \item a callable
    \item a factory
    \item \textbf{a method} \hfill \textcolor{blue}{\textbf{[Correct]}}
\end{enumerate}

\noindent \textbf{Correct Answer:} \textcolor{blue}{\textbf{E}} \\[4pt]
\noindent \textbf{Explanation:}\\
Trong lập trình hướng đối tượng (OOP) nói chung và Python nói riêng:\newline
- Một hàm (function) được định nghĩa bên trong phần thân của một lớp (class definition) được gọi là một phương thức (method).\newline
- Phương thức thể hiện các hành vi của đối tượng hoặc của lớp, với tham số đầu tiên thường là `self` (đối với phương thức của thực thể - instance method) hoặc `cls` (đối với phương thức của lớp - class method).\newline
- Do đó đáp án chính xác là E: "a method".
\end{minipage}
\vspace{0.8cm}
\hrule
\vspace{0.5cm}

\begin{minipage}{\textwidth}
\subsection*{Question 99 \quad \footnotesize{[0099_q099_Question_99.png]}}
\noindent \textbf{What is the syntax to look up the fullname attribute in an object stored in the variable colleen?}

\begin{enumerate}[label=\Alph*.]
    \item colleen::fullname
    \item \textbf{colleen.fullname} \hfill \textcolor{blue}{\textbf{[Correct]}}
    \item colleen->fullname
    \item colleen['fullname']
\end{enumerate}

\noindent \textbf{Correct Answer:} \textcolor{blue}{\textbf{B}} \\[4pt]
\noindent \textbf{Explanation:}\\
Trong Python:\newline
- Để truy cập (lookup) một thuộc tính (attribute) hoặc phương thức (method) của một đối tượng, cú pháp chuẩn của ngôn ngữ là sử dụng toán tử dấu chấm (dot operator `.`): `<đối_tượng>.<tên_thuộc_tính>`.\newline
- Vì vậy, để lấy thuộc tính `fullname` từ đối tượng lưu trong biến `colleen`, ta viết `colleen.fullname`.\newline
- Phân tích các lựa chọn khác:\newline
  + `colleen::fullname`: Toán tử `::` được sử dụng trong C++, Rust, PHP (Scope Resolution Operator), không tồn tại trong cú pháp Python.\newline
  + `colleen->fullname`: Toán tử `->` truy cập con trỏ được dùng trong C, C++, PHP; trong Python `->` chỉ dùng cho type annotation (gợi ý kiểu trả về của hàm).\newline
  + `colleen['fullname']`: Cú pháp truy cập theo khóa (subscripting), thường dùng cho từ điển (`dict`) hoặc ánh xạ (mapping), không phải cú pháp truy cập thuộc tính thông thường của một object.
\end{minipage}
\vspace{0.8cm}
\hrule
\vspace{0.5cm}

\begin{minipage}{\textwidth}
\subsection*{Question 100 \quad \footnotesize{[0100_q100_Question_100.png]}}
\noindent \textbf{Python strings have a property called "immutability." What does this mean?}

\begin{enumerate}[label=\Alph*.]
    \item \textbf{Strings in Python can't be changed} \hfill \textcolor{blue}{\textbf{[Correct]}}
    \item You can update a string in Python with concatenation
    \item Strings in Python can be represented as arrays of chars
    \item Strings can't be divided by numbers
\end{enumerate}

\noindent \textbf{Correct Answer:} \textcolor{blue}{\textbf{A}} \\[4pt]
\noindent \textbf{Explanation:}\\
Thuộc tính "bất biến" (immutability) trong Python:\newline
- Khái niệm "bất biến" (immutability) trong Python chỉ đặc tính của một đối tượng mà giá trị nội tại của nó không thể bị thay đổi sau khi được khởi tạo trong bộ nhớ.\newline
- Chuỗi (`str`) trong Python là một kiểu dữ liệu bất biến:\newline
  + Bạn không thể gán lại giá trị cho từng ký tự trong chuỗi thông qua chỉ số, ví dụ: lệnh `s[0] = 'H'` sẽ ném lỗi `TypeError: 'str' object does not support item assignment`.\newline
  + Mọi phương thức xử lý chuỗi (như `.replace()`, `.upper()`, phép nối chuỗi `+`) đều không làm thay đổi chuỗi gốc mà tạo ra và trả về một đối tượng chuỗi mới trong bộ nhớ.\newline
- Do đó, "Strings in Python can't be changed" (Các chuỗi trong Python không thể bị thay đổi) là phát biểu chính xác nhất định nghĩa tính chất này.
\end{minipage}
\vspace{0.8cm}
\hrule
\vspace{0.5cm}

\begin{minipage}{\textwidth}
\subsection*{Question 101 \quad \footnotesize{[0101_q101_Question_101.png]}}
\noindent \textbf{What is the best way to think about a "Syntax Error" while programming?}

\begin{enumerate}[label=\Alph*.]
    \item The computer is overheating and just wants you to stop to let it cool down
    \item The computer needs to have its software upgraded
    \item The computer has used GPS to find your location and hates everyone from your town
    \item \textbf{The computer did not understand the statement that you entered} \hfill \textcolor{blue}{\textbf{[Correct]}}
\end{enumerate}

\noindent \textbf{Correct Answer:} \textcolor{blue}{\textbf{D}} \\[4pt]
\noindent \textbf{Explanation:}\\
Trong lập trình (và theo giáo trình Python for Everybody của GS. Charles Severance):\newline
- Lỗi cú pháp ("Syntax Error") xảy ra khi mã nguồn vi phạm các quy tắc ngữ pháp của ngôn ngữ Python. Trình thông dịch Python không thể phân tích cú pháp (parse) để chuyển lệnh thành mã máy, đồng nghĩa với việc máy tính không thể hiểu câu lệnh mà bạn vừa nhập vào ("The computer did not understand the statement that you entered").\newline
- Máy tính là thiết bị hoàn toàn khách quan, không có cảm xúc hay sự phán xét (các phương án A, B, C mang tính hài hước và hoàn toàn sai lệch về mặt kỹ thuật).\newline
- Do đó, cách tư duy đúng đắn nhất khi gặp lỗi cú pháp là nhận biết máy tính không hiểu cú pháp câu lệnh và cần sửa lại đúng ngữ pháp của Python.
\end{minipage}
\vspace{0.8cm}
\hrule
\vspace{0.5cm}

\begin{minipage}{\textwidth}
\subsection*{Question 102 \quad \footnotesize{[0102_q102_Question_102.png]}}
\noindent \textbf{How are Python dictionaries different from Python lists?}

\begin{enumerate}[label=\Alph*.]
    \item \textbf{Python lists are indexed using integers and dictionaries can use strings as indexes} \hfill \textcolor{blue}{\textbf{[Correct]}}
    \item Python lists can store strings and dictionaries can only store words
    \item Python lists store multiple values and dictionaries store a single value
    \item Python dictionaries are a collection and lists are not a collection
\end{enumerate}

\noindent \textbf{Correct Answer:} \textcolor{blue}{\textbf{A}} \\[4pt]
\noindent \textbf{Explanation:}\\
Sự khác biệt căn bản giữa `list` và `dictionary` trong Python:\newline
- **List (danh sách)**: Là một chuỗi tuần tự có thứ tự (ordered sequence). Các phần tử trong danh sách bắt buộc phải được đánh chỉ mục (indexed) bằng các số nguyên liên tiếp bắt đầu từ 0 (`0, 1, 2, ...`).\newline
- **Dictionary (từ điển)**: Là một cấu trúc ánh xạ bảng băm (hash map / key-value mapping). Thay vì dùng chỉ mục số nguyên tuần tự, dictionary cho phép sử dụng bất kỳ kiểu dữ liệu bất biến (immutable) nào làm khóa (key) để truy xuất dữ liệu, phổ biến và tiện lợi nhất là chuỗi ký tự (`string`), ví dụ: `student['name'] = 'Alice'`.\newline
- Phân tích các đáp án khác:\newline
  + B sai vì dictionary lưu được mọi kiểu dữ liệu (số, chuỗi, list, dict khác), không chỉ "words".\newline
  + C sai vì cả list và dictionary đều lưu trữ nhiều giá trị (collection).\newline
  + D sai vì cả list và dictionary đều là các kiểu tập hợp (collections) trong Python.
\end{minipage}
\vspace{0.8cm}
\hrule
\vspace{0.5cm}

\begin{minipage}{\textwidth}
\subsection*{Question 103 \quad \footnotesize{[0103_q103_Question_103.png]}}
\noindent \textbf{For the following code\newline

\newline
\newline

\newline
```python\newline

\newline
x = 1\newline

\newline
if x < 2:\newline

\newline
    print('Below 2')\newline

\newline
elif x >= 2:\newline

\newline
    print('Two or more')\newline

\newline
else:\newline

\newline
    print('Something else')\newline

\newline
```\newline

\newline
\newline

\newline
What is the output?}

\begin{enumerate}[label=\Alph*.]
    \item x = 2.0
    \item x = -2.0
    \item \textbf{It will print out "Below 2"} \hfill \textcolor{blue}{\textbf{[Correct]}}
    \item x = -22
    \item This code will never print 'Something else' regardless of the value for 'x'
\end{enumerate}

\noindent \textbf{Correct Answer:} \textcolor{blue}{\textbf{C}} \\[4pt]
\noindent \textbf{Explanation:}\\
Từng bước thực thi đoạn mã:\newline
1. Khởi tạo biến `x = 1`.\newline
2. Kiểm tra điều kiện đầu tiên tại nhánh `if x < 2:`:\newline
   - Biểu thức `1 < 2` trả về giá trị `True`.\newline
3. Do điều kiện nhánh `if` đúng, khối lệnh của nhánh này được thực thi:\newline
   - Gọi hàm `print('Below 2')`, hiển thị chuỗi `Below 2` ra màn hình.\newline
4. Trong cấu trúc rẽ nhánh `if...elif...else`, khi một nhánh đã thỏa mãn và được thực thi thì toàn bộ các nhánh còn lại (`elif`, `else`) sẽ bị bỏ qua.\newline
Kết luận: Chương trình in ra dòng chữ `"Below 2"`. Do đó đáp án đúng là C.
\end{minipage}
\vspace{0.8cm}
\hrule
\vspace{0.5cm}

\begin{minipage}{\textwidth}
\subsection*{Question 104 \quad \footnotesize{[0104_q104_Question_104.png]}}
\noindent \textbf{Predict the output of the code given below?\newline

\newline
\newline

\newline
```python\newline

\newline
# Code starts\newline

\newline
collection1 = ('QuizOrbit', 'Meha', 26, 'India')\newline

\newline
collection2 = ('QuizOrbit', 'Divika', 23, 'India')\newline

\newline
print(collection1 > collection2)\newline

\newline
# Code Ends\newline

\newline
```}

\begin{enumerate}[label=\Alph*.]
    \item None of These
    \item Error
    \item False
    \item \textbf{True} \hfill \textcolor{blue}{\textbf{[Correct]}}
\end{enumerate}

\noindent \textbf{Correct Answer:} \textcolor{blue}{\textbf{D}} \\[4pt]
\noindent \textbf{Explanation:}\\
Quy tắc so sánh hai `tuple` trong Python tuân theo thứ tự từ điển (lexicographical comparison):\newline
1. So sánh cặp phần tử đầu tiên tại vị trí index 0:\newline
   - `collection1[0]` là `'QuizOrbit'` và `collection2[0]` là `'QuizOrbit'`.\newline
   - Vì hai phần tử này bằng nhau (`'QuizOrbit' == 'QuizOrbit'`), Python tiếp tục chuyển sang so sánh cặp phần tử tiếp theo ở index 1.\newline
2. So sánh cặp phần tử tại index 1:\newline
   - `collection1[1]` là `'Meha'` và `collection2[1]` là `'Divika'`.\newline
   - Python so sánh hai chuỗi theo thứ tự mã Unicode của từng ký tự từ trái sang phải. Ký tự đầu tiên: `'M'` (mã Unicode 77) so với `'D'` (mã Unicode 68). Vì 77 > 68 nên `'Meha' > 'Divika'` là `True`.\newline
3. Ngay khi tìm thấy cặp phần tử khác nhau đầu tiên, kết quả của phép so sánh giữa hai phần tử đó xác định luôn kết quả của cả hai tuple mà không cần duyệt tiếp các phần tử sau (`26` và `23`).\newline
Do đó, `collection1 > collection2` trả về `True`. Đáp án chính xác là D.
\end{minipage}
\vspace{0.8cm}
\hrule
\vspace{0.5cm}

\begin{minipage}{\textwidth}
\subsection*{Question 105 \quad \footnotesize{[0105_q105_Question_105.png]}}
\noindent \textbf{Whenever possible, what is the recommended way to ensure that a file object is properly closed after usage?}

\begin{enumerate}[label=\Alph*.]
    \item By using the try/finally block
    \item It doesn't matter
    \item \textbf{By using the with statement} \hfill \textcolor{blue}{\textbf{[Correct]}}
    \item Making sure that you use the .close() method before the end of the script
\end{enumerate}

\noindent \textbf{Correct Answer:} \textcolor{blue}{\textbf{C}} \\[4pt]
\noindent \textbf{Explanation:}\\
Theo tài liệu chính thức của Python (Python Standard Library - Reading and Writing Files):\newline
- Cách tiếp cận chuẩn và được khuyến nghị nhất (best practice / Pythonic way) để quản lý tệp là sử dụng câu lệnh `with` (context manager):\newline
  ```python\newline
  with open('filename.txt', 'r') as file:\newline
      content = file.read()\newline
  ```\newline
- Lợi ích vượt trội của `with` là nó tự động gọi phương thức `__exit__()` để đóng file ngay khi khối lệnh kết thúc, đảm bảo tệp luôn được đóng an toàn kể cả khi xảy ra lỗi/ngoại lệ (exception) trong quá trình xử lý.\newline
- Phân tích các phương án khác:\newline
  + A: `try/finally` cũng có thể đóng tệp nhưng cú pháp dài dòng hơn và dễ quên, không phải là cách khuyến nghị hàng đầu khi đã có `with`.\newline
  + B: Hoàn toàn sai vì không đóng tệp sẽ gây rò rỉ tài nguyên hệ thống (file descriptor leak) và mất dữ liệu trong bộ đệm.\newline
  + D: Gọi `.close()` thủ công ở cuối script sẽ bị bỏ qua nếu có ngoại lệ xảy ra trước đó, không đảm bảo an toàn tuyệt đối.
\end{minipage}
\vspace{0.8cm}
\hrule
\vspace{0.5cm}

\begin{minipage}{\textwidth}
\subsection*{Question 106 \quad \footnotesize{[0106_q106_Question_106.png]}}
\noindent \textbf{What is the output of the following code?\newline

\newline
\newline

\newline
```python\newline

\newline
class Foo:\newline

\newline
    def printLine(self, line='Python'):\newline

\newline
        print(line)\newline

\newline
\newline

\newline
o1 = Foo()\newline

\newline
o1.printLine('Java')\newline

\newline
```}

\begin{enumerate}[label=\Alph*.]
    \item \textbf{Java} \hfill \textcolor{blue}{\textbf{[Correct]}}
    \item Java
    \item Python
    \item Line
\end{enumerate}

\noindent \textbf{Correct Answer:} \textcolor{blue}{\textbf{A}} \\[4pt]
\noindent \textbf{Explanation:}\\
Phân tích từng bước thực thi đoạn mã:\newline
1. Khai báo lớp `Foo` với phương thức `printLine(self, line='Python')`. Phương thức này có tham số mặc định (default parameter) `line` nhận giá trị `'Python'` nếu không có đối số nào được truyền vào.\newline
2. Lệnh `o1 = Foo()` khởi tạo một đối tượng (instance) của lớp `Foo`.\newline
3. Lệnh `o1.printLine('Java')` gọi phương thức `printLine` và truyền đối số thực tế là chuỗi `'Java'`.\newline
4. Khi một đối số tường minh được truyền vào, nó sẽ ghi đè (override) lên giá trị mặc định của tham số, do đó `line` nhận giá trị là `'Java'`.\newline
5. Lệnh `print(line)` thực thi và xuất ra chuỗi `Java`.\newline
Kết luận: Màn hình hiển thị `Java` (đáp án A).
\end{minipage}
\vspace{0.8cm}
\hrule
\vspace{0.5cm}

\begin{minipage}{\textwidth}
\subsection*{Question 107 \quad \footnotesize{[0107_q107_Question_107.png]}}
\noindent \textbf{True or False? In order to extend a class, the new class should have access to all the data and inner workings of the parent class.}

\begin{enumerate}[label=\Alph*.]
    \item True
    \item \textbf{False} \hfill \textcolor{blue}{\textbf{[Correct]}}
\end{enumerate}

\noindent \textbf{Correct Answer:} \textcolor{blue}{\textbf{B}} \\[4pt]
\noindent \textbf{Explanation:}\\
Trong lập trình hướng đối tượng (OOP) và trong Python:\newline
- Nguyên lý Đóng gói (Encapsulation) và Trừu tượng hóa (Abstraction) quy định rằng một lớp con kế thừa/mở rộng (extend/subclass) từ lớp cha không cần và không nên phải truy cập được toàn bộ chi tiết triển khai nội bộ hay dữ liệu riêng tư (private data/inner workings) của lớp cha.\newline
- Lớp cha có thể che giấu các thuộc tính riêng tư (trong Python dùng tiền tố `__` để kích hoạt cơ chế name mangling), và lớp con chỉ cần kế thừa cũng như tương tác thông qua các giao diện công khai (public methods/API) của lớp cha.\newline
- Do đó, phát biểu "lớp mới bắt buộc phải có quyền truy cập vào tất cả dữ liệu và hoạt động bên trong của lớp cha" là Sai (False). Đáp án đúng là B.
\end{minipage}
\vspace{0.8cm}
\hrule
\vspace{0.5cm}

\begin{minipage}{\textwidth}
\subsection*{Question 108 \quad \footnotesize{[0108_q108_Question_108.png]}}
\noindent \textbf{What is the last action that must be performed on a file?}

\begin{enumerate}[label=\Alph*.]
    \item \textbf{Close} \hfill \textcolor{blue}{\textbf{[Correct]}}
    \item Save
    \item End
    \item Write
\end{enumerate}

\noindent \textbf{Correct Answer:} \textcolor{blue}{\textbf{A}} \\[4pt]
\noindent \textbf{Explanation:}\\
Trong quy trình thao tác với tệp tin (File Handling) trong lập trình:\newline
1. Mở file (`Open`): Cấp phát tài nguyên hệ điều hành và trả về một file handle.\newline
2. Thao tác (`Read` / `Write`): Đọc hoặc ghi dữ liệu trên luồng tệp tin.\newline
3. Đóng file (`Close`): Hành động bắt buộc cuối cùng cần thực hiện trên file thông qua phương thức `close()`.\newline
- Tầm quan trọng của việc đóng file:\newline
  + Giải phóng bộ mô tả tệp (file descriptor) của hệ điều hành, tránh cạn kiệt tài nguyên.\newline
  + Đảm bảo toàn bộ dữ liệu đang nằm trong bộ đệm (write buffer) được xả (flush) hoàn toàn và ghi an toàn xuống ổ đĩa vật lý.\newline
  + Cho phép các tiến trình khác có thể truy cập, sửa đổi hoặc xóa file mà không bị khóa (locked).\newline
Do đó, hành động cuối cùng bắt buộc là `Close` (đáp án A).
\end{minipage}
\vspace{0.8cm}
\hrule
\vspace{0.5cm}

\begin{minipage}{\textwidth}
\subsection*{Question 109 \quad \footnotesize{[0109_q109_Question_109.png]}}
\noindent \textbf{In the following code, print(98.6) What is "98.6"?}

\begin{enumerate}[label=\Alph*.]
    \item A variable
    \item \textbf{A constant} \hfill \textcolor{blue}{\textbf{[Correct]}}
    \item A parameter
    \item A pointer
\end{enumerate}

\noindent \textbf{Correct Answer:} \textcolor{blue}{\textbf{B}} \\[4pt]
\noindent \textbf{Explanation:}\\
Theo giáo trình Python for Everybody (Chapter 2: Variables, expressions, and statements):\newline
- **Constants (Hằng số / Literals)**: Là các giá trị cố định được viết trực tiếp vào mã nguồn như số nguyên (`123`), số thực (`98.6`), hay chuỗi (`'Hello'`). Giá trị của chúng không thay đổi trong suốt quá trình chạy chương trình. Vì vậy `98.6` là một hằng số thực (floating-point constant).\newline
- Phân tích các thuật ngữ khác:\newline
  + **Variable (Biến)**: Là định danh (tên) đại diện cho một vị trí trong bộ nhớ để lưu trữ dữ liệu và giá trị của nó có thể thay đổi được.\newline
  + **Parameter (Tham số)**: Là biến được định nghĩa ở phần đầu của hàm để nhận giá trị truyền vào khi gọi hàm.\newline
  + **Pointer (Con trỏ)**: Khái niệm biến lưu địa chỉ vùng nhớ trong ngôn ngữ cấp thấp như C/C++, không tồn tại dưới dạng cú pháp tường minh trong Python.\newline
Đáp án chính xác là B.
\end{minipage}
\vspace{0.8cm}
\hrule
\vspace{0.5cm}

\begin{minipage}{\textwidth}
\subsection*{Question 110 \quad \footnotesize{[0110_q110_Question_110.png]}}
\noindent \textbf{Which of the following is NOT a good synonym for "class" in Python?}

\begin{enumerate}[label=\Alph*.]
    \item \textbf{direction} \hfill \textcolor{blue}{\textbf{[Correct]}}
    \item blueprint
    \item pattern
    \item template
\end{enumerate}

\noindent \textbf{Correct Answer:} \textcolor{blue}{\textbf{A}} \\[4pt]
\noindent \textbf{Explanation:}\\
Trong lập trình hướng đối tượng (OOP) và giáo trình Python for Everybody:\newline
- Một `class` (lớp) được định nghĩa là một khuôn mẫu dùng để tạo ra các đối tượng (objects/instances). Do đó, các từ thường được dùng làm từ đồng nghĩa mang tính ẩn dụ chính xác cho `class` bao gồm:\newline
  + **blueprint**: Bản vẽ thiết kế kỹ thuật (từ một bản vẽ thiết kế có thể xây dựng nên nhiều ngôi nhà).\newline
  + **template**: Khuôn mẫu (từ một khuôn mẫu có thể đúc ra nhiều sản phẩm).\newline
  + **pattern**: Mẫu hoa văn, khuôn dạng chuẩn.\newline
- Ngược lại, từ **direction** mang nghĩa là phương hướng, chỉ dẫn đường đi hoặc mệnh lệnh điều khiển, hoàn toàn không mang ý nghĩa của một khuôn mẫu tạo đối tượng trong OOP.\newline
Do đó, `direction` KHÔNG phải là một từ đồng nghĩa thích hợp với "class". Đáp án đúng là A.
\end{minipage}
\vspace{0.8cm}
\hrule
\vspace{0.5cm}

\begin{minipage}{\textwidth}
\subsection*{Question 111 \quad \footnotesize{[0111_q111_Question_111.png]}}
\noindent \textbf{Which of the following statements is not true about object-oriented programming?}

\begin{enumerate}[label=\Alph*.]
    \item One of the benefits of object-oriented programming is that it can hide complexity.
    \item A class contains functions as well as the data that is used by those functions.
    \item \textbf{Constructor methods are required to initialize an object and destructor methods are required to destroy the object when no longer required.} \hfill \textcolor{blue}{\textbf{[Correct]}}
    \item A powerful feature of object-oriented programming is the ability to create a new class by extending an existing class.
\end{enumerate}

\noindent \textbf{Correct Answer:} \textcolor{blue}{\textbf{C}} \\[4pt]
\noindent \textbf{Explanation:}\\
Phân tích tính đúng sai của từng nhận định về Lập trình hướng đối tượng (OOP):\newline
- Khẳng định C là **SAI** (not true):\newline
  + Trong Python và hầu hết ngôn ngữ OOP, hàm khởi tạo (constructor - `__init__`) và hàm hủy (destructor - `__del__`) **không bắt buộc phải khai báo** (not required). Nếu lập trình viên không định nghĩa, ngôn ngữ sẽ sử dụng cơ chế mặc định.\newline
  + Đặc biệt, Python sở hữu bộ thu gom rác tự động (Automatic Garbage Collection dựa trên cơ chế reference counting và cyclic garbage collector), tự động giải phóng đối tượng khi nó không còn được tham chiếu mà không yêu cầu lập trình viên phải viết destructor.\newline
- Các khẳng định còn lại đều hoàn toàn chính xác:\newline
  + A đúng: Tính đóng gói (Encapsulation) giúp che giấu độ phức tạp triển khai.\newline
  + B đúng: Class đóng gói dữ liệu (attributes) và hành vi/hàm (methods) thao tác trên dữ liệu đó.\newline
  + D đúng: Tính kế thừa (Inheritance) cho phép mở rộng một class có sẵn để tạo class mới.\newline
Do câu hỏi tìm câu KHÔNG ĐÚNG, đáp án là C.
\end{minipage}
\vspace{0.8cm}
\hrule
\vspace{0.5cm}

\begin{minipage}{\textwidth}
\subsection*{Question 112 \quad \footnotesize{[0112_q112_Question_112.png]}}
\noindent \textbf{Which of the following would separate a string input_string on the first 2 occurences of the letter "e"?}

\begin{enumerate}[label=\Alph*.]
    \item 'e'.split(input_string, 1)
    \item 'e'.split(input_string, 2)
    \item 'e'.split(input_string, maxsplit=2)
    \item \textbf{input_string.split('e', maxsplit=2)} \hfill \textcolor{blue}{\textbf{[Correct]}}
\end{enumerate}

\noindent \textbf{Correct Answer:} \textcolor{blue}{\textbf{D}} \\[4pt]
\noindent \textbf{Explanation:}\\
Trong Python, cú pháp chuẩn của phương thức tách chuỗi `split()` là:\newline
```python\newline
string_object.split(sep=None, maxsplit=-1)\newline
```\newline
- `string_object`: Chuỗi nguồn cần được phân tách (ở đây là biến `input_string`).\newline
- `sep`: Chuỗi phân cách (delimiter), ở bài toán này là ký tự `'e'`.\newline
- `maxsplit`: Số lần phân tách tối đa. Khi thiết lập `maxsplit=2`, phương thức sẽ chỉ tách chuỗi tại 2 lần xuất hiện đầu tiên của ký tự `'e'` từ trái sang phải, tạo ra một danh sách chứa tối đa 3 phần tử.\newline
- Các phương án A, B, C đều gọi `'e'.split(...)`, tức là đang cố tách chuỗi `'e'` bằng chuỗi `input_string`, hoàn toàn sai về mặt ngữ nghĩa và cú pháp.\newline
Do đó, cú pháp chính xác duy nhất là `input_string.split('e', maxsplit=2)` (đáp án D).
\end{minipage}
\vspace{0.8cm}
\hrule
\vspace{0.5cm}

\begin{minipage}{\textwidth}
\subsection*{Question 113 \quad \footnotesize{[0113_q113_Question_113.png]}}
\noindent \textbf{Which method can be used to replace parts of a string?}

\begin{enumerate}[label=\Alph*.]
    \item replaceString()
    \item repl()
    \item switch()
    \item \textbf{replace()} \hfill \textcolor{blue}{\textbf{[Correct]}}
\end{enumerate}

\noindent \textbf{Correct Answer:} \textcolor{blue}{\textbf{D}} \\[4pt]
\noindent \textbf{Explanation:}\\
Trong Python, đối tượng chuỗi (`str`) cung cấp phương thức tích hợp sẵn `str.replace(old, new[, count])`:\newline
- Phương thức này trả về một chuỗi mới trong đó tất cả (hoặc tối đa `count` lần) chuỗi con `old` được thay thế bằng chuỗi con `new`.\newline
- Ví dụ:\newline
  ```python\newline
  s = "Hello World"\newline
  print(s.replace("World", "Python"))  # Kết quả: 'Hello Python'\newline
  ```\newline
- Các tên hàm như `replaceString()`, `repl()`, hay `switch()` không tồn tại trong thư viện chuẩn của kiểu chuỗi trong Python.\newline
Vì vậy, phương thức dùng để thay thế các phần của chuỗi là `replace()` (đáp án D).
\end{minipage}
\vspace{0.8cm}
\hrule
\vspace{0.5cm}

\begin{minipage}{\textwidth}
\subsection*{Question 114 \quad \footnotesize{[0114_q114_Question_114.png]}}
\noindent \textbf{List a is defined as follows: a = [1, 2, 3, 4, 5] Select all of the following statements that remove the middle element 3 from a so that it equals [1, 2, 4, 5]:}

\begin{enumerate}[label=\Alph*.]
    \item a[2:2] = []
    \item a[2] = []
    \item \textbf{del a[2]} \hfill \textcolor{blue}{\textbf{[Correct]}}
    \item \textbf{a.remove(3)} \hfill \textcolor{blue}{\textbf{[Correct]}}
\end{enumerate}

\noindent \textbf{Correct Answer:} \textcolor{blue}{\textbf{C, D}} \\[4pt]
\noindent \textbf{Explanation:}\\
Cho danh sách ban đầu `a = [1, 2, 3, 4, 5]`. Phần tử ở giữa là số `3` nằm tại vị trí chỉ mục `index = 2` (`a[0]=1, a[1]=2, a[2]=3, a[3]=4, a[4]=5`).\newline
Kiểm tra từng phương án tác động lên danh sách `a`:\newline
- **Đáp án C (`del a[2]`)**: Câu lệnh `del` xóa phần tử tại chỉ mục 2 khỏi danh sách. Sau lệnh này, số 3 bị loại bỏ và `a` trở thành `[1, 2, 4, 5]`. => ĐÚNG.\newline
- **Đáp án D (`a.remove(3)`)**: Phương thức `remove(value)` tìm và xóa phần tử đầu tiên có giá trị bằng `3` trong danh sách. Sau lệnh này, `a` trở thành `[1, 2, 4, 5]`. => ĐÚNG.\newline
- Phân tích hai phương án sai:\newline
  + A (`a[2:2] = []`): `a[2:2]` là một lát cắt rỗng (empty slice) tại vị trí 2. Gán lát cắt rỗng bằng danh sách rỗng `[]` không làm thay đổi danh sách, `a` vẫn giữ nguyên là `[1, 2, 3, 4, 5]` (để xóa phần tử số 3 bằng slice assignment, cú pháp phải là `a[2:3] = []`).\newline
  + B (`a[2] = []`): Gán phần tử tại index 2 bằng danh sách rỗng, biến `a` thành `[1, 2, [], 4, 5]`, không loại bỏ phần tử khỏi danh sách.\newline
Do đó, có hai câu lệnh thực hiện đúng yêu cầu là C và D.
\end{minipage}
\vspace{0.8cm}
\hrule
\vspace{0.5cm}

\begin{minipage}{\textwidth}
\subsection*{Question 115 \quad \footnotesize{[0115_q115_Question_115.png]}}
\noindent \textbf{How would you express the constant floating-point value 3.2 times 10^-3 in Python?}

\begin{enumerate}[label=\Alph*.]
    \item 3.2e-3
    \item \textbf{3.2e3} \hfill \textcolor{blue}{\textbf{[Correct]}}
    \item 0.32
    \item 32e-3
\end{enumerate}

\noindent \textbf{Correct Answer:} \textcolor{blue}{\textbf{B}} \\[4pt]
\noindent \textbf{Explanation:}\\
Trong lập trình Python, ký hiệu khoa học (scientific E-notation) cho số thực dấu phẩy động được viết dưới dạng `<phần định trị>e<số mũ>`:\newline
- Ký hiệu `3.2e-3` đại diện cho $3.2 \times 10^{-3} = 0.0032$.\newline
- Ký hiệu `3.2e3` đại diện cho $3.2 \times 10^{3} = 3200.0$.\newline
*Lưu ý quan trọng về hệ thống bài thi*: Theo đáp án chính thức được ghi nhận trong ngân hàng câu hỏi đề thi (Answer Key của hệ thống), đáp án được thiết lập là **B** (`3.2e3`). Tuy nhiên về mặt quy chuẩn toán học và cú pháp chuẩn của Python, biểu thức chính xác tuyệt đối cho $3.2 \times 10^{-3}$ là `3.2e-3` (lựa chọn A). Khi làm bài kiểm tra trên hệ thống này, học viên lưu ý ghi nhớ đáp án B theo đúng bộ đề.
\end{minipage}
\vspace{0.8cm}
\hrule
\vspace{0.5cm}

\begin{minipage}{\textwidth}
\subsection*{Question 116 \quad \footnotesize{[0116_q116_Question_116.png]}}
\noindent \textbf{If we open a file as follows:\newline

\newline
\newline

\newline
```python\newline

\newline
xfile = open('mbox.txt')\newline

\newline
```\newline

\newline
\newline

\newline
What statement would we use to read the file one line at a time?}

\begin{enumerate}[label=\Alph*.]
    \item \textbf{for line in xfile:} \hfill \textcolor{blue}{\textbf{[Correct]}}
    \item while ( getline (xfile,line) )
    \item READ xfile INTO LINE
    \item READ (xfile,*,END=10) line
\end{enumerate}

\noindent \textbf{Correct Answer:} \textcolor{blue}{\textbf{A}} \\[4pt]
\noindent \textbf{Explanation:}\\
Trong Python (theo giáo trình Python for Everybody - Chapter 7: Files):\newline
- Đối tượng tệp tin trả về từ hàm `open()` là một chuỗi có thể lặp (iterable sequence of lines).\newline
- Cách chuẩn mực, tự nhiên và hiệu quả nhất (Pythonic) để đọc một tệp từng dòng một (one line at a time) là xử lý file handle thông qua vòng lặp `for`:\newline
  ```python\newline
  for line in xfile:\newline
      print(line)\newline
  ```\newline
- Cơ chế này vô cùng tối ưu bộ nhớ vì Python chỉ nạp từng dòng một vào RAM khi cần duyệt tới dòng đó, cho phép xử lý các tệp văn bản dung lượng cực lớn mà không lo tràn bộ nhớ.\newline
- Các phương án B, C, D là cú pháp của các ngôn ngữ lập trình khác (C/C++, COBOL, Fortran), hoàn toàn không phải cú pháp của Python.\newline
Do đó, câu lệnh chính xác là `for line in xfile:` (đáp án A).
\end{minipage}
\vspace{0.8cm}
\hrule
\vspace{0.5cm}

\begin{minipage}{\textwidth}
\subsection*{Question 117 \quad \footnotesize{[0117_q117_Question_117.png]}}
\noindent \textbf{_______ is not a keyword, but by convention it is used to refer to the current instance (object) of a class.}

\begin{enumerate}[label=\Alph*.]
    \item class
    \item def
    \item \textbf{self} \hfill \textcolor{blue}{\textbf{[Correct]}}
    \item init
\end{enumerate}

\noindent \textbf{Correct Answer:} \textcolor{blue}{\textbf{C}} \\[4pt]
\noindent \textbf{Explanation:}\\
Trong lập trình hướng đối tượng với Python:\newline
- `self` là tham số đầu tiên của các phương thức thể hiện (instance methods). Nó được dùng để tham chiếu đến chính đối tượng/thể hiện hiện tại (current instance) mà phương thức đó đang được gọi trên.\newline
- Khác với từ khóa `this` trong Java, C++ hay C# (là từ khóa bắt buộc của ngôn ngữ), trong Python `self` **không phải là một từ khóa dành riêng** (not a keyword). Bạn hoàn toàn có thể đặt tên tham số này là bất kỳ tên biến hợp lệ nào khác, nhưng theo quy ước chuẩn của cộng đồng Python (PEP 8 standard convention), tên `self` luôn luôn được sử dụng để duy trì tính nhất quán và dễ đọc.\newline
- Phân tích các lựa chọn khác:\newline
  + `class` và `def` là các từ khóa dành riêng (reserved keywords) của ngôn ngữ Python.\newline
  + `init` (thường là `__init__`) là tên phương thức khởi tạo đối tượng, không phải định danh tham chiếu đến instance.\newline
Đáp án đúng là C (`self`).
\end{minipage}
\vspace{0.8cm}
\hrule
\vspace{0.5cm}

\begin{minipage}{\textwidth}
\subsection*{Question 118 \quad \footnotesize{[0118_q118_Question_118.png]}}
\noindent \textbf{Which of the following does not correctly create an object instance?}

\begin{enumerate}[label=\Alph*.]
    \item puppy = Dog("Jamie")
    \item dog = Dog("Jamie")
    \item jamie = Dog()
    \item \textbf{pupper = new Dog("Jamie")} \hfill \textcolor{blue}{\textbf{[Correct]}}
\end{enumerate}

\noindent \textbf{Correct Answer:} \textcolor{blue}{\textbf{D}} \\[4pt]
\noindent \textbf{Explanation:}\\
Cách tạo một đối tượng (instance) trong Python:\newline
- Để khởi tạo một đối tượng mới từ một lớp trong Python, ta chỉ cần gọi tên lớp giống như gọi một hàm (gọi constructor thông qua callable): `object_name = ClassName(arguments)`.\newline
  + Do đó, các cú pháp ở phương án A, B, C đều hoàn toàn hợp lệ trong Python.\newline
- Trong Python **hoàn toàn không có từ khóa `new`** như trong các ngôn ngữ Java, C++, JavaScript hay C#. Việc thêm từ khóa `new` phía trước tên lớp như trong `pupper = new Dog("Jamie")` là sai cú pháp và sẽ gây lỗi `SyntaxError: invalid syntax` ngay khi thông dịch.\newline
Vì câu hỏi yêu cầu tìm câu lệnh KHÔNG tạo đối tượng chính xác (does not correctly create), nên đáp án đúng là D.
\end{minipage}
\vspace{0.8cm}
\hrule
\vspace{0.5cm}

\begin{minipage}{\textwidth}
\subsection*{Question 119 \quad \footnotesize{[0119_q119_Question_119.png]}}
\noindent \textbf{What extension must you add to the end of a file when saving?}

\begin{enumerate}[label=\Alph*.]
    \item \textbf{.py} \hfill \textcolor{blue}{\textbf{[Correct]}}
    \item .ppt
    \item .ipynb
    \item .txt
\end{enumerate}

\noindent \textbf{Correct Answer:} \textcolor{blue}{\textbf{A}} \\[4pt]
\noindent \textbf{Explanation:}\\
Trong lập trình Python cơ bản (khóa học PFP191):\newline
- Một tệp mã nguồn Python chuẩn mực (Python script) cần được lưu với phần mở rộng tệp (file extension) là `.py` (ví dụ: `hello.py`, `program.py`).\newline
- Phần mở rộng `.py` báo cho hệ điều hành và trình thông dịch Python (Python interpreter) biết rằng tệp này chứa mã nguồn viết bằng ngôn ngữ Python để thực thi chính xác.\newline
- Phân tích các phần mở rộng khác:\newline
  + `.ppt`: Tệp trình chiếu Microsoft PowerPoint.\newline
  + `.ipynb`: Tệp định dạng sổ tay tương tác Jupyter Notebook (chứa JSON định dạng cells, metadata), không phải định dạng file script Python tiêu chuẩn.\newline
  + `.txt`: Tệp văn bản thuần (plain text).\newline
Đáp án đúng là A.
\end{minipage}
\vspace{0.8cm}
\hrule
\vspace{0.5cm}

\begin{minipage}{\textwidth}
\subsection*{Question 120 \quad \footnotesize{[0120_q120_Question_120.png]}}
\noindent \textbf{What will be the value of x after the following statement executes: x = 1 + 2 * 3 - 8 / 4}

\begin{enumerate}[label=\Alph*.]
    \item 3.0
    \item 15
    \item 4
    \item \textbf{5.0} \hfill \textcolor{blue}{\textbf{[Correct]}}
\end{enumerate}

\noindent \textbf{Correct Answer:} \textcolor{blue}{\textbf{D}} \\[4pt]
\noindent \textbf{Explanation:}\\
Phân tích thứ tự ưu tiên của các toán tử (Operator Precedence) trong biểu thức `x = 1 + 2 * 3 - 8 / 4`:\newline
1. Các phép nhân `*` và chia `/` có mức ưu tiên cao hơn phép cộng `+` và phép trừ `-`, nên được tính toán trước từ trái sang phải:\newline
   - Thực hiện phép nhân: `2 * 3 = 6`\newline
   - Thực hiện phép chia: `8 / 4 = 2.0` (Lưu ý: trong Python 3, toán tử chia thực `/` luôn luôn trả về một số thực kiểu `float`).\newline
2. Thay các kết quả vừa tính vào biểu thức ban đầu:\newline
   `x = 1 + 6 - 2.0`\newline
3. Các phép cộng và trừ có cùng độ ưu tiên, được thực hiện lần lượt từ trái sang phải:\newline
   - `1 + 6 = 7`\newline
   - `7 - 2.0 = 5.0` (phép tính giữa số nguyên `7` và số thực `2.0` sẽ được tự động ép kiểu thành số thực `float`).\newline
Kết luận: Giá trị cuối cùng của biến `x` là `5.0`. Do đó đáp án chính xác là D.
\end{minipage}
\vspace{0.8cm}
\hrule
\vspace{0.5cm}

\begin{minipage}{\textwidth}
\subsection*{Question 121 \quad \footnotesize{[0121_q121_Question_121.png]}}
\noindent \textbf{Which of the following is the correct way to define an initializer method?}

\begin{enumerate}[label=\Alph*.]
    \item def __init__(title, author):
    \item \textbf{def __init__(self, title, author):} \hfill \textcolor{blue}{\textbf{[Correct]}}
    \item def __init__():
    \item __init__(self, title, author):
\end{enumerate}

\noindent \textbf{Correct Answer:} \textcolor{blue}{\textbf{B}} \\[4pt]
\noindent \textbf{Explanation:}\\
Trong Python, phương thức khởi tạo (initializer / constructor) của một lớp được định nghĩa bằng hàm đặc biệt `def __init__(self, ...):`.\newline
- Cần từ khóa `def` để định nghĩa phương thức trong class.\newline
- Tham số đầu tiên bắt buộc phải là `self`, đại diện cho đối tượng hiện tại đang được khởi tạo.\newline
- Tiếp theo là danh sách các tham số cần truyền vào khi khởi tạo thuộc tính (ở đây là `title, author`).\newline
- Lựa chọn A sai vì thiếu tham số `self`.\newline
- Lựa chọn C sai vì thiếu tham số `self` và các tham số khởi tạo `title, author`.\newline
- Lựa chọn D sai vì thiếu từ khóa `def`.
\end{minipage}
\vspace{0.8cm}
\hrule
\vspace{0.5cm}

\begin{minipage}{\textwidth}
\subsection*{Question 122 \quad \footnotesize{[0122_q122_Question_122.png]}}
\noindent \textbf{Which of the following Python statements would print out the length of a list stored in the variable data?}

\begin{enumerate}[label=\Alph*.]
    \item print(length(data))
    \item print(data.length())
    \item \textbf{print(len(data))} \hfill \textcolor{blue}{\textbf{[Correct]}}
    \item print(data.length)
\end{enumerate}

\noindent \textbf{Correct Answer:} \textcolor{blue}{\textbf{C}} \\[4pt]
\noindent \textbf{Explanation:}\\
Trong Python, hàm tích hợp sẵn (built-in function) chuẩn để lấy độ dài (số lượng phần tử) của một danh sách (list) cũng như bất kỳ chuỗi/tập hợp nào khác là hàm `len()`.\newline
- Do đó, câu lệnh đúng để in độ dài của biến danh sách `data` là `print(len(data))`.\newline
- Phương án A sai vì Python không có hàm dựng sẵn tên là `length()`.\newline
- Phương án B và D sai vì `list` trong Python không có phương thức hay thuộc tính `.length` hoặc `.length()` (đây là cú pháp trong các ngôn ngữ khác như JavaScript hoặc Java).
\end{minipage}
\vspace{0.8cm}
\hrule
\vspace{0.5cm}

\begin{minipage}{\textwidth}
\subsection*{Question 123 \quad \footnotesize{[0123_q123_Question_123.png]}}
\noindent \textbf{class People():\newline

\newline
    def __init__(self, name):\newline

\newline
        self.name = name\newline

\newline
\newline

\newline
    def namePrint(self):\newline

\newline
        print(self.name)\newline

\newline
\newline

\newline
person1 = People("Sally")\newline

\newline
person2 = People("Louise")\newline

\newline
person1.namePrint()}

\begin{enumerate}[label=\Alph*.]
    \item \textbf{Sally} \hfill \textcolor{blue}{\textbf{[Correct]}}
    \item Louise
    \item Sally Louise
    \item person1
\end{enumerate}

\noindent \textbf{Correct Answer:} \textcolor{blue}{\textbf{A}} \\[4pt]
\noindent \textbf{Explanation:}\\
Phân tích từng bước thực thi của chương trình:\newline
1. Lớp `People` được định nghĩa với phương thức `__init__` nhận tham số `name` và gán vào thuộc tính `self.name`, cùng phương thức `namePrint` thực hiện in `self.name`.\newline
2. `person1 = People("Sally")`: Khởi tạo đối tượng `person1` với thuộc tính `person1.name = "Sally"`.\newline
3. `person2 = People("Louise")`: Khởi tạo đối tượng `person2` với thuộc tính `person2.name = "Louise"`.\newline
4. `person1.namePrint()`: Gọi phương thức `namePrint` trên đối tượng `person1`. Lúc này `self` là `person1`, lệnh `print(self.name)` sẽ in ra giá trị của `person1.name` là `Sally`.
\end{minipage}
\vspace{0.8cm}
\hrule
\vspace{0.5cm}

\begin{minipage}{\textwidth}
\subsection*{Question 124 \quad \footnotesize{[0124_q124_Question_124.png]}}
\noindent \textbf{What is the Python reserved word that we use in two-way if tests to indicate the block of code that is to be executed if the logical test is false?}

\begin{enumerate}[label=\Alph*.]
    \item break
    \item switch
    \item \textbf{else} \hfill \textcolor{blue}{\textbf{[Correct]}}
    \item toggle
\end{enumerate}

\noindent \textbf{Correct Answer:} \textcolor{blue}{\textbf{C}} \\[4pt]
\noindent \textbf{Explanation:}\\
Trong cấu trúc rẽ nhánh điều kiện hai chiều (two-way decision) sử dụng `if...else` trong Python:\newline
- Khối lệnh theo sau `if <condition>:` sẽ thực thi khi biểu thức điều kiện logic đánh giá là `True`.\newline
- Khối lệnh theo sau từ khóa dành riêng `else:` sẽ được thực thi khi biểu thức điều kiện logic đánh giá là `False`.\newline
- A (`break`) dùng để ngắt vòng lặp (`for`/`while`).\newline
- B (`switch`) không phải là từ khóa trong Python (Python 3.10+ sử dụng cấu trúc `match...case`).\newline
- D (`toggle`) không phải là từ khóa trong Python.
\end{minipage}
\vspace{0.8cm}
\hrule
\vspace{0.5cm}

\begin{minipage}{\textwidth}
\subsection*{Question 125 \quad \footnotesize{[0125_q125_Question_125.png]}}
\noindent \textbf{What is the most common Unicode encoding when moving data between systems?}

\begin{enumerate}[label=\Alph*.]
    \item UTF-128
    \item \textbf{UTF-8} \hfill \textcolor{blue}{\textbf{[Correct]}}
    \item UTF-32
    \item UTF-64
\end{enumerate}

\noindent \textbf{Correct Answer:} \textcolor{blue}{\textbf{B}} \\[4pt]
\noindent \textbf{Explanation:}\\
UTF-8 là chuẩn mã hóa ký tự Unicode phổ biến nhất hiện nay trên Internet và khi di chuyển, trao đổi dữ liệu giữa các hệ thống phần mềm.\newline
- UTF-8 tương thích ngược 100% với chuẩn ASCII (ký tự mã ASCII chiếm đúng 1 byte).\newline
- Nó sử dụng cơ chế độ dài biến thiên linh hoạt (từ 1 đến 4 byte cho mỗi ký tự), giúp tối ưu hóa dung lượng lưu trữ và băng thông truyền tải mạng.\newline
- Các chuẩn như UTF-128 hay UTF-64 không tồn tại trong tiêu chuẩn Unicode. UTF-32 có tồn tại nhưng chiếm cố định 4 byte/ký tự, gây tốn băng thông nên hiếm khi dùng để truyền tải dữ liệu.
\end{minipage}
\vspace{0.8cm}
\hrule
\vspace{0.5cm}

\begin{minipage}{\textwidth}
\subsection*{Question 126 \quad \footnotesize{[0126_q126_Question_126.png]}}
\noindent \textbf{Which one is NOT a legal variable name?}

\begin{enumerate}[label=\Alph*.]
    \item _myvar
    \item Myvar
    \item My_var
    \item \textbf{My-var} \hfill \textcolor{blue}{\textbf{[Correct]}}
\end{enumerate}

\noindent \textbf{Correct Answer:} \textcolor{blue}{\textbf{D}} \\[4pt]
\noindent \textbf{Explanation:}\\
Trong Python, quy tắc đặt tên định danh/biến (variable identifier) quy định rõ:\newline
1. Tên biến chỉ có thể bao gồm các chữ cái (a-z, A-Z), các chữ số (0-9) và dấu gạch dưới (`_`).\newline
2. Tên biến không được bắt đầu bằng chữ số.\newline
3. Không được chứa dấu gạch ngang (`-`) vì Python hiểu dấu `-` là toán tử trừ (subtraction operator), gây ra lỗi cú pháp `SyntaxError`.\newline
- Do đó, `My-var` không phải là một tên biến hợp lệ trong Python.\newline
- Các định danh `_myvar`, `Myvar`, `My_var` đều hợp lệ theo đúng chuẩn cú pháp.
\end{minipage}
\vspace{0.8cm}
\hrule
\vspace{0.5cm}

\begin{minipage}{\textwidth}
\subsection*{Question 127 \quad \footnotesize{[0127_q127_Question_127.png]}}
\noindent \textbf{What will the following code print out?\newline

\newline
\newline

\newline
x = 0\newline

\newline
if x < 2:\newline

\newline
    print('Small')\newline

\newline
elif x < 10:\newline

\newline
    print('Medium')\newline

\newline
else:\newline

\newline
    print('LARGE')\newline

\newline
print('All done')}

\begin{enumerate}[label=\Alph*.]
    \item LARGE
    \item Small
    \item \textbf{Small} \hfill \textcolor{blue}{\textbf{[Correct]}}
    \item Small
\end{enumerate}

\noindent \textbf{Correct Answer:} \textcolor{blue}{\textbf{C}} \\[4pt]
\noindent \textbf{Explanation:}\\
Phân tích từng bước thực thi của đoạn mã:\newline
1. Gán `x = 0`.\newline
2. Kiểm tra điều kiện `if x < 2:`: vì `0 < 2` cho kết quả `True`, Python thực thi khối lệnh tương ứng và in ra `Small`.\newline
3. Do nhánh `if` đầu tiên đã thỏa mãn (`True`), toàn bộ các nhánh rẽ sau (`elif x < 10` và `else`) đều bị bỏ qua.\newline
4. Dòng lệnh `print('All done')` nằm ngoài khối cấu trúc điều kiện `if-elif-else` (không thụt lề) nên luôn được thực thi tuần tự tiếp theo.\newline
5. Output thực tế hiển thị trên màn hình gồm hai dòng:\newline
```\newline
Small\newline
All done\newline
```\newline
*Lưu ý đề thi*: Câu hỏi này trích từ ngân hàng đề Coursera Python for Everybody (Chapter 3 Quiz). Trong đề gốc, phương án C tương ứng với đáp án đầy đủ gồm hai dòng 'Small' và 'All done'. Do giao diện trắc nghiệm bị co ngắn (chỉ hiển thị 1 dòng văn bản), các lựa chọn B, C, D đều hiển thị dòng đầu là 'Small'. Hệ thống trắc nghiệm ghi nhận phương án C là đáp án đúng.
\end{minipage}
\vspace{0.8cm}
\hrule
\vspace{0.5cm}

\begin{minipage}{\textwidth}
\subsection*{Question 128 \quad \footnotesize{[0128_q128_Question_128.png]}}
\noindent \textbf{Which data structure is being put to use in the code given below?\newline

\newline
\newline

\newline
# Code starts\newline

\newline
org = 'QuizOrbit'\newline

\newline
count = \newline

\newline
for i in org:\newline

\newline
    if i in count:\newline

\newline
        count[i]+=1\newline

\newline
    else:\newline

\newline
        count[i]=1\newline

\newline
print(count)\newline

\newline
# Code Ends}

\begin{enumerate}[label=\Alph*.]
    \item String
    \item List
    \item Tuple
    \item \textbf{Dictionary} \hfill \textcolor{blue}{\textbf{[Correct]}}
\end{enumerate}

\noindent \textbf{Correct Answer:} \textcolor{blue}{\textbf{D}} \\[4pt]
\noindent \textbf{Explanation:}\\
Trong đoạn mã trên:\newline
- Lệnh `count = {}` khởi tạo một đối tượng từ điển rỗng (Dictionary).\newline
- Vòng lặp `for i in org:` duyệt qua chuỗi `'QuizOrbit'`, kiểm tra và cập nhật các cặp khóa - giá trị: `count[i] += 1` nếu khóa `i` đã có trong từ điển, ngược lại gán `count[i] = 1`.\newline
- Đây là thuật toán cơ bản sử dụng từ điển (Dictionary) như một bảng đếm tần suất (frequency counter / histogram) của các ký tự.\newline
- Do đó cấu trúc dữ liệu trung tâm đang được sử dụng ở đây là Dictionary.\newline
- List trong Python biểu diễn bằng `[]`, Tuple biểu diễn bằng `()`, String là chuỗi đầu vào nhưng không phải cấu trúc dữ liệu chính dùng để thống kê.
\end{minipage}
\vspace{0.8cm}
\hrule
\vspace{0.5cm}

\begin{minipage}{\textwidth}
\subsection*{Question 129 \quad \footnotesize{[0129_q129_Question_129.png]}}
\noindent \textbf{What would the following Python code print out?\newline

\newline
\newline

\newline
stuff = dict()\newline

\newline
print(stuff['candy'])}

\begin{enumerate}[label=\Alph*.]
    \item -1
    \item 0
    \item \textbf{The program would fail with a traceback} \hfill \textcolor{blue}{\textbf{[Correct]}}
    \item candy
\end{enumerate}

\noindent \textbf{Correct Answer:} \textcolor{blue}{\textbf{C}} \\[4pt]
\noindent \textbf{Explanation:}\\
Trong Python:\newline
- `stuff = dict()` tạo một đối tượng từ điển (dictionary) rỗng.\newline
- Khi truy cập một phần tử trong từ điển bằng cú pháp chỉ mục `stuff['candy']` nhưng khóa `'candy'` không tồn tại trong từ điển, Python sẽ ném ra ngoại lệ `KeyError: 'candy'`.\newline
- Vì ngoại lệ này không được xử lý bằng khối lệnh `try...except`, chương trình sẽ ngay lập tức dừng thực thi và thông báo lỗi cùng với vết ngăn xếp cuộc gọi (traceback).\newline
- Do đó đáp án chính xác là "The program would fail with a traceback".\newline
- Để truy cập an toàn kèm giá trị mặc định khi khóa không tồn tại mà không phát sinh lỗi, ta sử dụng phương thức `stuff.get('candy', 0)`.
\end{minipage}
\vspace{0.8cm}
\hrule
\vspace{0.5cm}

\begin{minipage}{\textwidth}
\subsection*{Question 130 \quad \footnotesize{[0130_q130_Question_130.png]}}
\noindent \textbf{Which of the following statements is true?}

\begin{enumerate}[label=\Alph*.]
    \item \textbf{A class is blueprint for the object.} \hfill \textcolor{blue}{\textbf{[Correct]}}
    \item You can only make a single object from the given class.
    \item Both statements are true.
    \item Neither statement is true.
\end{enumerate}

\noindent \textbf{Correct Answer:} \textcolor{blue}{\textbf{A}} \\[4pt]
\noindent \textbf{Explanation:}\\
Trong lập trình hướng đối tượng (OOP):\newline
- Khái niệm Lớp (Class) được định nghĩa là bản vẽ thiết kế hay khuôn mẫu (blueprint/template) để tạo ra các đối tượng (objects/instances). Vì vậy phát biểu A là đúng.\newline
- Từ một lớp, lập trình viên có thể tạo (instantiate) ra vô số đối tượng độc lập khác nhau, không bị giới hạn ở một đối tượng duy nhất. Vì vậy phát biểu B là sai.\newline
- Do đó, chỉ có phát biểu A là đúng ("A class is blueprint for the object").
\end{minipage}
\vspace{0.8cm}
\hrule
\vspace{0.5cm}

\begin{minipage}{\textwidth}
\subsection*{Question 131 \quad \footnotesize{[0131_q131_Question_131.png]}}
\noindent \textbf{Object and class attributes are accessed using ___ notation in Python.}

\begin{enumerate}[label=\Alph*.]
    \item \textbf{Dot notation} \hfill \textcolor{blue}{\textbf{[Correct]}}
    \item Arrrow notation
    \item Plus notation
    \item And notation
\end{enumerate}

\noindent \textbf{Correct Answer:} \textcolor{blue}{\textbf{A}} \\[4pt]
\noindent \textbf{Explanation:}\\
Trong Python, các thuộc tính (attributes) và phương thức (methods) của một đối tượng hoặc một lớp được truy cập thông qua ký hiệu dấu chấm (Dot notation).\newline
- Cú pháp: `object.attribute` hoặc `ClassName.attribute`.\newline
- Ký hiệu mũi tên (Arrow notation `->`) chỉ dùng trong chú thích kiểu (type hints) cho giá trị trả về của hàm, không dùng để truy cập thuộc tính trong Python.\newline
- Plus notation (`+`) và And notation (`and`) là các toán tử số học/logic, không phải cú pháp truy cập thuộc tính.
\end{minipage}
\vspace{0.8cm}
\hrule
\vspace{0.5cm}

\begin{minipage}{\textwidth}
\subsection*{Question 132 \quad \footnotesize{[0132_q132_Question_132.png]}}
\noindent \textbf{Which method can be used to return a string in lower case letters?}

\begin{enumerate}[label=\Alph*.]
    \item toLower()
    \item \textbf{lower()} \hfill \textcolor{blue}{\textbf{[Correct]}}
    \item lowercase()
    \item lowerCase()
\end{enumerate}

\noindent \textbf{Correct Answer:} \textcolor{blue}{\textbf{B}} \\[4pt]
\noindent \textbf{Explanation:}\\
Trong Python, phương thức dựng sẵn của kiểu dữ liệu chuỗi (`str`) dùng để chuyển đổi tất cả các ký tự chữ hoa trong chuỗi thành chữ thường là phương thức `str.lower()`.\newline
- Ví dụ: `'PYTHON'.lower()` trả về `'python'`.\newline
- `toLower()` hoặc `toLowerCase()` là các phương thức trong các ngôn ngữ khác như Java, JavaScript, C#, không có trong Python.\newline
- `lowercase()` không phải là phương thức hợp lệ của kiểu chuỗi trong Python.
\end{minipage}
\vspace{0.8cm}
\hrule
\vspace{0.5cm}

\begin{minipage}{\textwidth}
\subsection*{Question 133 \quad \footnotesize{[0133_q133_Question_133.png]}}
\noindent \textbf{How would you use the index operator [] to print out the letter q from the following string?\newline

\newline
\newline

\newline
x = 'From marquard@uct.ac.za'}

\begin{enumerate}[label=\Alph*.]
    \item print(x[q])
    \item print x[10]
    \item print(x[9])
    \item \textbf{print(x[8])} \hfill \textcolor{blue}{\textbf{[Correct]}}
\end{enumerate}

\noindent \textbf{Correct Answer:} \textcolor{blue}{\textbf{D}} \\[4pt]
\noindent \textbf{Explanation:}\\
Trong Python, các ký tự trong chuỗi được đánh chỉ mục từ 0 (zero-based index):\newline
- x[0] = 'F'\newline
- x[1] = 'r'\newline
- x[2] = 'o'\newline
- x[3] = 'm'\newline
- x[4] = ' ' (khoảng trắng)\newline
- x[5] = 'm'\newline
- x[6] = 'a'\newline
- x[7] = 'r'\newline
- x[8] = 'q'\newline
Ký tự `'q'` nằm tại vị trí chỉ số 8. Do đó câu lệnh chính xác để in ra ký tự này là `print(x[8])`.\newline
- A sai vì `q` là biến chưa định nghĩa và index phải là số nguyên.\newline
- B sai vì chỉ số 10 là chữ `'a'`, đồng thời trong Python 3 `print` là một hàm nên cần dấu ngoặc đơn.\newline
- C sai vì chỉ số 9 là chữ `'u'`.
\end{minipage}
\vspace{0.8cm}
\hrule
\vspace{0.5cm}

\begin{minipage}{\textwidth}
\subsection*{Question 134 \quad \footnotesize{[0134_q134_Question_134.png]}}
\noindent \textbf{What is a correct syntax to output "Hello World" in Python?}

\begin{enumerate}[label=\Alph*.]
    \item echo "Hello World"
    \item p("Hello World")
    \item \textbf{print("Hello World")} \hfill \textcolor{blue}{\textbf{[Correct]}}
    \item echo("Hello World");
\end{enumerate}

\noindent \textbf{Correct Answer:} \textcolor{blue}{\textbf{C}} \\[4pt]
\noindent \textbf{Explanation:}\\
Trong Python (đặc biệt là Python 3), hàm chuẩn tích hợp sẵn dùng để in dữ liệu ra màn hình/dòng lệnh là `print()`. Vì vậy cú pháp chính xác là `print("Hello World")`.\newline
- `echo` là lệnh của hệ vỏ (Shell/Bash) hoặc PHP, không có trong Python.\newline
- `p()` không phải là hàm dựng sẵn của Python.\newline
- Dấu chấm phẩy `;` kết thúc câu lệnh trong phương án D là cú pháp của C/C++/PHP, không phải cú pháp chuẩn của Python.
\end{minipage}
\vspace{0.8cm}
\hrule
\vspace{0.5cm}

\begin{minipage}{\textwidth}
\subsection*{Question 135 \quad \footnotesize{[0135_q135_Question_135.png]}}
\noindent \textbf{What does the following Python code do?\newline

\newline
\newline

\newline
fhand = open('mbox-short.txt')\newline

\newline
inp = fhand.read()}

\begin{enumerate}[label=\Alph*.]
    \item Checks to see if the file exists and can be written
    \item Turns the text in the file into a graphic image like a PNG or JPG
    \item \textbf{Reads the entire file into the variable inp as a string} \hfill \textcolor{blue}{\textbf{[Correct]}}
    \item Prompts the user for a file name
\end{enumerate}

\noindent \textbf{Correct Answer:} \textcolor{blue}{\textbf{C}} \\[4pt]
\noindent \textbf{Explanation:}\\
Trong Python, phương thức `.read()` của đối tượng file (`file handle`) đọc toàn bộ nội dung của tệp tin từ vị trí con trỏ hiện tại đến cuối tệp (kể cả các ký tự xuống dòng `\n`) thành một chuỗi ký tự duy nhất (string) và gán vào biến `inp`.\newline
- Lựa chọn A sai vì `open()` mở file mặc định ở chế độ đọc (`'r'`), không kiểm tra khả năng ghi (`'w'`).\newline
- Lựa chọn B sai vì phương thức `read()` chỉ xử lý văn bản thành chuỗi ký tự, không chuyển đổi thành ảnh đồ họa.\newline
- Lựa chọn D sai vì tên tệp tin `'mbox-short.txt'` được truyền trực tiếp vào hàm `open()`, không hề có lệnh `input()` để nhắc người dùng nhập tên file.
\end{minipage}
\vspace{0.8cm}
\hrule
\vspace{0.5cm}

\begin{minipage}{\textwidth}
\subsection*{Question 136 \quad \footnotesize{[0136_q136_Question_136.png]}}
\noindent \textbf{What is the purpose of the newline character in text files?}

\begin{enumerate}[label=\Alph*.]
    \item It adds a new network connection to retrieve files from the network
    \item \textbf{It indicates the end of one line of text and the beginning of another line of text} \hfill \textcolor{blue}{\textbf{[Correct]}}
    \item It enables random movement throughout the file
    \item It allows us to open more than one files and read them in a synchronized manner
\end{enumerate}

\noindent \textbf{Correct Answer:} \textcolor{blue}{\textbf{B}} \\[4pt]
\noindent \textbf{Explanation:}\\
Trong các tệp văn bản (text files), ký tự xuống dòng (newline character, biểu diễn bằng `\n` trong Python) là một ký tự điều khiển dùng để đánh dấu điểm kết thúc của một dòng văn bản hiện tại và bắt đầu dòng văn bản tiếp theo.\newline
- A sai vì ký tự xuống dòng không liên quan gì đến kết nối mạng.\newline
- C sai vì di chuyển ngẫu nhiên trong file sử dụng phương thức `seek()`.\newline
- D sai vì việc mở nhiều file độc lập hoàn toàn với ký tự xuống dòng.
\end{minipage}
\vspace{0.8cm}
\hrule
\vspace{0.5cm}

\begin{minipage}{\textwidth}
\subsection*{Question 137 \quad \footnotesize{[0137_q137_Question_137.png]}}
\noindent \textbf{If you write a Python program to read a text file and you see extra blank lines in the output that are not present in the file input as shown below, what Python string function will likely solve the problem?\newline

\newline
\newline

\newline
From: stephen.marquard@uct.ac.za\newline

\newline
\newline

\newline
From: louis@media.berkeley.edu\newline

\newline
\newline

\newline
From: zqian@umich.edu\newline

\newline
\newline

\newline
From: rjlowe@iupui.edu\newline

\newline
\newline

\newline
...}

\begin{enumerate}[label=\Alph*.]
    \item find()
    \item startswith()
    \item \textbf{rstrip()} \hfill \textcolor{blue}{\textbf{[Correct]}}
    \item split()
\end{enumerate}

\noindent \textbf{Correct Answer:} \textcolor{blue}{\textbf{C}} \\[4pt]
\noindent \textbf{Explanation:}\\
Khi duyệt qua từng dòng của tệp văn bản bằng vòng lặp (như `for line in fhand:`), chuỗi `line` đã kết thúc sẵn bằng một ký tự xuống dòng (`\n`). Khi truyền chuỗi này vào hàm `print(line)`, `print` mặc định lại tự động chèn thêm một ký tự xuống dòng nữa (`end='\n'`). Kết quả là có hai ký tự xuống dòng liên tiếp, làm xuất hiện các dòng trống thừa xen kẽ.\newline
- Để giải quyết vấn đề này, phương thức `line.rstrip()` được sử dụng để cắt bỏ ký tự xuống dòng và khoảng trắng ở phía bên phải (cuối chuỗi).\newline
- `find()` dùng để tìm vị trí xuất hiện của chuỗi con.\newline
- `startswith()` kiểm tra chuỗi có bắt đầu bằng tiền tố nhất định hay không.\newline
- `split()` chia chuỗi thành danh sách các từ.
\end{minipage}
\vspace{0.8cm}
\hrule
\vspace{0.5cm}

\begin{minipage}{\textwidth}
\subsection*{Question 138 \quad \footnotesize{[0138_q138_Question_138.png]}}
\noindent \textbf{Which of the following statements assigns the value 100 to the variable x in Python:}

\begin{enumerate}[label=\Alph*.]
    \item x := 100
    \item let x = 100
    \item x << 100
    \item \textbf{x = 100} \hfill \textcolor{blue}{\textbf{[Correct]}}
    \item x == 100
\end{enumerate}

\noindent \textbf{Correct Answer:} \textcolor{blue}{\textbf{D}} \\[4pt]
\noindent \textbf{Explanation:}\\
Trong Python, câu lệnh gán giá trị tiêu chuẩn cho một biến sử dụng toán tử gán đơn `=`. Cú pháp `x = 100` gán giá trị số nguyên 100 cho biến `x`.\newline
- Phương án A (`x := 100`): `:=` là toán tử gán biểu thức (walrus operator), dùng trong biểu thức (expression) chứ không phải câu lệnh gán biến độc lập thông thường.\newline
- Phương án B (`let x = 100`): `let` là từ khóa khai báo biến trong JavaScript/Rust, không tồn tại trong Python.\newline
- Phương án C (`x << 100`): `<<` là toán tử dịch bit sang trái (bitwise left shift).\newline
- Phương án E (`x == 100`): `==` là toán tử so sánh bằng (equality operator), trả về `True` hoặc `False`, không dùng để gán giá trị.
\end{minipage}
\vspace{0.8cm}
\hrule
\vspace{0.5cm}

\begin{minipage}{\textwidth}
\subsection*{Question 139 \quad \footnotesize{[0139_q139_Question_139.png]}}
\noindent \textbf{What is the output of the following python code snippet?\newline

\newline
\newline

\newline
# Code starts\newline

\newline
emp1 = ('Ajay',25,10000)\newline

\newline
emp2 = ('Sidhi',23,15000)\newline

\newline
\newline

\newline
emp1[2] = 20000\newline

\newline
\newline

\newline
print(emp1[2]+emp2[2])\newline

\newline
# Code Ends}

\begin{enumerate}[label=\Alph*.]
    \item 35000
    \item 25000
    \item \textbf{Type Error} \hfill \textcolor{blue}{\textbf{[Correct]}}
    \item None of the above.
\end{enumerate}

\noindent \textbf{Correct Answer:} \textcolor{blue}{\textbf{C}} \\[4pt]
\noindent \textbf{Explanation:}\\
Phân tích các bước thực thi:\newline
1. `emp1` và `emp2` được khởi tạo dưới dạng `tuple` (chứa các phần tử đặt trong dấu ngoặc tròn `()`).\newline
2. Trong Python, `tuple` là kiểu dữ liệu bất biến (immutable sequence), tức là không thể sửa đổi, thêm hoặc gán lại giá trị cho các phần tử sau khi đã tạo.\newline
3. Dòng lệnh `emp1[2] = 20000` cố gắng thay đổi phần tử tại chỉ số 2 của tuple `emp1`. Điều này sẽ kích hoạt ngoại lệ `TypeError: 'tuple' object does not support item assignment`.\newline
4. Do đó chương trình sẽ dừng lại và báo lỗi `Type Error`.
\end{minipage}
\vspace{0.8cm}
\hrule
\vspace{0.5cm}

\begin{minipage}{\textwidth}
\subsection*{Question 140 \quad \footnotesize{[0140_q140_Question_140.png]}}
\noindent \textbf{Find the output of the code given below?\newline

\newline
\newline

\newline
# Code starts\newline

\newline
\newline

\newline
content = (21, 26, 37, 94)\newline

\newline
content.append( (75, 62, 79) )\newline

\newline
print(len(content))\newline

\newline
\newline

\newline
# Code Ends}

\begin{enumerate}[label=\Alph*.]
    \item 7
    \item 2
    \item 1
    \item \textbf{Error} \hfill \textcolor{blue}{\textbf{[Correct]}}
\end{enumerate}

\noindent \textbf{Correct Answer:} \textcolor{blue}{\textbf{D}} \\[4pt]
\noindent \textbf{Explanation:}\\
Phân tích các bước thực thi:\newline
1. `content = (21, 26, 37, 94)`: Khởi tạo biến `content` là một `tuple` chứa 4 số nguyên.\newline
2. `tuple` trong Python là kiểu dữ liệu bất biến (immutable) và đối tượng tuple hoàn toàn không có phương thức `.append()` (phương thức này chỉ có ở kiểu danh sách `list`).\newline
3. Lệnh `content.append((75, 62, 79))` sẽ phát sinh ngoại lệ `AttributeError: 'tuple' object has no attribute 'append'`.\newline
4. Do đó chương trình dừng lại do lỗi (`Error`).
\end{minipage}
\vspace{0.8cm}
\hrule
\vspace{0.5cm}

\begin{minipage}{\textwidth}
\subsection*{Question 141 \quad \footnotesize{[0141_q141_Question_141.png]}}
\noindent \textbf{What is the correct way to create a function in Python?}

\begin{enumerate}[label=\Alph*.]
    \item def function()
    \item define function()
    \item function()
    \item creat function()
    \item \textbf{def function():} \hfill \textcolor{blue}{\textbf{[Correct]}}
\end{enumerate}

\noindent \textbf{Correct Answer:} \textcolor{blue}{\textbf{E}} \\[4pt]
\noindent \textbf{Explanation:}\\
Trong Python, cú pháp chuẩn để định nghĩa một hàm bao gồm từ khóa `def`, theo sau là tên hàm, cặp dấu ngoặc đơn `()` chứa các tham số (nếu có) và kết thúc bằng dấu hai chấm `:`. Khối lệnh của thân hàm được thụt lề (indentation).\newline
- Cú pháp đúng: `def function():` (Đáp án E).\newline
- Lựa chọn A sai vì thiếu dấu hai chấm `:` ở cuối dòng định nghĩa.\newline
- Lựa chọn B, C, D sai vì Python không sử dụng từ khóa `define`, `creat` hay cặp dấu ngoặc nhọn `{}` để bao bọc khối lệnh hàm như trong C, Java hay JavaScript.
\end{minipage}
\vspace{0.8cm}
\hrule
\vspace{0.5cm}

\begin{minipage}{\textwidth}
\subsection*{Question 142 \quad \footnotesize{[0142_q142_Question_142.png]}}
\noindent \textbf{What is the result of the following statement:\newline

\newline
\newline

\newline
```python\newline

\newline
list((b'abcde' + b'fghi')[3:6])\newline

\newline
```}

\begin{enumerate}[label=\Alph*.]
    \item [b'd', b'e', b'f']
    \item b'def'
    \item [100, 101, 102]
    \item \textbf{It raises an exception} \hfill \textcolor{blue}{\textbf{[Correct]}}
\end{enumerate}

\noindent \textbf{Correct Answer:} \textcolor{blue}{\textbf{D}} \\[4pt]
\noindent \textbf{Explanation:}\\
- Theo đáp án chính thức của hệ thống đề thi / ngân hàng câu hỏi PFP191, câu trả lời được chọn là D ('It raises an exception').\newline
- Phân tích kỹ thuật chuyên sâu theo chuẩn Python 3:\newline
  1. Phép cộng hai đối tượng bytes `b'abcde' + b'fghi'` tạo thành `b'abcdefghi'`. Cắt lát (slicing) `[3:6]` lấy các phần tử từ chỉ số 3 đến 5 (các ký tự byte tại vị trí 3, 4, 5 tương ứng là 'd', 'e', 'f'), kết quả là đối tượng bytes `b'def'`.\newline
  2. Trong Python 3, khi chuyển đổi một đối tượng `bytes` thành `list` bằng hàm `list()`, Python lặp qua các byte và trả về danh sách mã ASCII số nguyên của từng byte: `ord('d') = 100`, `ord('e') = 101`, `ord('f') = 102`, do đó kết quả thực tế trên Python 3 console là `[100, 101, 102]` (tương ứng với phương án C). Tuy nhiên, trong đề thi trắc nghiệm nguồn, câu hỏi này được quy ước đáp án là D.
\end{minipage}
\vspace{0.8cm}
\hrule
\vspace{0.5cm}

\begin{minipage}{\textwidth}
\subsection*{Question 143 \quad \footnotesize{[0143_q143_Question_143.png]}}
\noindent \textbf{Which of these collections defines a LIST?}

\begin{enumerate}[label=\Alph*.]
    \item "name": "apple", "color": "green"
    \item \textbf{["apple", "banana", "cherry"]} \hfill \textcolor{blue}{\textbf{[Correct]}}
    \item "apple", "banana", "cherry"
    \item ("apple", "banana", "cherry")
\end{enumerate}

\noindent \textbf{Correct Answer:} \textcolor{blue}{\textbf{B}} \\[4pt]
\noindent \textbf{Explanation:}\\
Trong Python, các kiểu tập hợp dữ liệu (collections) được phân biệt bằng các loại dấu ngoặc:\newline
- `[...]` (cặp ngoặc vuông): Định nghĩa kiểu danh sách (`list`). Do đó `["apple", "banana", "cherry"]` là một List (Đáp án B đúng).\newline
- `{"name": "apple", "color": "green"}`: Sử dụng ngoặc nhọn kèm các cặp `key: value`, định nghĩa kiểu từ điển (`dict`).\newline
- `{"apple", "banana", "cherry"}`: Sử dụng ngoặc nhọn chứa các phần tử phân tách bởi dấu phẩy, định nghĩa kiểu tập hợp (`set`).\newline
- `("apple", "banana", "cherry")`: Sử dụng ngoặc tròn, định nghĩa kiểu bộ dữ liệu không thay đổi (`tuple`).
\end{minipage}
\vspace{0.8cm}
\hrule
\vspace{0.5cm}

\begin{minipage}{\textwidth}
\subsection*{Question 144 \quad \footnotesize{[0144_q144_Question_144.png]}}
\noindent \textbf{If you want to transform a list of strings input_list into a single string with a comma between each item, which of the following would you give as the input to join()?}

\begin{enumerate}[label=\Alph*.]
    \item ','
    \item string
    \item \textbf{input_list} \hfill \textcolor{blue}{\textbf{[Correct]}}
    \item str
\end{enumerate}

\noindent \textbf{Correct Answer:} \textcolor{blue}{\textbf{C}} \\[4pt]
\noindent \textbf{Explanation:}\\
Cú pháp của phương thức nối chuỗi trong Python là `str_separator.join(iterable)`.\newline
- Trong đó, `str_separator` là chuỗi phân cách gọi phương thức `join()`, ví dụ: `','`.\newline
- Đối số đầu vào truyền vào trong dấu ngoặc của phương thức `join(...)` phải là một iterable chứa các chuỗi cần nối, ở đây chính là danh sách `input_list`.\newline
- Câu lệnh hoàn chỉnh là: `','.join(input_list)`.\newline
- Do câu hỏi hỏi 'giá trị nào được truyền làm đầu vào (input/argument) cho `join()`?', đáp án chính xác là C (`input_list`).
\end{minipage}
\vspace{0.8cm}
\hrule
\vspace{0.5cm}

\begin{minipage}{\textwidth}
\subsection*{Question 145 \quad \footnotesize{[0145_q145_Question_145.png]}}
\noindent \textbf{What do we do to a Python statement that is immediately after an if statement to indicate that the statement is to be executed only when the if statement is true?}

\begin{enumerate}[label=\Alph*.]
    \item Underline all of the conditional code
    \item \textbf{Indent the line below the if statement} \hfill \textcolor{blue}{\textbf{[Correct]}}
    \item Start the statement with a "#" character
    \item Begin the statement with a curly brace
\end{enumerate}

\noindent \textbf{Correct Answer:} \textcolor{blue}{\textbf{B}} \\[4pt]
\noindent \textbf{Explanation:}\\
Trong Python, cú pháp khối lệnh không sử dụng cặp ngoặc nhọn `{}` hay từ khóa `begin`/`end` như các ngôn ngữ khác, mà sử dụng việc thụt lề (indentation):\newline
- Dòng lệnh nằm ngay sau mệnh đề `if` phải được thụt lề (thường là 4 dấu cách) để chỉ định rằng dòng lệnh đó thuộc khối mã (suite) của `if` và chỉ được thực thi khi điều kiện `if` đánh giá là `True`.\newline
- Do đó, đáp án B ('Indent the line below the if statement') là chính xác.\newline
- Các lựa chọn khác đều sai: gạch chân không phải cú pháp code, dấu `#` dùng để tạo chú thích (comment), và `{` dùng cho dictionary/set.
\end{minipage}
\vspace{0.8cm}
\hrule
\vspace{0.5cm}

\begin{minipage}{\textwidth}
\subsection*{Question 146 \quad \footnotesize{[0146_q146_Question_146.png]}}
\noindent \textbf{What will be the value of x after the following statement executes:\newline

\newline
\newline

\newline
```python\newline

\newline
x = 2 * 3 - 8 / 4 % 2\newline

\newline
```}

\begin{enumerate}[label=\Alph*.]
    \item 1
    \item 2
    \item 4
    \item \textbf{6} \hfill \textcolor{blue}{\textbf{[Correct]}}
\end{enumerate}

\noindent \textbf{Correct Answer:} \textcolor{blue}{\textbf{D}} \\[4pt]
\noindent \textbf{Explanation:}\\
Xét thứ tự ưu tiên của các toán tử (operator precedence) trong biểu thức:\newline
1. Các toán tử nhân `*`, chia `/`, chia lấy dư `%` có độ ưu tiên cao hơn toán tử trừ `-` và có cùng mức ưu tiên với nhau, nên được tính toán lần lượt từ trái sang phải:\newline
   - `2 * 3 = 6`\newline
   - `8 / 4 = 2.0` (toán tử chia `/` trong Python 3 luôn trả về kiểu float)\newline
   - `2.0 % 2 = 0.0`\newline
2. Tiếp theo thực hiện phép trừ: `6 - 0.0 = 6.0`.\newline
3. Giá trị kết quả là 6 (hoặc 6.0). Trong các lựa chọn trắc nghiệm, đáp án đúng là D (6).
\end{minipage}
\vspace{0.8cm}
\hrule
\vspace{0.5cm}

\begin{minipage}{\textwidth}
\subsection*{Question 147 \quad \footnotesize{[0147_q147_Question_147.png]}}
\noindent \textbf{Which of the following Python statements would print out the length of a list stored in the variable data?}

\begin{enumerate}[label=\Alph*.]
    \item print(length(data))
    \item print(data.length())
    \item \textbf{print(len(data))} \hfill \textcolor{blue}{\textbf{[Correct]}}
    \item print(data.length)
\end{enumerate}

\noindent \textbf{Correct Answer:} \textcolor{blue}{\textbf{C}} \\[4pt]
\noindent \textbf{Explanation:}\\
Trong Python, hàm chuẩn tích hợp sẵn (built-in function) để lấy số lượng phần tử (độ dài) của một tập hợp như list, tuple, string, dictionary... là `len()`.\newline
- Do đó, câu lệnh chính xác để in ra độ dài của biến danh sách `data` là `print(len(data))` (Đáp án C).\newline
- Python không hỗ trợ hàm `length()` (gây lỗi `NameError`), cũng không có thuộc tính hoặc phương thức `.length` hay `.length()` trên đối tượng list (gây lỗi `AttributeError`).
\end{minipage}
\vspace{0.8cm}
\hrule
\vspace{0.5cm}

\begin{minipage}{\textwidth}
\subsection*{Question 148 \quad \footnotesize{[0148_q148_Question_148.png]}}
\noindent \textbf{What is true about the following code segment:\newline

\newline
\newline

\newline
```python\newline

\newline
if x == 5:\newline

\newline
    print('Is 5')\newline

\newline
    print('Is Still 5')\newline

\newline
    print('Third 5')\newline

\newline
```}

\begin{enumerate}[label=\Alph*.]
    \item \textbf{Depending on the value of x, either all three of the print statements will execute or none of the statements will execute} \hfill \textcolor{blue}{\textbf{[Correct]}}
    \item The string 'Is 5' will always print out regardless of the value for x.
    \item The string 'Is 5' will never print out regardless of the value for x.
    \item Only two of the three print statements will print out if the value of x is less than zero.
\end{enumerate}

\noindent \textbf{Correct Answer:} \textcolor{blue}{\textbf{A}} \\[4pt]
\noindent \textbf{Explanation:}\\
Trong đoạn mã trên, cả ba câu lệnh `print` đều được thụt lề cùng cấp và nằm trong khối lệnh (body) của cấu trúc điều kiện `if x == 5:`:\newline
- Nếu giá trị của biến `x` bằng 5, điều kiện `x == 5` nhận giá trị `True`, toàn bộ cả ba câu lệnh `print` sẽ được thực thi tuần tự.\newline
- Nếu giá trị của biến `x` khác 5 (bất kể lớn hơn hay nhỏ hơn 5), điều kiện nhận giá trị `False`, toàn bộ khối lệnh thụt lề sẽ bị bỏ qua và không có câu lệnh `print` nào được chạy.\newline
- Vì vậy, đáp án A ('Depending on the value of x, either all three of the print statements will execute or none of the statements will execute') là khẳng định hoàn toàn chính xác.
\end{minipage}
\vspace{0.8cm}
\hrule
\vspace{0.5cm}

\begin{minipage}{\textwidth}
\subsection*{Question 149 \quad \footnotesize{[0149_q149_Question_149.png]}}
\noindent \textbf{```python\newline

\newline
class Person:\newline

\newline
    def __init__(self, id):\newline

\newline
        self.id = id\newline

\newline
\newline

\newline
sam = Person(100)\newline

\newline
sam.__dict__['age'] = 49\newline

\newline
print (sam.age + len(sam.__dict__))\newline

\newline
```}

\begin{enumerate}[label=\Alph*.]
    \item 49
    \item 50
    \item \textbf{51} \hfill \textcolor{blue}{\textbf{[Correct]}}
    \item 48
\end{enumerate}

\noindent \textbf{Correct Answer:} \textcolor{blue}{\textbf{C}} \\[4pt]
\noindent \textbf{Explanation:}\\
Từng bước thực thi đoạn mã:\newline
1. Khởi tạo đối tượng `sam = Person(100)`: Phương thức `__init__` gán `self.id = 100`. Lúc này từ điển thuộc tính của đối tượng `sam.__dict__` chứa: `{'id': 100}`.\newline
2. Lệnh `sam.__dict__['age'] = 49`: Thêm trực tiếp một mục mới vào từ điển thuộc tính của đối tượng với khóa `'age'` và giá trị 49. Bây giờ `sam.__dict__` trở thành `{'id': 100, 'age': 49}`.\newline
3. Truy cập thuộc tính `sam.age`: Python tìm kiếm trong `sam.__dict__` và lấy ra giá trị tương ứng là 49.\newline
4. Hàm `len(sam.__dict__)`: Đếm số lượng cặp key-value trong từ điển thuộc tính của `sam`. Từ điển có 2 khóa (`'id'` và `'age'`), do đó hàm trả về giá trị 2.\newline
5. Thực hiện phép cộng: `sam.age + len(sam.__dict__) = 49 + 2 = 51`.\newline
Do đó kết quả in ra màn hình là 51 (Đáp án C).
\end{minipage}
\vspace{0.8cm}
\hrule
\vspace{0.5cm}

\begin{minipage}{\textwidth}
\subsection*{Question 150 \quad \footnotesize{[0150_q150_Question_150.png]}}
\noindent \textbf{What's the output of the following code snippet?\newline

\newline
\newline

\newline
```python\newline

\newline
>>> class Dog:\newline

\newline
...     def walk(self):\newline

\newline
...         return "*walking*"\newline

\newline
...\newline

\newline
...     def speak(self):\newline

\newline
...         return "Woof!"\newline

\newline
...\newline

\newline
>>> class JackRussellTerrier(Dog):\newline

\newline
...     def talk(self):\newline

\newline
...         return super().speak()\newline

\newline
...\newline

\newline
>>> bobo = JackRussellTerrier()\newline

\newline
>>> bobo.talk()\newline

\newline
```}

\begin{enumerate}[label=\Alph*.]
    \item \textbf{Woof!} \hfill \textcolor{blue}{\textbf{[Correct]}}
    \item walking
    \item CanineError: Tail curvature exceeded
\end{enumerate}

\noindent \textbf{Correct Answer:} \textcolor{blue}{\textbf{A}} \\[4pt]
\noindent \textbf{Explanation:}\\
Phân tích đoạn mã kế thừa trong Python:\newline
1. Lớp `Dog` định nghĩa hai phương thức là `walk(self)` trả về `"*walking*"` và `speak(self)` trả về `"Woof!"`.\newline
2. Lớp `JackRussellTerrier` kế thừa từ lớp `Dog` thông qua cú pháp `class JackRussellTerrier(Dog):`.\newline
3. Trong lớp con `JackRussellTerrier`, phương thức `talk(self)` sử dụng hàm `super().speak()` để gọi phương thức `speak()` của lớp cha (`Dog`).\newline
4. Khi tạo instance `bobo = JackRussellTerrier()` và gọi phương thức `bobo.talk()`, hàm `super().speak()` được kích hoạt và trả về chuỗi `"Woof!"`.\newline
5. Trên interactive console (REPL), giá trị trả về của `bobo.talk()` sẽ được hiển thị là `'Woof!'` (phương án A: Woof!).
\end{minipage}
\vspace{0.8cm}
\hrule
\vspace{0.5cm}

\begin{minipage}{\textwidth}
\subsection*{Question 151 \quad \footnotesize{[0151_q151_Question_151.png]}}
\noindent \textbf{What would the following mean in a regular expression? [a-z0-9]}

\begin{enumerate}[label=\Alph*.]
    \item Match any number of lowercase letters followed by any number of digits
    \item \textbf{Match a lowercase letter or a digit} \hfill \textcolor{blue}{\textbf{[Correct]}}
    \item Match an entire line as long as it is lowercase letters or digits
    \item Match any text that is surrounded by square braces
\end{enumerate}

\noindent \textbf{Correct Answer:} \textcolor{blue}{\textbf{B}} \\[4pt]
\noindent \textbf{Explanation:}\\
Trong biểu thức chính quy (Regular Expressions):\newline
- Cặp ngoặc vuông `[...]` đại diện cho một tập ký tự (character class). Nó khớp với chính xác một ký tự đơn xuất hiện trong tập ký tự đó.\newline
- Dấu gạch ngang `-` trong ngoặc vuông định nghĩa một dải ký tự (range). Cụ thể, `a-z` đại diện cho bất kỳ chữ cái thường nào từ 'a' đến 'z', và `0-9` đại diện cho bất kỳ chữ số nào từ '0' đến '9'.\newline
- Do không có bộ định lượng số lượng (quantifier như `*`, `+`) đi kèm, `[a-z0-9]` có ý nghĩa là khớp với một ký tự chữ cái thường hoặc một chữ số ('Match a lowercase letter or a digit').\newline
- Vì vậy, đáp án B là chính xác.
\end{minipage}
\vspace{0.8cm}
\hrule
\vspace{0.5cm}

\begin{minipage}{\textwidth}
\subsection*{Question 152 \quad \footnotesize{[0152_q152_Question_152.png]}}
\noindent \textbf{What does the following Python Program print out?\newline

\newline
\newline

\newline
```python\newline

\newline
str1 = "Hello"\newline

\newline
str2 = 'there'\newline

\newline
bob = str1 + str2\newline

\newline
print(bob)\newline

\newline
```}

\begin{enumerate}[label=\Alph*.]
    \item Hello
    \item 0
    \item \textbf{Hellothere} \hfill \textcolor{blue}{\textbf{[Correct]}}
    \item Hello
\end{enumerate}

\noindent \textbf{Correct Answer:} \textcolor{blue}{\textbf{C}} \\[4pt]
\noindent \textbf{Explanation:}\\
Trong Python, toán tử `+` khi áp dụng giữa hai chuỗi ký tự (`str`) thực hiện thao tác ghép chuỗi (string concatenation):\newline
- Biến `str1` mang giá trị `"Hello"`.\newline
- Biến `str2` mang giá trị `'there'`.\newline
- Biểu thức `str1 + str2` nối hai chuỗi lại với nhau mà không tự động chèn thêm khoảng trắng ở giữa, tạo thành chuỗi `'Hellothere'`.\newline
- Lệnh `print(bob)` in chuỗi này ra màn hình, kết quả là `Hellothere` (Đáp án C).
\end{minipage}
\vspace{0.8cm}
\hrule
\vspace{0.5cm}

\begin{minipage}{\textwidth}
\subsection*{Question 153 \quad \footnotesize{[0153_q153_Question_153.png]}}
\noindent \textbf{Which method can be used to remove any whitespace from both the beginning and the end of a string?}

\begin{enumerate}[label=\Alph*.]
    \item ltrip()
    \item \textbf{strip()} \hfill \textcolor{blue}{\textbf{[Correct]}}
    \item ptrip()
    \item strim()
\end{enumerate}

\noindent \textbf{Correct Answer:} \textcolor{blue}{\textbf{B}} \\[4pt]
\noindent \textbf{Explanation:}\\
Trong Python, đối tượng chuỗi (`str`) cung cấp các phương thức tích hợp để loại bỏ khoảng trắng (whitespace: khoảng cách, tab, xuống dòng):\newline
- `strip()`: Loại bỏ khoảng trắng ở cả hai đầu (đầu và cuối) chuỗi. Đáp án B đúng.\newline
- `lstrip()`: Chỉ loại bỏ khoảng trắng ở đầu bên trái (leading whitespace).\newline
- `rstrip()`: Chỉ loại bỏ khoảng trắng ở cuối bên phải (trailing whitespace).\newline
- Các tên hàm như `ltrip()`, `ptrip()`, `strim()` không tồn tại trong Python.
\end{minipage}
\vspace{0.8cm}
\hrule
\vspace{0.5cm}

\begin{minipage}{\textwidth}
\subsection*{Question 154 \quad \footnotesize{[0154_q154_Question_154.png]}}
\noindent \textbf{What is the output for: 'Tutorials Point' [100:200]?}

\begin{enumerate}[label=\Alph*.]
    \item Index error.
    \item \textbf{''} \hfill \textcolor{blue}{\textbf{[Correct]}}
    \item 'Tutorials Point'
    \item Syntax error
\end{enumerate}

\noindent \textbf{Correct Answer:} \textcolor{blue}{\textbf{B}} \\[4pt]
\noindent \textbf{Explanation:}\\
Trong Python, quy tắc truy xuất chuỗi:\newline
- Khi truy cập một chỉ số đơn lẻ nằm ngoài phạm vi chuỗi, ví dụ `'abc'[10]`, Python sẽ ném ngoại lệ `IndexError: string index out of range`.\newline
- Tuy nhiên, khi thực hiện phép cắt lát chuỗi (slicing) với cú pháp `s[start:end]`, Python xử lý các chỉ số vượt biên một cách an toàn (graceful handling): nếu cả hai chỉ số `start` và `end` đều vượt quá độ dài của chuỗi (chuỗi `'Tutorials Point'` chỉ có độ dài 15 ký tự, trong khi `start=100` và `end=200`), Python sẽ trả về một chuỗi rỗng `''` thay vì báo lỗi.\newline
- Do đó, kết quả là `''` (Đáp án B).
\end{minipage}
\vspace{0.8cm}
\hrule
\vspace{0.5cm}

\begin{minipage}{\textwidth}
\subsection*{Question 155 \quad \footnotesize{[0155_q155_Question_155.png]}}
\noindent \textbf{Which of the following command is used to open a file "c:temp .txt" in write-mode only?}

\begin{enumerate}[label=\Alph*.]
    \item \textbf{outfile = open("c:temp .txt", "w")} \hfill \textcolor{blue}{\textbf{[Correct]}}
    \item outfile = open("c: temp.txt", "w")
    \item outfile = open(file = "c:temp .txt", "w+")
    \item outfile = open(file = "c: temp.txt", "w+")
\end{enumerate}

\noindent \textbf{Correct Answer:} \textcolor{blue}{\textbf{A}} \\[4pt]
\noindent \textbf{Explanation:}\\
- Để mở tệp ở chế độ 'chỉ ghi' (write-mode only), đối số mode được sử dụng là `"w"`. Chế độ `"w+"` là chế độ vừa ghi vừa đọc (read and write), do đó loại trừ các đáp án C và D.\newline
- Theo quy ước câu hỏi và đáp án chính thức của hệ thống đề thi FUO / ngân hàng câu hỏi PFP191, đáp án được chọn là A: `outfile = open("c:\temp.txt", "w")`.\newline
- Lưu ý về mặt kỹ thuật trong thực tế: Do ký tự `\t` có thể bị hiểu là ký tự tab escape sequence trong chuỗi Python, khi làm việc thực tế trên Windows nên sử dụng ký tự escape `"c:\\temp.txt"` (hoặc raw string `r"c:\temp.txt"`, hoặc dấu gạch chéo xuôi `"c:/temp.txt"`) để tránh lỗi phân giải đường dẫn. Tuy nhiên với đề thi trắc nghiệm lý thuyết gốc, đáp án chính thức là A.
\end{minipage}
\vspace{0.8cm}
\hrule
\vspace{0.5cm}

\begin{minipage}{\textwidth}
\subsection*{Question 156 \quad \footnotesize{[0156_q156_Question_156.png]}}
\noindent \textbf{Which operator can be used to compare two values?}

\begin{enumerate}[label=\Alph*.]
    \item <>
    \item =
    \item ><
    \item \textbf{==} \hfill \textcolor{blue}{\textbf{[Correct]}}
\end{enumerate}

\noindent \textbf{Correct Answer:} \textcolor{blue}{\textbf{D}} \\[4pt]
\noindent \textbf{Explanation:}\\
Trong Python, các toán tử được phân loại rõ ràng theo chức năng:\newline
- Toán tử `==` (equality operator) được dùng để so sánh giá trị của hai toán hạng xem chúng có bằng nhau hay không. Biểu thức trả về giá trị boolean `True` hoặc `False`. Do đó đáp án đúng là D.\newline
- Toán tử `=` là toán tử gán (assignment operator), dùng để gán giá trị cho một biến, không dùng để so sánh.\newline
- Ký hiệu `<>` là toán tử so sánh khác từng có trong Python 2 nhưng đã bị loại bỏ hoàn toàn trong Python 3 (thay thế bằng toán tử `!=`).\newline
- `><` không phải là một toán tử hợp lệ trong Python (gây lỗi cú pháp `SyntaxError`).
\end{minipage}
\vspace{0.8cm}
\hrule
\vspace{0.5cm}

\begin{minipage}{\textwidth}
\subsection*{Question 157 \quad \footnotesize{[0157_q157_Question_157.png]}}
\noindent \textbf{What is the proper way to say "good-bye" to Python?}

\begin{enumerate}[label=\Alph*.]
    \item \textbf{quit()} \hfill \textcolor{blue}{\textbf{[Correct]}}
    \item stop()
    \item exit()
    \item end()
\end{enumerate}

\noindent \textbf{Correct Answer:} \textcolor{blue}{\textbf{A}} \\[4pt]
\noindent \textbf{Explanation:}\\
Trong giáo trình Python for Everybody (Chương 1 - của Giáo sư Charles Severance / Dr. Chuck, giáo trình nền tảng của môn học PFP191 tại FPT University):\newline
- Tác giả đặt câu hỏi: 'How do you tell Python good-bye? You type quit() at the interactive chevron >>> prompt.'\newline
- Khi làm việc trong chế độ thông dịch tương tác (interactive Python shell), câu lệnh chuẩn được giới thiệu để thoát khỏi môi trường Python là gọi hàm `quit()` (hoặc `exit()`, nhưng câu hỏi trắc nghiệm của giáo trình chọn `quit()`).\newline
- Các hàm như `stop()` hay `end()` không phải là lệnh chuẩn để thoát môi trường tương tác của Python.\newline
- Do đó, đáp án đúng theo giáo trình là A (`quit()`).
\end{minipage}
\vspace{0.8cm}
\hrule
\vspace{0.5cm}

\begin{minipage}{\textwidth}
\subsection*{Question 158 \quad \footnotesize{[0158_q158_Question_158.png]}}
\noindent \textbf{A USB memory stick is an example of which of the following components of computer architecture?}

\begin{enumerate}[label=\Alph*.]
    \item Central Processing Unit
    \item Main Memory
    \item \textbf{Secondary Memory} \hfill \textcolor{blue}{\textbf{[Correct]}}
    \item None of these
\end{enumerate}

\noindent \textbf{Correct Answer:} \textcolor{blue}{\textbf{C}} \\[4pt]
\noindent \textbf{Explanation:}\\
Trong kiến trúc máy tính cơ bản (được giảng dạy trong chương mở đầu về Computer Hardware Architecture của môn PFP191):\newline
1. Central Processing Unit (CPU): Bộ xử lý trung tâm, có nhiệm vụ thực thi các chỉ lệnh và điều khiển các thành phần khác.\newline
2. Main Memory (RAM): Bộ nhớ chính, lưu trữ thông tin tạm thời tốc độ cao khi máy tính đang hoạt động; dữ liệu sẽ mất đi khi mất nguồn điện (volatile memory).\newline
3. Secondary Memory (Bộ nhớ thứ cấp/lưu trữ ngoài): Dùng để lưu trữ dữ liệu vĩnh viễn, không bị mất khi tắt máy (non-volatile memory). Các thiết bị thứ cấp bao gồm ổ cứng (HDD/SSD), đĩa quang, thẻ nhớ và thẻ nhớ flash USB (USB memory stick).\newline
Do đó, USB memory stick là một ví dụ tiêu biểu của Secondary Memory (Đáp án C).
\end{minipage}
\vspace{0.8cm}
\hrule
\vspace{0.5cm}

\begin{minipage}{\textwidth}
\subsection*{Question 159 \quad \footnotesize{[0159_q159_Question_159.png]}}
\noindent \textbf{How would you express the constant floating-point value 3.2 × 10³ in Python?}

\begin{enumerate}[label=\Alph*.]
    \item 3.2e-3
    \item \textbf{3.2e3} \hfill \textcolor{blue}{\textbf{[Correct]}}
    \item 0.32
    \item 32e-3
\end{enumerate}

\noindent \textbf{Correct Answer:} \textcolor{blue}{\textbf{B}} \\[4pt]
\noindent \textbf{Explanation:}\\
Trong Python, số thực dấu phẩy động (float literal) có thể được biểu diễn dưới dạng ký hiệu khoa học (scientific E-notation) theo định dạng `<hệ số>e<số mũ>` hoặc `<hệ số>E<số mũ>`, đại diện cho phép tính `<hệ số> × 10^(số mũ)`:\newline
- Giá trị 3.2 × 10³ có hệ số là 3.2 và số mũ là +3, do đó được viết trong Python là `3.2e3` (hoặc `3.2e+3`), mang giá trị kiểu float tương đương `3200.0`.\newline
- Phương án A: `3.2e-3` biểu diễn cho 3.2 × 10⁻³ = 0.0032.\newline
- Phương án C: `0.32`.\newline
- Phương án D: `32e-3` biểu diễn cho 32 × 10⁻³ = 0.032.\newline
Vì vậy, đáp án đúng là B.
\end{minipage}
\vspace{0.8cm}
\hrule
\vspace{0.5cm}

\begin{minipage}{\textwidth}
\subsection*{Question 160 \quad \footnotesize{[0160_q160_Question_160.png]}}
\noindent \textbf{What is the output of the following code?\newline

\newline
\newline

\newline
```python\newline

\newline
class Point:\newline

\newline
    def __init__(self, x = 0, y = 0):\newline

\newline
        self.x = x\newline

\newline
        self.y = y\newline

\newline
\newline

\newline
    def __sub__(self, other):\newline

\newline
        x = self.x + other.x\newline

\newline
        y = self.y + other.y\newline

\newline
        return Point(x,y)\newline

\newline
\newline

\newline
p1 = Point(3, 4)\newline

\newline
p2 = Point(1, 2)\newline

\newline
result = p1-p2\newline

\newline
print(result.x, result.y)\newline

\newline
```}

\begin{enumerate}[label=\Alph*.]
    \item 2 2
    \item \textbf{4 6} \hfill \textcolor{blue}{\textbf{[Correct]}}
    \item 0 0
    \item 1 1
\end{enumerate}

\noindent \textbf{Correct Answer:} \textcolor{blue}{\textbf{B}} \\[4pt]
\noindent \textbf{Explanation:}\\
Theo dõi từng bước thực thi chương trình:\newline
1. Định nghĩa lớp `Point`: Phương thức nạp chồng toán tử trừ `__sub__(self, other)` được định nghĩa lại. Chú ý rằng bên trong thân hàm, tác giả cố tình viết logic cộng tọa độ thay vì trừ:\newline
   `x = self.x + other.x`\newline
   `y = self.y + other.y`\newline
   `return Point(x, y)`\newline
2. Khởi tạo đối tượng `p1 = Point(3, 4)`: `p1.x = 3`, `p1.y = 4`.\newline
3. Khởi tạo đối tượng `p2 = Point(1, 2)`: `p2.x = 1`, `p2.y = 2`.\newline
4. Thực hiện phép toán `result = p1 - p2`: Biểu thức này tương đương với lời gọi phương thức `result = p1.__sub__(p2)`.\newline
   - `x = p1.x + p2.x = 3 + 1 = 4`\newline
   - `y = p1.y + p2.y = 4 + 2 = 6`\newline
   - Trả về đối tượng mới `Point(4, 6)` gán cho biến `result`.\newline
5. Lệnh `print(result.x, result.y)` in ra giá trị của hai thuộc tính: `4 6`.\newline
Do đó, đáp án đúng là B (`4 6`).
\end{minipage}
\vspace{0.8cm}
\hrule
\vspace{0.5cm}

\begin{minipage}{\textwidth}
\subsection*{Question 161 \quad \footnotesize{[0161_q161_Question_161.png]}}
\noindent \textbf{What does the __init__() the function do in Python?}

\begin{enumerate}[label=\Alph*.]
    \item Initializes the class for use.
    \item \textbf{This function is called when a new object is instantiated.} \hfill \textcolor{blue}{\textbf{[Correct]}}
    \item Initializes all the data attributes to zero when called.
    \item None of the above.
\end{enumerate}

\noindent \textbf{Correct Answer:} \textcolor{blue}{\textbf{B}} \\[4pt]
\noindent \textbf{Explanation:}\\
Trong lập trình hướng đối tượng (OOP) với Python, phương thức `__init__()` là một hàm khởi tạo (constructor / initializer) đặc biệt. Khi một đối tượng mới (instance) được khởi tạo từ lớp (ví dụ: `p = Person()`), Python sẽ tự động gọi phương thức `__init__()` này để gán giá trị và thiết lập trạng thái ban đầu cho các thuộc tính của đối tượng.\newline
- Phương án A sai vì `__init__()` khởi tạo thực thể (object/instance), không phải khởi tạo bản thân class.\newline
- Phương án C sai vì `__init__()` không tự động gán tất cả thuộc tính về 0, mà việc khởi tạo phụ thuộc vào các tham số và câu lệnh gán do lập trình viên định nghĩa.\newline
- Phương án D sai vì B là đáp án chính xác.
\end{minipage}
\vspace{0.8cm}
\hrule
\vspace{0.5cm}

\begin{minipage}{\textwidth}
\subsection*{Question 162 \quad \footnotesize{[0162_q162_Question_162.png]}}
\noindent \textbf{List a is defined as follows: a = ['a', 'b', 'c'] Which of the following statements adds 'd' and 'e' to the end of a, so that it then equals ['a', 'b', 'c', 'd', 'e']:}

\begin{enumerate}[label=\Alph*.]
    \item \textbf{a.extend(['d', 'e'])} \hfill \textcolor{blue}{\textbf{[Correct]}}
    \item a += ['d', 'e']
    \item a += 'de'
    \item a[len(a):] = ['d', 'e']
\end{enumerate}

\noindent \textbf{Correct Answer:} \textcolor{blue}{\textbf{A}} \\[4pt]
\noindent \textbf{Explanation:}\\
Phương thức chuẩn và phổ biến nhất được thiết kế riêng trong Python để thêm lần lượt từng phần tử của một iterable vào cuối một danh sách có sẵn (in-place) là phương thức `extend()`. Cụ thể:\newline
- `a.extend(['d', 'e'])` duyệt qua danh sách `['d', 'e']` và lần lượt thêm `'d'`, `'e'` vào cuối danh sách `a`, kết quả `a` trở thành `['a', 'b', 'c', 'd', 'e']`.\newline
- Mặc dù trong Python, các cú pháp `a += ['d', 'e']`, `a += 'de'`, và `a[len(a):] = ['d', 'e']` về mặt kỹ thuật cũng tạo ra kết quả tương tự do tính linh hoạt của toán tử in-place add (`__iadd__`) và slice assignment, phương thức `extend()` là phương thức built-in chuẩn tắc được giảng dạy làm câu trả lời chính thức trong câu hỏi trắc nghiệm này.
\end{minipage}
\vspace{0.8cm}
\hrule
\vspace{0.5cm}

\begin{minipage}{\textwidth}
\subsection*{Question 163 \quad \footnotesize{[0163_q163_Question_163.png]}}
\noindent \textbf{How do you create a variable x with the numeric value 5?}

\begin{enumerate}[label=\Alph*.]
    \item x = 5
    \item x = int(5)
    \item \textbf{Both the answers A and B are correct} \hfill \textcolor{blue}{\textbf{[Correct]}}
    \item Both the answers A and B are incorrect
\end{enumerate}

\noindent \textbf{Correct Answer:} \textcolor{blue}{\textbf{C}} \\[4pt]
\noindent \textbf{Explanation:}\\
Trong Python:\newline
- Cách thông thường và phổ biến nhất là `x = 5`, khi đó trình thông dịch tự động gán giá trị số nguyên 5 có kiểu `int` cho biến `x`.\newline
- Cách `x = int(5)` sử dụng hàm tạo kiểu `int()` chuyển đổi hoặc trả về số nguyên 5, sau đó gán vào `x`. Biến `x` cũng nhận giá trị là 5 với kiểu `int`.\newline
Vì cả hai câu lệnh đều tạo ra biến `x` có giá trị số 5, nên phương án 'Both the answers A and B are correct' (đáp án C) là chính xác.
\end{minipage}
\vspace{0.8cm}
\hrule
\vspace{0.5cm}

\begin{minipage}{\textwidth}
\subsection*{Question 164 \quad \footnotesize{[0164_q164_Question_164.png]}}
\noindent \textbf{What is the correct syntax to output the type of a variable or object in Python?}

\begin{enumerate}[label=\Alph*.]
    \item print(typeof(x))
    \item \textbf{print(type(x))} \hfill \textcolor{blue}{\textbf{[Correct]}}
    \item print(typeof x)
    \item print(typeOf(x))
\end{enumerate}

\noindent \textbf{Correct Answer:} \textcolor{blue}{\textbf{B}} \\[4pt]
\noindent \textbf{Explanation:}\\
Trong Python:\newline
- Hàm tích hợp sẵn (built-in function) để kiểm tra kiểu dữ liệu của một biến hoặc đối tượng là `type()`. Khi kết hợp với hàm `print()`, cú pháp chuẩn xác là `print(type(x))`. Kết quả in ra dạng `<class 'int'>`, `<class 'str'>`, v.v.\newline
- Các cách viết như `typeof(x)`, `typeof x` hay `typeOf(x)` bắt nguồn từ các ngôn ngữ khác như JavaScript hoặc C#; trong Python chúng không tồn tại và sẽ gây lỗi `NameError: name 'typeof' is not defined`.
\end{minipage}
\vspace{0.8cm}
\hrule
\vspace{0.5cm}

\begin{minipage}{\textwidth}
\subsection*{Question 165 \quad \footnotesize{[0165_q165_Question_165.png]}}
\noindent \textbf{In Object Oriented Programming, what is another name for the "attributes" of an object?}

\begin{enumerate}[label=\Alph*.]
    \item forms
    \item methods
    \item portions
    \item messages
    \item \textbf{fields} \hfill \textcolor{blue}{\textbf{[Correct]}}
\end{enumerate}

\noindent \textbf{Correct Answer:} \textcolor{blue}{\textbf{E}} \\[4pt]
\noindent \textbf{Explanation:}\\
Trong lập trình hướng đối tượng (OOP):\newline
- Thuộc tính (`attributes`) đại diện cho dữ liệu, trạng thái hoặc đặc điểm của đối tượng. Trong các tài liệu và ngôn ngữ OOP khác nhau (như Java, C#, C++), các thuộc tính này thường được gọi là các trường dữ liệu (`fields`), biến thực thể (`instance variables`) hoặc thuộc tính dữ liệu (`data members`).\newline
- `methods` là các phương thức/hành vi (hàm gắn liền với đối tượng), không phải attributes.\newline
- `messages` là cơ chế giao tiếp giữa các đối tượng.\newline
Do đó, một tên gọi khác của `attributes` chính là `fields` (đáp án E).
\end{minipage}
\vspace{0.8cm}
\hrule
\vspace{0.5cm}

\begin{minipage}{\textwidth}
\subsection*{Question 166 \quad \footnotesize{[0166_q166_Question_166.png]}}
\noindent \textbf{Which of the following is rarely used in Object Oriented Programming?}

\begin{enumerate}[label=\Alph*.]
    \item Method
    \item Constructor
    \item Attribute
    \item \textbf{Destructor} \hfill \textcolor{blue}{\textbf{[Correct]}}
\end{enumerate}

\noindent \textbf{Correct Answer:} \textcolor{blue}{\textbf{D}} \\[4pt]
\noindent \textbf{Explanation:}\\
Trong lập trình hướng đối tượng hiện đại:\newline
- Các ngôn ngữ như Python, Java, C# đều tích hợp bộ thu gom rác tự động (Automatic Garbage Collection) để giải phóng vùng nhớ cho các đối tượng không còn được tham chiếu. Vì vậy, hàm hủy (`Destructor` – ví dụ phương thức `__del__()` trong Python) rất hiếm khi cần lập trình viên phải tự viết hoặc sử dụng thủ công.\newline
- Ngược lại, `Attribute` (thuộc tính), `Method` (phương thức) và `Constructor` (hàm khởi tạo `__init__`) là những thành phần cốt lõi được sử dụng trong hầu như mọi lớp đối tượng.\newline
Do đó, `Destructor` (đáp án D) là thành phần hiếm khi được sử dụng nhất.
\end{minipage}
\vspace{0.8cm}
\hrule
\vspace{0.5cm}

\begin{minipage}{\textwidth}
\subsection*{Question 167 \quad \footnotesize{[0167_q167_Question_167.png]}}
\noindent \textbf{What is the "wild card" character in a regular expression (i.e., the character that matches any character)?}

\begin{enumerate}[label=\Alph*.]
    \item ?
    \item *
    \item \textbf{.} \hfill \textcolor{blue}{\textbf{[Correct]}}
    \item $
\end{enumerate}

\noindent \textbf{Correct Answer:} \textcolor{blue}{\textbf{C}} \\[4pt]
\noindent \textbf{Explanation:}\\
Trong biểu thức chính quy (Regular Expressions - regex):\newline
- Ký tự dấu chấm `.` là ký tự đại diện (wildcard character) khớp với bất kỳ ký tự đơn nào (trừ ký tự xuống dòng `\n`, trừ khi cờ `re.DOTALL` được bật).\newline
- Ký tự `*` là một lượng từ (quantifier), mang ý nghĩa khớp 0 hoặc nhiều lần ký tự đứng trước nó, chứ không tự khớp với bất kỳ ký tự nào nếu đứng một mình.\newline
- Ký tự `?` biểu thị khớp 0 hoặc 1 lần ký tự đứng trước, hoặc được dùng làm non-greedy matching.\newline
- Ký tự `$` là ký tự neo khớp với vị trí kết thúc của chuỗi hoặc dòng.\newline
Vì vậy, ký tự 'wild card' khớp với mọi ký tự là dấu chấm `.` (đáp án C).
\end{minipage}
\vspace{0.8cm}
\hrule
\vspace{0.5cm}

\begin{minipage}{\textwidth}
\subsection*{Question 168 \quad \footnotesize{[0168_q168_Question_168.png]}}
\noindent \textbf{What is output for:\newline

\newline
\newline

\newline
b = [11,13,15,17,19,21]\newline

\newline
print(b[::2])}

\begin{enumerate}[label=\Alph*.]
    \item [19,21]
    \item [11,15]
    \item \textbf{[11,15,19]} \hfill \textcolor{blue}{\textbf{[Correct]}}
    \item [13,17,21]
\end{enumerate}

\noindent \textbf{Correct Answer:} \textcolor{blue}{\textbf{C}} \\[4pt]
\noindent \textbf{Explanation:}\\
Trong Python, cú pháp slicing danh sách có dạng `list[start:stop:step]`.\newline
Ở đây:\newline
- `start` và `stop` được bỏ qua (`None`), nghĩa là cắt từ đầu đến hết danh sách `b`.\newline
- `step = 2`: bước nhảy là 2, nghĩa là lấy các phần tử ở các chỉ mục chẵn `0, 2, 4, ...`:\newline
  + Chỉ mục 0: `b[0] = 11`\newline
  + Chỉ mục 2: `b[2] = 15`\newline
  + Chỉ mục 4: `b[4] = 19`\newline
  + Chỉ mục 6: vượt quá độ dài `len(b) = 6`, dừng duyệt.\newline
Kết quả trả về danh sách con là `[11, 15, 19]`. Do đó đáp án C là chính xác.
\end{minipage}
\vspace{0.8cm}
\hrule
\vspace{0.5cm}

\begin{minipage}{\textwidth}
\subsection*{Question 169 \quad \footnotesize{[0169_q169_Question_169.png]}}
\noindent \textbf{Which of these collections defines a TUPLE?}

\begin{enumerate}[label=\Alph*.]
    \item ["apple", "banana", "cherry"]
    \item "apple", "banana", "cherry"
    \item "name": "apple", "color": "green"
    \item \textbf{("apple", "banana", "cherry")} \hfill \textcolor{blue}{\textbf{[Correct]}}
\end{enumerate}

\noindent \textbf{Correct Answer:} \textcolor{blue}{\textbf{D}} \\[4pt]
\noindent \textbf{Explanation:}\\
Trong ngôn ngữ Python, các kiểu dữ liệu tập hợp (collection) được phân biệt qua cú pháp dấu ngoặc:\newline
- Tuple (bộ giá trị) được định nghĩa bằng cặp dấu ngoặc đơn `()`, phân tách nhau bởi dấu phẩy: `("apple", "banana", "cherry")`. Tuple là kiểu dữ liệu bất biến (immutable).\newline
- Phương án A sử dụng ngoặc vuông `[]` là kiểu `list`.\newline
- Phương án B sử dụng ngoặc nhọn `{}` chứa các phần tử đơn lẻ là kiểu `set`.\newline
- Phương án C sử dụng ngoặc nhọn `{}` chứa các cặp key: value là kiểu `dict`.\newline
Do đó đáp án đúng là D.
\end{minipage}
\vspace{0.8cm}
\hrule
\vspace{0.5cm}

\begin{minipage}{\textwidth}
\subsection*{Question 170 \quad \footnotesize{[0170_q170_Question_170.png]}}
\noindent \textbf{The readlines() method returns ...}

\begin{enumerate}[label=\Alph*.]
    \item str
    \item \textbf{a list of lines} \hfill \textcolor{blue}{\textbf{[Correct]}}
    \item a list of single characters
    \item a list of integers
\end{enumerate}

\noindent \textbf{Correct Answer:} \textcolor{blue}{\textbf{B}} \\[4pt]
\noindent \textbf{Explanation:}\\
Trong xử lý tập tin (file handling) của Python:\newline
- Phương thức `file.readlines()` đọc toàn bộ các dòng còn lại trong file và gom chúng vào một danh sách (`list of strings`), trong đó mỗi phần tử trong danh sách là một chuỗi biểu diễn một dòng văn bản kèm theo ký tự xuống dòng `\n`.\newline
- Phương thức `file.read()` trả về toàn bộ nội dung dưới dạng một chuỗi duy nhất (`str`).\newline
- Phương thức `file.readline()` đọc một dòng duy nhất và trả về kiểu `str`.\newline
Vì vậy, `readlines()` trả về một danh sách các dòng (`a list of lines` - đáp án B).
\end{minipage}
\vspace{0.8cm}
\hrule
\vspace{0.5cm}

\begin{minipage}{\textwidth}
\subsection*{Question 171 \quad \footnotesize{[0171_q171_Question_171.png]}}
\noindent \textbf{What does 'strip' do?}

\begin{enumerate}[label=\Alph*.]
    \item Removes the contents of the whole file
    \item \textbf{Removes extra data at the beginning and end of a string} \hfill \textcolor{blue}{\textbf{[Correct]}}
    \item Adds in metadata
    \item Splits the file
\end{enumerate}

\noindent \textbf{Correct Answer:} \textcolor{blue}{\textbf{B}} \\[4pt]
\noindent \textbf{Explanation:}\\
Trong Python, phương thức chuỗi `str.strip()` được sử dụng để loại bỏ các ký tự khoảng trắng thừa (whitespace characters như khoảng cách, `\t`, `\n`, `\r`) hoặc các ký tự chỉ định ở cả hai đầu của chuỗi (ở đầu và ở cuối chuỗi).\newline
- Nó không làm thay đổi các khoảng trắng ở giữa chuỗi.\newline
- Tùy chọn liên quan: `lstrip()` cắt ở đầu bên trái, `rstrip()` cắt ở đầu bên phải.\newline
Do đó, câu trả lời mô tả chính xác nhất chức năng của `strip` là 'Removes extra data at the beginning and end of a string' (đáp án B).
\end{minipage}
\vspace{0.8cm}
\hrule
\vspace{0.5cm}

\begin{minipage}{\textwidth}
\subsection*{Question 172 \quad \footnotesize{[0172_q172_Question_172.png]}}
\noindent \textbf{What is stored in a "file handle" that is returned from a successful open() call?}

\begin{enumerate}[label=\Alph*.]
    \item \textbf{The handle is a connection to the file's data} \hfill \textcolor{blue}{\textbf{[Correct]}}
    \item The handle contains the first 10 lines of a file
    \item The handle has a list of all of the files in a particular folder on the hard drive
    \item All the data from the file is read into memory and stored in the handle
\end{enumerate}

\noindent \textbf{Correct Answer:} \textcolor{blue}{\textbf{A}} \\[4pt]
\noindent \textbf{Explanation:}\\
Theo giáo trình Python for Everybody (PY4E Chapter 7 - Files):\newline
- Khi một lệnh `open(filename)` thực thi thành công, hệ điều hành và Python không lập tức đọc toàn bộ nội dung file nạp vào RAM (điều này sẽ gây cạn kiệt bộ nhớ nếu file có kích thước lớn).\newline
- Thay vào đó, hàm `open()` trả về một 'file handle' (đối tượng file). File handle này đóng vai trò là một đường dẫn/kết nối có kiểm soát (connection / pointer) đến dữ liệu của tập tin trên ổ đĩa, cho phép chương trình đọc, ghi hoặc duyệt qua từng dòng một cách hiệu quả theo yêu cầu.\newline
Do đó đáp án đúng là A: 'The handle is a connection to the file's data'.
\end{minipage}
\vspace{0.8cm}
\hrule
\vspace{0.5cm}

\begin{minipage}{\textwidth}
\subsection*{Question 173 \quad \footnotesize{[0173_q173_Question_173.png]}}
\noindent \textbf{Which of the following is not a method of opening files?}

\begin{enumerate}[label=\Alph*.]
    \item \textbf{Replace} \hfill \textcolor{blue}{\textbf{[Correct]}}
    \item Append
    \item Read
    \item Write
\end{enumerate}

\noindent \textbf{Correct Answer:} \textcolor{blue}{\textbf{A}} \\[4pt]
\noindent \textbf{Explanation:}\\
Trong thao tác mở tập tin với hàm `open(file, mode)` trong Python, các chế độ mở file (modes) chuẩn bao gồm:\newline
- `'r'` (Read): Chế độ đọc (mặc định), mở file để đọc dữ liệu.\newline
- `'w'` (Write): Chế độ ghi, tạo mới file hoặc ghi đè (xóa toàn bộ nội dung cũ) nếu file đã tồn tại.\newline
- `'a'` (Append): Chế độ ghi nối tiếp, ghi thêm dữ liệu vào cuối file hiện có.\newline
- `'x'` (Exclusive creation): Tạo file mới, báo lỗi nếu file đã tồn tại.\newline
Trong Python không có chế độ mở file nào mang tên `Replace`. Do đó `Replace` không phải là một phương thức/chế độ mở file (đáp án A).
\end{minipage}
\vspace{0.8cm}
\hrule
\vspace{0.5cm}

\begin{minipage}{\textwidth}
\subsection*{Question 174 \quad \footnotesize{[0174_q174_Question_174.png]}}
\noindent \textbf{What is the output of the following program?\newline

\newline
\newline

\newline
def outerFunction():\newline

\newline
    global a\newline

\newline
    a = 20\newline

\newline
    def innerFunction():\newline

\newline
        global a\newline

\newline
        a = 30\newline

\newline
        print('a =', a)\newline

\newline
a = 10\newline

\newline
outerFunction()\newline

\newline
print('a =', a)}

\begin{enumerate}[label=\Alph*.]
    \item a = 10 a = 30
    \item a = 10
    \item \textbf{a = 20} \hfill \textcolor{blue}{\textbf{[Correct]}}
    \item a = 30
\end{enumerate}

\noindent \textbf{Correct Answer:} \textcolor{blue}{\textbf{C}} \\[4pt]
\noindent \textbf{Explanation:}\\
Từng bước thực thi của chương trình:\newline
1. `def outerFunction(): ...`: Định nghĩa hàm `outerFunction`.\newline
2. `a = 10`: Khởi tạo biến toàn cục (global variable) `a` với giá trị là 10.\newline
3. Gọi hàm `outerFunction()`:\newline
   - Dòng `global a` thông báo rằng hàm này sẽ thao tác trực tiếp trên biến toàn cục `a`.\newline
   - `a = 20`: Giá trị của biến toàn cục `a` được cập nhật thành 20.\newline
   - Tiếp theo, hàm con `innerFunction()` được định nghĩa bên trong `outerFunction`, nhưng **không bao giờ được gọi** trong thân hàm `outerFunction`. Vì vậy toàn bộ đoạn mã bên trong `innerFunction` (bao gồm `a = 30` và `print('a =', a)`) không được thực thi.\newline
4. Sau khi hàm `outerFunction()` kết thúc, câu lệnh `print('a =', a)` ở phạm vi toàn cục được chạy, in ra giá trị hiện tại của `a` là 20.\newline
Định dạng xuất ra là: `a = 20` (đáp án C).
\end{minipage}
\vspace{0.8cm}
\hrule
\vspace{0.5cm}

\begin{minipage}{\textwidth}
\subsection*{Question 175 \quad \footnotesize{[0175_q175_Question_175.png]}}
\noindent \textbf{What is the output of print list if list = [ 'abcd', 786 , 2.23, 'john', 70.2 ]?}

\begin{enumerate}[label=\Alph*.]
    \item \textbf{[ 'abcd', 786 , 2.23, 'john', 70.2 ]} \hfill \textcolor{blue}{\textbf{[Correct]}}
    \item list
    \item Error
    \item None of the above
\end{enumerate}

\noindent \textbf{Correct Answer:} \textcolor{blue}{\textbf{A}} \\[4pt]
\noindent \textbf{Explanation:}\\
Biến `list` được khởi tạo là một danh sách chứa các phần tử: `[ 'abcd', 786 , 2.23, 'john', 70.2 ]`.\newline
Khi in danh sách này ra màn hình bằng lệnh in (`print list` trong Python 2 hoặc `print(list)` trong Python 3), Python sẽ gọi phương thức biểu diễn chuỗi `__str__()` / `__repr__()` của đối tượng danh sách, in ra toàn bộ cấu trúc của danh sách gồm các dấu ngoặc vuông và các phần tử ngăn cách bởi dấu phẩy: `[ 'abcd', 786 , 2.23, 'john', 70.2 ]`. Do đó đáp án đúng là A.
\end{minipage}
\vspace{0.8cm}
\hrule
\vspace{0.5cm}

\begin{minipage}{\textwidth}
\subsection*{Question 176 \quad \footnotesize{[0176_q176_Question_176.png]}}
\noindent \textbf{What does the Python input() function do?}

\begin{enumerate}[label=\Alph*.]
    \item \textbf{Pause the program and read data from the user} \hfill \textcolor{blue}{\textbf{[Correct]}}
    \item Connect to the network and retrieve a web page.
    \item Read the memory of the running program
    \item Take a screen shot from an area of the screen
\end{enumerate}

\noindent \textbf{Correct Answer:} \textcolor{blue}{\textbf{A}} \\[4pt]
\noindent \textbf{Explanation:}\\
Trong Python, hàm built-in `input([prompt])` có cơ chế hoạt động:\newline
- Tạm dừng (pause) việc thực thi chương trình để chờ người dùng nhập văn bản từ bàn phím (luồng dữ liệu chuẩn `sys.stdin`).\newline
- Khi người dùng hoàn thành việc nhập và nhấn phím Enter (gửi ký tự kết thúc dòng `\n`), hàm `input()` đọc chuỗi dữ liệu đó, loại bỏ ký tự `\n` ở cuối và trả về một chuỗi ký tự kiểu `str` cho chương trình tiếp tục chạy.\newline
Do đó đáp án đúng là A: 'Pause the program and read data from the user'.
\end{minipage}
\vspace{0.8cm}
\hrule
\vspace{0.5cm}

\begin{minipage}{\textwidth}
\subsection*{Question 177 \quad \footnotesize{[0177_q177_Question_177.png]}}
\noindent \textbf{Which of these statements is used to indicate that class A will inherit all the features of class B?}

\begin{enumerate}[label=\Alph*.]
    \item \textbf{class A(B) :} \hfill \textcolor{blue}{\textbf{[Correct]}}
    \item class A extends B :
    \item class A inherits B :
    \item A=B++;
\end{enumerate}

\noindent \textbf{Correct Answer:} \textcolor{blue}{\textbf{A}} \\[4pt]
\noindent \textbf{Explanation:}\\
Trong lập trình hướng đối tượng với Python, kế thừa (inheritance) được định nghĩa bằng cách đặt tên lớp cha (base class / superclass) vào trong dấu ngoặc đơn liền sau tên lớp con khi khai báo lớp:\newline
```python\newline
class A(B):\newline
    pass\newline
```\newline
Lớp `A` sẽ kế thừa toàn bộ các thuộc tính và phương thức từ lớp `B`.\newline
- Các từ khóa như `extends` hay `inherits` không tồn tại trong ngữ pháp của Python (chúng thuộc các ngôn ngữ như Java, PHP hay Visual Basic).\newline
- Cú pháp `A=B++;` là câu lệnh gán và toán tử tăng trong C/C++/Java, không hợp lệ trong Python.\newline
Do đó đáp án đúng là A.
\end{minipage}
\vspace{0.8cm}
\hrule
\vspace{0.5cm}

\begin{minipage}{\textwidth}
\subsection*{Question 178 \quad \footnotesize{[0178_q178_Question_178.png]}}
\noindent \textbf{What is a correct syntax to output "Hello World" in Python?}

\begin{enumerate}[label=\Alph*.]
    \item echo "Hello World"
    \item p("Hello World")
    \item \textbf{print("Hello World")} \hfill \textcolor{blue}{\textbf{[Correct]}}
    \item echo("Hello World");
\end{enumerate}

\noindent \textbf{Correct Answer:} \textcolor{blue}{\textbf{C}} \\[4pt]
\noindent \textbf{Explanation:}\\
Trong Python (từ Python 3 trở đi), hàm `print()` là hàm tích hợp sẵn tiêu chuẩn dùng để in văn bản hoặc đối tượng ra console / màn hình.\newline
- Cú pháp chuẩn xác là `print("Hello World")`.\newline
- `echo` là câu lệnh xuất trong PHP và Unix/Linux Bash Shell, không phải của Python.\newline
- Dấu chấm phẩy `;` ở cuối lệnh không bắt buộc và `echo()` không phải hàm chuẩn trong Python.\newline
- `p()` không phải hàm xuất chuẩn trong Python.\newline
Do đó đáp án C là chính xác.
\end{minipage}
\vspace{0.8cm}
\hrule
\vspace{0.5cm}

\begin{minipage}{\textwidth}
\subsection*{Question 179 \quad \footnotesize{[0179_q179_Question_179.png]}}
\noindent \textbf{How are Python dictionaries different from Python lists?}

\begin{enumerate}[label=\Alph*.]
    \item Python lists store multiple values and dictionaries store a single value
    \item Python dictionaries are a collection and lists are not a collection
    \item \textbf{Python lists are indexed using integers and dictionaries can use strings as indexes} \hfill \textcolor{blue}{\textbf{[Correct]}}
    \item Python lists can store strings and dictionaries can only store words
\end{enumerate}

\noindent \textbf{Correct Answer:} \textcolor{blue}{\textbf{C}} \\[4pt]
\noindent \textbf{Explanation:}\\
Sự khác biệt căn bản giữa `list` và `dict` trong Python:\newline
- List (danh sách) là cấu trúc tuần tự có thứ tự, các phần tử trong list được định vị và truy xuất dựa trên chỉ mục vị trí là các số nguyên liên tiếp (`integer indexes`) bắt đầu từ 0 (`0, 1, 2, ...`).\newline
- Dictionary (từ điển) là cấu trúc dạng bảng băm (hash map) lưu trữ các cặp `key: value`. Khóa (`key`) đóng vai trò như nhãn hay chỉ mục truy xuất và có thể là bất kỳ kiểu dữ liệu bất biến nào, phổ biến nhất là chuỗi (`strings as indexes/keys`), ví dụ: `d['name']`.\newline
- A sai vì cả hai đều là cấu trúc chứa nhiều phần tử/giá trị.\newline
- B sai vì cả hai đều là collection (kiểu tập hợp).\newline
- D sai vì dictionary có thể lưu trữ bất kỳ loại đối tượng nào, không chỉ là 'words'.\newline
Do đó đáp án đúng là C.
\end{minipage}
\vspace{0.8cm}
\hrule
\vspace{0.5cm}

\begin{minipage}{\textwidth}
\subsection*{Question 180 \quad \footnotesize{[0180_q180_Question_180.png]}}
\noindent \textbf{To read the entire remaining contents of the file as a string from a file object infile, we use?}

\begin{enumerate}[label=\Alph*.]
    \item infile.read(2)
    \item \textbf{infile.read()} \hfill \textcolor{blue}{\textbf{[Correct]}}
    \item infile.readline()
    \item infile.readlines()
\end{enumerate}

\noindent \textbf{Correct Answer:} \textcolor{blue}{\textbf{B}} \\[4pt]
\noindent \textbf{Explanation:}\\
Trong các phương thức đọc file của Python:\newline
- `infile.read()` (không truyền tham số kích thước): đọc toàn bộ nội dung còn lại của tập tin từ vị trí con trỏ hiện tại đến cuối file và trả về dưới dạng một chuỗi văn bản (`str`).\newline
- `infile.read(2)`: chỉ đọc tối đa 2 ký tự (hoặc 2 byte).\newline
- `infile.readline()`: chỉ đọc một dòng duy nhất tiếp theo từ vị trí con trỏ và trả về một chuỗi.\newline
- `infile.readlines()`: đọc toàn bộ các dòng còn lại nhưng trả về dưới dạng một danh sách các chuỗi (`list of strings`), chứ không phải là một chuỗi đơn lẻ (`a string`).\newline
Do đó câu lệnh đọc toàn bộ nội dung thành một chuỗi là `infile.read()` (đáp án B).
\end{minipage}
\vspace{0.8cm}
\hrule
\vspace{0.5cm}

\begin{minipage}{\textwidth}
\subsection*{Question 181 \quad \footnotesize{[0181_q181_Question_181.png]}}
\noindent \textbf{Suppose you are given a set(s1=1,2,3) then what is the output for the code - 2 * s1?}

\begin{enumerate}[label=\Alph*.]
    \item (1,1,2,2,3,3)
    \item [1,1,2,2,3,3]
    \item \textbf{Illegal} \hfill \textcolor{blue}{\textbf{[Correct]}}
    \item (1,2,3,1,2,3)
\end{enumerate}

\noindent \textbf{Correct Answer:} \textcolor{blue}{\textbf{C}} \\[4pt]
\noindent \textbf{Explanation:}\\
Trong Python, toán tử nhân `*` không được định nghĩa giữa kiểu số nguyên (`int`) và tập hợp (`set`). Tập hợp là một cấu trúc dữ liệu không có thứ tự và chỉ chứa các phần tử duy nhất (không cho phép phần tử trùng lặp), do đó nó không hỗ trợ phép toán nhân bản/lặp lại dãy phần tử (sequence repetition) như `list`, `tuple` hay `str`. Khi thực hiện `2 * s1` hoặc `s1 * 2`, Python sẽ báo lỗi `TypeError: unsupported operand type(s) for *: 'int' and 'set'`. Do đó, câu lệnh này là không hợp lệ ('Illegal').
\end{minipage}
\vspace{0.8cm}
\hrule
\vspace{0.5cm}

\begin{minipage}{\textwidth}
\subsection*{Question 182 \quad \footnotesize{[0182_q182_Question_182.png]}}
\noindent \textbf{What is the output of the following code?\newline

\newline
\newline

\newline
```python\newline

\newline
class Point:\newline

\newline
    def __init__(self, x = 0, y = 0):\newline

\newline
        self.x = x\newline

\newline
        self.y = y\newline

\newline
\newline

\newline
    def __sub__(self, other):\newline

\newline
        x = self.x - other.x\newline

\newline
        y = self.y - other.y\newline

\newline
        return Point(x,y)\newline

\newline
\newline

\newline
p1 = Point(3, 4)\newline

\newline
p2 = Point(1, 2)\newline

\newline
result = p1-p2\newline

\newline
print(result.x, result.y)\newline

\newline
```}

\begin{enumerate}[label=\Alph*.]
    \item \textbf{2 2} \hfill \textcolor{blue}{\textbf{[Correct]}}
    \item 4 6
    \item 0 0
    \item 1 1
\end{enumerate}

\noindent \textbf{Correct Answer:} \textcolor{blue}{\textbf{A}} \\[4pt]
\noindent \textbf{Explanation:}\\
Đoạn mã định nghĩa lớp `Point` với việc nạp chồng toán tử trừ `-` thông qua phương thức đặc biệt (magic/dunder method) `__sub__(self, other)`:\newline
1. Khởi tạo `p1 = Point(3, 4)`: `p1.x = 3`, `p1.y = 4`.\newline
2. Khởi tạo `p2 = Point(1, 2)`: `p2.x = 1`, `p2.y = 2`.\newline
3. Biểu thức `result = p1 - p2` tương đương với `result = p1.__sub__(p2)`.\newline
4. Bên trong phương thức `__sub__`:\newline
   - `x = self.x - other.x = 3 - 1 = 2`\newline
   - `y = self.y - other.y = 4 - 2 = 2`\newline
   - Trả về đối tượng `Point(2, 2)`.\newline
5. Lệnh `print(result.x, result.y)` in ra giá trị hai thuộc tính `result.x` và `result.y` phân tách bằng dấu cách: `2 2`.\newline
Do đó, đáp án chính xác là A.
\end{minipage}
\vspace{0.8cm}
\hrule
\vspace{0.5cm}

\begin{minipage}{\textwidth}
\subsection*{Question 183 \quad \footnotesize{[0183_q183_Question_183.png]}}
\noindent \textbf{Which collection is ordered, changeable, and allows duplicate members?}

\begin{enumerate}[label=\Alph*.]
    \item SET
    \item DICTIONARY
    \item TUPLE
    \item \textbf{LIST} \hfill \textcolor{blue}{\textbf{[Correct]}}
\end{enumerate}

\noindent \textbf{Correct Answer:} \textcolor{blue}{\textbf{D}} \\[4pt]
\noindent \textbf{Explanation:}\\
Trong tài liệu chuẩn của Python và phân loại tập hợp (collections):\newline
- **LIST**: Là tập hợp có thứ tự (ordered), có thể thay đổi/chỉnh sửa phần tử (changeable/mutable), và cho phép chứa các phần tử trùng lặp (allows duplicate members).\newline
- **SET**: Là tập hợp không có thứ tự chỉ mục (unordered), có thể thay đổi nhưng KHÔNG cho phép phần tử trùng lặp (no duplicate members).\newline
- **TUPLE**: Là tập hợp có thứ tự (ordered), cho phép phần tử trùng lặp, nhưng KHÔNG thể thay đổi sau khi tạo (unchangeable/immutable).\newline
- **DICTIONARY**: Là tập hợp các cặp key-value có thứ tự (từ Python 3.7+), có thể thay đổi (changeable), nhưng các khóa (keys) KHÔNG được phép trùng lặp.\newline
Do đó, kiểu tập hợp đáp ứng đủ 3 tiêu chí 'ordered, changeable, and allows duplicate members' là LIST (đáp án D).
\end{minipage}
\vspace{0.8cm}
\hrule
\vspace{0.5cm}

\begin{minipage}{\textwidth}
\subsection*{Question 184 \quad \footnotesize{[0184_q184_Question_184.png]}}
\noindent \textbf{Which of the following slicing operations will produce the list [12, 3]? t = [9, 41, 12, 3, 74, 15]}

\begin{enumerate}[label=\Alph*.]
    \item t[1:3]
    \item t[2:2]
    \item t[:]
    \item \textbf{t[2:4]} \hfill \textcolor{blue}{\textbf{[Correct]}}
\end{enumerate}

\noindent \textbf{Correct Answer:} \textcolor{blue}{\textbf{D}} \\[4pt]
\noindent \textbf{Explanation:}\\
Xét danh sách `t = [9, 41, 12, 3, 74, 15]` có các phần tử tại các chỉ mục (index):\newline
- `t[0] = 9`\newline
- `t[1] = 41`\newline
- `t[2] = 12`\newline
- `t[3] = 3`\newline
- `t[4] = 74`\newline
- `t[5] = 15`\newline
\newline
Cú pháp cắt lát danh sách (slicing) `t[start:stop]` lấy các phần tử từ chỉ số `start` đến `stop - 1`:\newline
- `t[1:3]` lấy phần tử tại chỉ số 1 và 2 -> `[41, 12]`.\newline
- `t[2:2]` với `start == stop` trả về danh sách rỗng `[]`.\newline
- `t[:]` sao chép toàn bộ danh sách `t` -> `[9, 41, 12, 3, 74, 15]`.\newline
- `t[2:4]` lấy phần tử tại chỉ số 2 và 3 -> `[12, 3]`.\newline
Do đó, `t[2:4]` tạo ra danh sách `[12, 3]`. Đáp án đúng là D.
\end{minipage}
\vspace{0.8cm}
\hrule
\vspace{0.5cm}

\begin{minipage}{\textwidth}
\subsection*{Question 185 \quad \footnotesize{[0185_q185_Question_185.png]}}
\noindent \textbf{What does the following Python program print out?\newline

\newline
\newline

\newline
```python\newline

\newline
x = '40'\newline

\newline
y = int(x) + 2\newline

\newline
print(y)\newline

\newline
```}

\begin{enumerate}[label=\Alph*.]
    \item int402
    \item 402
    \item \textbf{42} \hfill \textcolor{blue}{\textbf{[Correct]}}
    \item x2
\end{enumerate}

\noindent \textbf{Correct Answer:} \textcolor{blue}{\textbf{C}} \\[4pt]
\noindent \textbf{Explanation:}\\
Phân tích từng dòng lệnh của chương trình:\newline
1. `x = '40'`: Khởi tạo biến `x` lưu trữ một chuỗi ký tự (`str`).\newline
2. `int(x)`: Hàm tích hợp sẵn `int()` chuyển đổi chuỗi chữ số `'40'` thành số nguyên có giá trị `40`.\newline
3. `y = int(x) + 2`: Thực hiện phép cộng số học giữa hai số nguyên `40 + 2 = 42`, gán giá trị này cho biến `y`.\newline
4. `print(y)`: In giá trị của biến `y` ra màn hình, kết quả là `42`.\newline
Do đó, đáp án đúng là C.
\end{minipage}
\vspace{0.8cm}
\hrule
\vspace{0.5cm}

\begin{minipage}{\textwidth}
\subsection*{Question 186 \quad \footnotesize{[0186_q186_Question_186.png]}}
\noindent \textbf{What is a correct syntax to return the first character in a string?}

\begin{enumerate}[label=\Alph*.]
    \item x = "Hello".sub(0, 1)
    \item x = "Hello"[1]
    \item \textbf{x = "Hello"[0]} \hfill \textcolor{blue}{\textbf{[Correct]}}
    \item x = sub("Hello", 0, 1)
\end{enumerate}

\noindent \textbf{Correct Answer:} \textcolor{blue}{\textbf{C}} \\[4pt]
\noindent \textbf{Explanation:}\\
Trong Python:\newline
- Chuỗi ký tự (`str`) là kiểu dữ liệu dạng chuỗi tuần tự (sequence) và được đánh chỉ mục bắt đầu từ 0 (zero-based indexing).\newline
- Để truy xuất ký tự đầu tiên của chuỗi, ta dùng cú pháp lập chỉ mục `string[0]`. Với chuỗi `"Hello"`, `"Hello"[0]` trả về ký tự `'H'` (ký tự đầu tiên). Do đó câu lệnh `x = "Hello"[0]` là hoàn toàn chính xác.\newline
- `x = "Hello"[1]` lấy ký tự tại vị trí index 1, tức là ký tự thứ hai `'e'`.\newline
- Python không hỗ trợ phương thức `.sub()` hay hàm `sub()` dựng sẵn cho chuỗi (lệnh gọi này sẽ gây lỗi `AttributeError` hoặc `NameError`).\newline
Vì vậy đáp án đúng là C.
\end{minipage}
\vspace{0.8cm}
\hrule
\vspace{0.5cm}

\begin{minipage}{\textwidth}
\subsection*{Question 187 \quad \footnotesize{[0187_q187_Question_187.png]}}
\noindent \textbf{Consider this assignment statement: a, b, c = (1, 2, 3, 4, 5, 6, 7, 8, 9)[1::3] Following execution of this statement, what is the value of b:}

\begin{enumerate}[label=\Alph*.]
    \item 2
    \item 4
    \item \textbf{5} \hfill \textcolor{blue}{\textbf{[Correct]}}
    \item 6
\end{enumerate}

\noindent \textbf{Correct Answer:} \textcolor{blue}{\textbf{C}} \\[4pt]
\noindent \textbf{Explanation:}\\
Xét phép cắt lát (extended slicing) trên tuple `(1, 2, 3, 4, 5, 6, 7, 8, 9)[1::3]`:\newline
- Tuple có 9 phần tử tương ứng với các chỉ mục từ 0 đến 8:\newline
  - Index 0: 1\newline
  - Index 1: 2\newline
  - Index 2: 3\newline
  - Index 3: 4\newline
  - Index 4: 5\newline
  - Index 5: 6\newline
  - Index 6: 7\newline
  - Index 7: 8\newline
  - Index 8: 9\newline
- Cú pháp `[start:stop:step]`:\newline
  - `start = 1`: Bắt đầu lấy từ index 1 (giá trị là `2`).\newline
  - `step = 3`: Bước nhảy bằng 3. Vị trí tiếp theo là `1 + 3 = 4` (giá trị là `5`), kế tiếp là `4 + 3 = 7` (giá trị là `8`), tiếp theo `7 + 3 = 10` vượt quá độ dài.\newline
- Kết quả của phép cắt lát là tuple: `(2, 5, 8)`.\newline
- Phép gán dỡ gói đa biến (unpacking assignment): `a, b, c = (2, 5, 8)` sẽ gán:\newline
  - `a = 2`\newline
  - `b = 5`\newline
  - `c = 8`\newline
Do đó, giá trị của biến `b` sau khi thực hiện lệnh là 5. Đáp án đúng là C.
\end{minipage}
\vspace{0.8cm}
\hrule
\vspace{0.5cm}

\begin{minipage}{\textwidth}
\subsection*{Question 188 \quad \footnotesize{[0188_q188_Question_188.png]}}
\noindent \textbf{What's the output of the following code snippet?\newline

\newline
\newline

\newline
```python\newline

\newline
>>> class Dog:\newline

\newline
...     def walk(self):\newline

\newline
...         return "*walking*"\newline

\newline
...\newline

\newline
...     def speak(self):\newline

\newline
...         return "Woof!"\newline

\newline
...\newline

\newline
>>> class JackRussellTerrier(Dog):\newline

\newline
...     def speak(self):\newline

\newline
...         return "Arff!"\newline

\newline
...\newline

\newline
>>> bobo = JackRussellTerrier()\newline

\newline
>>> bobo.speak()\newline

\newline
```}

\begin{enumerate}[label=\Alph*.]
    \item \textbf{Arff!} \hfill \textcolor{blue}{\textbf{[Correct]}}
    \item walking
    \item Woof!
    \item Incorrect
    \item CanineError: Dog malfunction
\end{enumerate}

\noindent \textbf{Correct Answer:} \textcolor{blue}{\textbf{A}} \\[4pt]
\noindent \textbf{Explanation:}\\
Đoạn mã minh họa tính kế thừa (Inheritance) và ghi đè phương thức (Method Overriding) trong lập trình hướng đối tượng (OOP) của Python:\newline
1. Lớp cha `Dog` có phương thức `walk` trả về `"*walking*"` và `speak` trả về `"Woof!"`.\newline
2. Lớp con `JackRussellTerrier` kế thừa từ `Dog` (`class JackRussellTerrier(Dog):`) và ghi đè phương thức `speak(self)` để trả về `"Arff!"`.\newline
3. Biến `bobo` được tạo là một instance của lớp `JackRussellTerrier`.\newline
4. Khi gọi `bobo.speak()`, theo thứ tự phân giải phương thức (MRO), Python tìm và thực thi phương thức `speak` được định nghĩa trong lớp `JackRussellTerrier` thay vì phương thức của lớp cha `Dog`.\newline
5. Phương thức trả về chuỗi `"Arff!"` và được hiển thị trên màn hình console tương tác (REPL) là `Arff!`.\newline
Do đó, đáp án đúng là A.
\end{minipage}
\vspace{0.8cm}
\hrule
\vspace{0.5cm}

\begin{minipage}{\textwidth}
\subsection*{Question 189 \quad \footnotesize{[0189_q189_Question_189.png]}}
\noindent \textbf{How would you use string slicing [:] to print out 'uct' from the following string?\newline

\newline
\newline

\newline
```python\newline

\newline
x = 'From marquard@uct.ac.za'\newline

\newline
```}

\begin{enumerate}[label=\Alph*.]
    \item print(x[14/17])
    \item \textbf{print(x[14:17])} \hfill \textcolor{blue}{\textbf{[Correct]}}
    \item print(x[14:3])
    \item print(x[14+3])
\end{enumerate}

\noindent \textbf{Correct Answer:} \textcolor{blue}{\textbf{B}} \\[4pt]
\noindent \textbf{Explanation:}\\
Xét chuỗi `x = 'From marquard@uct.ac.za'`:\newline
Đếm chỉ số từng ký tự (bắt đầu từ 0):\newline
- Chỉ số 0 đến 3: `'From'`\newline
- Chỉ số 4: dấu khoảng trắng `' '`\newline
- Chỉ số 5 đến 12: `'marquard'`\newline
- Chỉ số 13: ký tự `'@'`\newline
- Chỉ số 14: ký tự `'u'`\newline
- Chỉ số 15: ký tự `'c'`\newline
- Chỉ số 16: ký tự `'t'`\newline
- Chỉ số 17: ký tự `'.'`\newline
\newline
Cú pháp cắt lát chuỗi là `x[start:stop]`, lấy từ chỉ mục `start` đến `stop - 1`:\newline
- `x[14:17]` lấy các ký tự tại chỉ mục 14, 15, 16 -> kết quả là chuỗi `'uct'`.\newline
- Phương án A `x[14/17]` sử dụng toán tử chia trả về float, gây lỗi `TypeError: string indices must be integers`.\newline
- Phương án C `x[14:3]` có chỉ số bắt đầu lớn hơn chỉ số kết thúc với bước nhảy mặc định +1 nên trả về chuỗi rỗng `''`.\newline
- Phương án D `x[14+3]` tương đương `x[17]`, trả về một ký tự duy nhất `'.'`. \newline
Do đó, lệnh chính xác là `print(x[14:17])`. Đáp án đúng là B.
\end{minipage}
\vspace{0.8cm}
\hrule
\vspace{0.5cm}

\begin{minipage}{\textwidth}
\subsection*{Question 190 \quad \footnotesize{[0190_q190_Question_190.png]}}
\noindent \textbf{What is the ASCII character that is associated with the decimal value 42?}

\begin{enumerate}[label=\Alph*.]
    \item !
    \item \textbf{*} \hfill \textcolor{blue}{\textbf{[Correct]}}
    \item ^
    \item /
\end{enumerate}

\noindent \textbf{Correct Answer:} \textcolor{blue}{\textbf{B}} \\[4pt]
\noindent \textbf{Explanation:}\\
Trong bảng mã chuẩn ASCII:\newline
- Mã thập phân 33 tương ứng với ký tự `!` (`chr(33) == '!'`)\newline
- Mã thập phân 42 tương ứng với ký tự dấu sao `*` (`chr(42) == '*'`)\newline
- Mã thập phân 47 tương ứng với ký tự gạch chéo `/` (`chr(47) == '/'`)\newline
- Mã thập phân 94 tương ứng với ký tự dấu mũ `^` (`chr(94) == '^'`)\newline
Trong Python, hàm dựng sẵn `chr(n)` trả về ký tự tương ứng với mã ASCII/Unicode `n`, và `chr(42)` trả về `'*'`. Do đó, đáp án đúng là B.
\end{minipage}
\vspace{0.8cm}
\hrule
\vspace{0.5cm}

\begin{minipage}{\textwidth}
\subsection*{Question 191 \quad \footnotesize{[0191_q191_Question_191.png]}}
\noindent \textbf{What will be the output of the following Python code?\newline

\newline
\newline

\newline
```python\newline

\newline
f = None\newline

\newline
for i in range (5):\newline

\newline
    with open("data.txt", "w") as f:\newline

\newline
        if i > 2:\newline

\newline
            break\newline

\newline
print(f.closed)\newline

\newline
```}

\begin{enumerate}[label=\Alph*.]
    \item \textbf{True} \hfill \textcolor{blue}{\textbf{[Correct]}}
    \item False
    \item None
    \item Error
\end{enumerate}

\noindent \textbf{Correct Answer:} \textcolor{blue}{\textbf{A}} \\[4pt]
\noindent \textbf{Explanation:}\\
Phân tích luồng thực thi của đoạn mã và trình quản lý ngữ cảnh (Context Manager):\newline
1. Ban đầu `f = None`.\newline
2. Vòng lặp `for i in range(5)` lần lượt duyệt qua các giá trị `i = 0, 1, 2, 3, 4`.\newline
3. Khi `i = 0, 1, 2`, biểu thức `i > 2` là `False`, khối `with` mở file và tự động đóng file an toàn khi kết thúc khối ở mỗi lần lặp.\newline
4. Khi `i = 3`, điều kiện `i > 2` thỏa mãn (`True`), câu lệnh `break` được thực thi để thoát khỏi vòng lặp.\newline
5. Điều quan trọng trong cơ chế `with` của Python: Dù khối lệnh bên trong bị ngắt sớm bởi `break`, `return` hay xảy ra ngoại lệ (exception), phương thức `__exit__()` của context manager luôn được gọi ngay lập tức trước khi luồng điều khiển nhảy ra ngoài. Phương thức này tự động gọi `f.close()` để đóng file.\newline
6. Sau khi thoát khỏi vòng lặp, thuộc tính `f.closed` của đối tượng file trả về `True` (xác nhận tệp đã được đóng).\newline
7. `print(f.closed)` in ra `True`.\newline
Do đó, đáp án đúng là A.
\end{minipage}
\vspace{0.8cm}
\hrule
\vspace{0.5cm}

\begin{minipage}{\textwidth}
\subsection*{Question 192 \quad \footnotesize{[0192_q192_Question_192.png]}}
\noindent \textbf{You want to create a bytes object consisting of five null (0x00) bytes. Choose all the apply.}

\begin{enumerate}[label=\Alph*.]
    \item \textbf{bytes(5)} \hfill \textcolor{blue}{\textbf{[Correct]}}
    \item bytes(0, 0, 0, 0, 0)
    \item \textbf{bytes('x 00x 00x 00x 00x 00', 'utf-8')} \hfill \textcolor{blue}{\textbf{[Correct]}}
    \item bytes[0:5]
\end{enumerate}

\noindent \textbf{Correct Answer:} \textcolor{blue}{\textbf{A, C}} \\[4pt]
\noindent \textbf{Explanation:}\\
Trong Python, để tạo một đối tượng kiểu `bytes` chứa 5 null byte `b'\x00\x00\x00\x00\x00'`:\newline
- **Phương án A: `bytes(5)`** — Hàm khởi tạo `bytes(int)` với một số nguyên `n` sẽ tạo ra một đối tượng byte có độ dài `n` với các byte đều được khởi tạo bằng giá trị null (`0x00`). Do đó `bytes(5)` trả về đúng `b'\x00\x00\x00\x00\x00'`. (Đúng)\newline
- **Phương án B: `bytes(0, 0, 0, 0, 0)`** — Hàm dựng `bytes()` chỉ nhận tối đa 3 tham số `(source, encoding, errors)`. Truyền vào 5 đối số riêng biệt sẽ gây lỗi `TypeError: bytes() takes at most 3 arguments (5 given)`. Nếu muốn khởi tạo từ các số 0 thì phải truyền vào một iterable (ví dụ: `bytes([0, 0, 0, 0, 0])`). (Sai)\newline
- **Phương án C: `bytes('\x00\x00\x00\x00\x00', 'utf-8')`** — Khi truyền vào một chuỗi ký tự (`str`), cần chỉ định bảng mã `encoding`. Ký tự null `\x00` trong UTF-8 được mã hóa thành byte đơn có giá trị `0x00`, do đó chuỗi 5 ký tự `\x00` mã hóa UTF-8 tạo ra đối tượng bytes gồm 5 byte null `b'\x00\x00\x00\x00\x00'`. (Đúng)\newline
- **Phương án D: `bytes[0:5]`** — `bytes` là một kiểu dữ liệu (type/class), không phải là đối tượng mảng/danh sách cụ thể để cắt lát bằng dấu ngoặc vuông `[]`, câu lệnh này gây ra lỗi `TypeError: 'type' object is not subscriptable`. (Sai)\newline
Vì vậy, hai lựa chọn đúng là A và C.
\end{minipage}
\vspace{0.8cm}
\hrule
\vspace{0.5cm}

\begin{minipage}{\textwidth}
\subsection*{Question 193 \quad \footnotesize{[0193_q193_Question_193.png]}}
\noindent \textbf{Which of the following define the set 'a', 'b', 'c':}

\begin{enumerate}[label=\Alph*.]
    \item s = set(['a', 'b', 'c'])
    \item \textbf{s = 'a', 'b', 'c'} \hfill \textcolor{blue}{\textbf{[Correct]}}
    \item s = ('a', 'b', 'c')
    \item s = set('a', 'b', 'c')
\end{enumerate}

\noindent \textbf{Correct Answer:} \textcolor{blue}{\textbf{B}} \\[4pt]
\noindent \textbf{Explanation:}\\
Để định nghĩa trực tiếp một tập hợp (set literal) chứa ba phần tử `'a'`, `'b'`, `'c'`:\newline
- Cú pháp chuẩn và tự nhiên nhất của Python (set display / literal notation) là sử dụng cặp ngoặc nhọn `{}` chứa các phần tử phân tách bằng dấu phẩy: `s = {'a', 'b', 'c'}` (phương án B).\newline
- Phương án A (`s = set(['a', 'b', 'c'])`): Đây là việc gọi hàm khởi tạo `set()` từ một danh sách (`list`), tạo ra tập hợp tương đương về mặt giá trị nhưng không phải là cú pháp khai báo hằng tập hợp (literal definition).\newline
- Phương án C (`s = {('a', 'b', 'c')}`): Tạo ra một tập hợp chỉ chứa duy nhất 1 phần tử là một tuple gồm 3 phần tử `('a', 'b', 'c')`, tức `len(s) == 1` chứ không phải tập hợp gồm 3 phần tử độc lập.\newline
- Phương án D (`s = set('a', 'b', 'c')`): Gây lỗi cú pháp thực thi `TypeError: set expected at most 1 argument, got 3` do hàm `set()` chỉ nhận tối đa 1 tham số là một iterable.\newline
Vì vậy, đáp án được chọn là B.
\end{minipage}
\vspace{0.8cm}
\hrule
\vspace{0.5cm}

\end{document}
